\documentclass[11pt,a4paper]{article}
\usepackage{fontspec}
\usepackage{amsmath,amssymb,amsthm,mathtools}
\usepackage[margin=24mm]{geometry}
\usepackage{xspace}
\usepackage{xcolor,booktabs,longtable,array,enumitem,fancyhdr,titlesec}
\usepackage{tocloft}\setlength{\cftsecnumwidth}{3.2em}\setlength{\cftpartnumwidth}{3.2em}
\usepackage[colorlinks=true,linkcolor=blue!50!black,citecolor=blue!50!black,urlcolor=blue!50!black]{hyperref}
\setlength{\parindent}{1em}
\setlength{\parskip}{2pt}
\linespread{1.08}
\emergencystretch=4em
\setlist{itemsep=1pt,topsep=3pt}
\setlength{\tabcolsep}{4pt}

%%% --- status tags ---
\newcommand{\tP}{\textcolor{green!45!black}{\textbf{[Proved]}}\xspace}
\newcommand{\tC}{\textcolor{blue!55!black}{\textbf{[Lit.]}}\xspace}
\newcommand{\tV}{\textcolor{violet!80!black}{\textbf{[Verified numerically]}}\xspace}
\newcommand{\tK}{\textcolor{red!70!black}{\textbf{[Conjecture]}}\xspace}
\newcommand{\tQ}{\textcolor{orange!85!black}{\textbf{[To be verified]}}\xspace}

%%% --- theorem-like environments ---
\theoremstyle{definition}
\newtheorem{defn}{Definition}
\newtheorem{thm}[defn]{Theorem}
\newtheorem{prop}[defn]{Proposition}
\newtheorem{lem}[defn]{Lemma}
\newtheorem{cor}[defn]{Corollary}
\newtheorem{ex}[defn]{Example}
\newtheorem{rem}[defn]{Remark}
\newtheorem{obs}[defn]{Observation}
\newtheorem{conj}[defn]{Conjecture}
\newcommand{\src}[1]{\par\noindent{\small\textcolor{blue!55!black}{\textbf{[Source]}~#1}}\par}

\newenvironment{hyptab}{\par\medskip\noindent{\small\textbf{Hypotheses and their roles}}\par\nopagebreak\small\noindent\begin{tabular}{>{\raggedright\arraybackslash}p{3.3cm}>{\raggedright\arraybackslash}p{5.5cm}>{\raggedright\arraybackslash}p{5.9cm}}\toprule Hypothesis & Why it is needed & Where it enters the proof\\\midrule}{\\\bottomrule\end{tabular}\par\medskip}

%%% --- notation ---
\newcommand{\R}{\mathbb R}\newcommand{\C}{\mathbb C}\newcommand{\Z}{\mathbb Z}
\newcommand{\bT}{\mathbb T}
\newcommand{\Lb}{\overline L}
\newcommand{\Jc}{\mathcal J}
\newcommand{\dbar}{\bar\partial}
\newcommand{\ch}{\mathrm{ch}}\newcommand{\Td}{\mathrm{Td}}
\newcommand{\rk}{\mathrm{rk}}
\newcommand{\HH}{\mathrm{HH}}
\newcommand{\Res}{\mathrm{Res}}
\newcommand{\dlog}{\mathrm{dlog}\,}
\newcommand{\sh}{\sharp}

\pagestyle{fancy}\fancyhf{}
\fancyhead[L]{\small The three treasures of generalized geometry}
\fancyhead[R]{\small Research note 2026-11-22 (integrated)}
\fancyfoot[C]{\small\thepage}
\titleformat{\section}{\large\bfseries\sffamily}{\thesection}{0.8em}{}
\titleformat{\subsection}{\normalsize\bfseries\sffamily}{\thesubsection}{0.8em}{}


\newcommand{\Q}{\mathbb Q}\newcommand{\N}{\mathbb N}
\newcommand{\Hc}{H_c}
\newcommand{\we}{\wedge}
\setcounter{tocdepth}{1}
\renewcommand{\proofname}{Proof}

\begin{document}
\sloppy
\begin{center}
{\LARGE\bfseries\sffamily Approaching the ``three treasures'' of generalized complex geometry}\\[6pt]
{\Large\bfseries\sffamily From the example of holomorphic Poisson structures (integrated paper)}\\[12pt]
{\large Research and writing: Claude\ /\ Problem posed by Mori}\\[3pt]
22 November 2026 (integrated version v1; English translation of the Japanese original)\\[3pt]
{\small (An unreviewed research note. The comprehensive version, Parts I--V, is rewritten by theme.\\ Unverified points are marked in the text by \tQ.)}
\end{center}
\begin{abstract}\noindent
We study the generalized complex versions of the three treasures of complex geometry (Riemann--Roch, Serre vanishing, Kodaira embedding) starting from the simplest non-trivial example: the type-jumping generalized complex structure obtained as the $\beta$-transform of a holomorphic Poisson structure $\sigma$, together with a generalized holomorphic line bundle $L$ on it.
(1) A generalized holomorphic structure is a Poisson module; if $\sigma$ is log-symplectic it is a flat logarithmic connection and the residues $r_i$ are constants.
(2) Serre-type vanishing holds in transversal degree, and fails in all degrees. Duality $\HH^k(L)^*\cong\HH^{2n-k}(L^\vee\otimes K^2)$ holds in any dimension.
(3) The total Euler characteristic is $(-1)^n\rk(E)\chi(X)$, independent of $c_1$. For surfaces and toric varieties we write the Poincar\'e polynomial in terms of the residues, and prove a local formula for quasi-homogeneous singularities.
(4) In general dimension the problem reduces, by a reduction theorem, to a condition on the quotient complex, and for non-integral residues we show $H^k(U;L_U)=0$ for $k\ne n$.
(5) In regular type, the dimensions in each degree are polynomial in $m$; in log-symplectic type they are periodic.
(6) For Kodaira type, embedding by generalized holomorphic sections is impossible; we prove a leaf-space version and the toric Poisson--Kodaira theorem.
The generalization of the index theorem (weighted $b$-calculus) remains a conjecture, and general Poisson--Kodaira is open.
\end{abstract}
\tableofcontents
\clearpage
%==============================================================
\section{Introduction and main results}\label{sec:intro}
%==============================================================
\subsection{Motivation and setting}
The ``three treasures'' of complex geometry are the Hirzebruch--Riemann--Roch theorem (HRR), Serre's vanishing theorem, and Kodaira's embedding theorem. The ultimate goal of this study is to ask whether they extend to generalized complex geometry (Hitchin, Gualtieri). As far as we could find, there are Kodaira--Nakano vanishing and Riemann--Roch--Hirzebruch on Lie groupoids by Hu \cite{Hu1,Hu2}, and vanishing theorems for strong generalized holomorphic bundles (regular type) by Pal--Poddar \cite{PP}, but we found no Kodaira embedding theorem or Riemann--Roch theorem for a general generalized complex structure \tC.

We therefore start from the simplest non-trivial example: the generalized complex structure $\Jc_\sigma$ obtained as the $\beta$-transform of a holomorphic Poisson structure $\sigma$ on a complex manifold $X$. This example has two advantages. (i) The type jumps along the divisor $D$ where $\sigma$ degenerates (complex type and non-degenerate symplectic type coexist). (ii) Generalized holomorphic line bundles can be described concretely as Poisson modules (holomorphic bundles with a flat $A$-connection), and generalized Dolbeault cohomology can be computed as holomorphic Poisson hypercohomology. This paper collects the results obtained for this example into a single paper organized by theme.

\paragraph{Policy of this paper.}
(a) The paper is self-contained: every claim whose proof is given is proved in the text, and a claim whose proof is omitted for reasons of space is marked as such, with the script or computation named, and tagged \tV. (b) Every claim carries a status tag (\S\ref{subsec:tags}). A status is never upgraded beyond what is proved in the text (no conjecture is turned into a proved statement). (c) Definitions, hypotheses, the table ``Hypotheses and their roles'' (why each hypothesis is needed and where it enters the proof), and the symbol table are kept. (d) Numbering of theorems, lemmas, etc.\ is a single global numbering.

\subsection{Overview of the main results}
Throughout, $X$ is a compact complex manifold ($\dim_\C X=n$), $\sigma$ is a holomorphic Poisson structure, $E$ is a generalized holomorphic bundle (Poisson module), $\HH^k(E)$ is the generalized Dolbeault cohomology, and $P_E(t)=\sum_k\dim\HH^k(E)\,t^k$ is the Poincar\'e polynomial. If $\sigma$ is log-symplectic along a simple normal crossing divisor $D$, then $E$ carries a flat logarithmic connection along $D$ whose residues $r_i$ are constants.

\begin{enumerate}
\item[(1)] \textbf{Preliminaries (\S\ref{sec:prelim}).} A generalized holomorphic bundle is a Poisson module, and $\HH^*(E)=H^*(X;(\wedge^\bullet T\otimes E,\ [\sigma,\cdot]+V\wedge))$. If $\sigma$ is log-symplectic, a Poisson module is a flat logarithmic connection with constant residues and $c_1(E)=-\sum r_i[D_i]$ (Theorem \ref{thm:main1}). Generalized holomorphic sections are flat (Casimir-type) sections.
\item[(2)] \textbf{Serre-type vanishing and duality (\S\ref{sec:serre}).} If $X$ is projective and $L_0$ is an ample Poisson module, then for $m\gg0$ all cohomology of transversal degree $a>0$ vanishes (Theorem \ref{thm:serre}). Vanishing in all degrees fails even for $\sigma=0$. Duality $\HH^k(L)^*\cong\HH^{2n-k}(L^\vee\otimes K_X^2)$ holds in any dimension for any $\sigma$ (Theorem \ref{thm:dual}).
\item[(3)] \textbf{Euler characteristic (\S\ref{sec:euler}).} The total Euler characteristic is $(-1)^n\rk(E)\chi(X)$, independent of $c_1(E)$ (Theorem \ref{thm:E}). The Euler characteristic of the log de Rham complex is $\chi(U)$, $U=X\setminus D$ (Theorem \ref{thm:K1}).
\item[(4)] \textbf{Log-symplectic surfaces (\S\ref{sec:surface}).} For $r_i\neq1$, $\HH^*(E)=H^*(X;\mathcal F)\oplus\C^\mu[-2]$ ($\mathcal F$ is an $Rj_*$/$j_!$-type extension of $L_U$, $\mu$ is the number of nodes), a closed formula for the Poincar\'e polynomial (Theorems \ref{thm:main2}, \ref{thm:B}), and periodic stabilization in $m$ (Theorem \ref{thm:A2}).
\item[(5)] \textbf{Toric varieties and local terms (\S\ref{sec:toric}).} The weight formula for toric varieties, a closed formula for general $\sigma$ (Theorem \ref{thm:T}), surfaces with unequal residues (Theorem \ref{thm:C2}). We prove the stalk formula for plane curves with quasi-homogeneous isolated singularities for $s\neq1$ (Theorem \ref{thm:local}).
\item[(6)] \textbf{Reduction theorem in general dimension and Theorem V (\S\ref{sec:reduction}).} The general-dimensional problem reduces to a condition (GEN$_Q$) on the quotient complex $Q$ (Theorem \ref{thm:R}). If all residues are non-integers and $U$ is affine, then $H^k(U;L_U)=0$ for $k\ne n$ (Theorem \ref{thm:V}).
\item[(7)] \textbf{Comparison with regular type (\S\ref{sec:regular}).} For regular type (flat symplectic fibre bundles) the dimensions in each degree are polynomial in $m$ (Corollary \ref{cor:L}); in the type-jumping case they are periodic.
\item[(8)] \textbf{Kodaira type (\S\ref{sec:kodaira}).} Embedding by generalized holomorphic sections is impossible; we prove a leaf-space version of Kodaira and the toric Poisson--Kodaira theorem (Proposition \ref{prop:one}, Theorems \ref{thm:K1leaf}, \ref{thm:PKtoric}).
\item[(9)] \textbf{RR-type formulas and the index conjecture (\S\ref{sec:RR}).} The Poincar\'e polynomial formula $P(t)=\chi(X)t^n+\sum_{k<n}h^k(U;L_U)(t^k-(-1)^kt^n)$ (Theorem \ref{thm:K}). Unification by an index theorem (weighted $b$-calculus) remains a conjecture (Conjecture \ref{conj:index}).
\end{enumerate}

\subsection{Status table}\label{subsec:ledger}
\begin{longtable}{>{\raggedright\arraybackslash}p{5.3cm}>{\raggedright\arraybackslash}p{2.6cm}>{\raggedright\arraybackslash}p{7.2cm}}
\caption{Status of the main claims (meaning of tags: \S\ref{subsec:tags})}\label{tab:ledger}\\
\toprule Claim & Status & Main hypotheses / remarks\\\midrule\endfirsthead
\toprule Claim & Status & Main hypotheses / remarks\\\midrule\endhead
\bottomrule\endfoot
\footnotesize Prop.\ \ref{prop:ident} (GH bundle $=$ Poisson module, identification of $\HH$) & \tC \tP \tQ & Depends on LGSX Thm.~4.19 and Mokri's theorem; not re-derived\\
\footnotesize Thm.\ \ref{thm:main1} (Poisson module $=$ flat log connection) & \tP & Any even dimension, SNC. The residue formula $c_1=-\sum r_i[D_i]$ is \tC\\
\footnotesize Thm.\ \ref{thm:serre} (Serre-type vanishing, transversal degree) & \tP & $X$ projective, $\sigma$ arbitrary. False in all degrees (Remark \ref{rem:total})\\
\footnotesize Thm.\ \ref{thm:dual} (duality, any dimension) & \tP \tV \tQ & Lemmas \ref{lem:D}, \ref{lem:E} machine-checked for $n\le4$. Consistency of Koszul signs between the $\dbar$-part and the $d$-part checked by hand only\\
\footnotesize Thm.\ \ref{thm:E} (total Euler characteristic) & \tP & HRR only. Independent of $c_1(E)$\\
\footnotesize Thm.\ \ref{thm:K1} (log version, $\chi(U)$) & \tP \tV & SNC, flat log connection\\
\footnotesize Thm.\ \ref{thm:main2} (main theorem for surfaces) & \tP \tV & $r_i\ne1$. Depends on Prop.\ \ref{prop:ident}\\
\footnotesize Thm.\ \ref{thm:B} (closed formula for surfaces) & \tP \tV & (H1)--(H4). (H4) checked in 8 examples; in general a condition on $V^1(U)$\\
\footnotesize Thm.\ \ref{thm:A2} (periodic stabilization) & \tP & Theorem \ref{thm:main2} and $X$ simply connected, $U$ affine\\
\footnotesize Thms.\ \ref{thm:weight}, \ref{thm:T}, \ref{thm:C2} (toric) & \tP \tV & (V) (follows from Thm.\ \ref{thm:serre} for $L_0^m$, $m\gg0$), (A), (GEN), (R)\\
\footnotesize Obs.\ \ref{obs:res} (formula at resonance $r_i=1$) & \tV \tK & Definition of $N$ copied from earlier records, \tQ. Fails in some cases\\
\footnotesize Thm.\ \ref{thm:local} (quasi-homogeneous local formula, $s\ne1$) & \tP \tQ & Link between local model and global module \tQ (in the earlier version: observation \tV \tK; updated, see the closing appendix)\\
\footnotesize Local formula (non-quasi-homogeneous, $s=1$) & \tK & Unproved\\
\footnotesize Conj.\ \ref{conj:refined}, Obs.\ \ref{obs:triple}, \ref{obs:nr}, Conj.\ \ref{conj:NR} & \tV \tK & Data fits. Triple points, non-reduced\\
\footnotesize Thm.\ \ref{thm:R} (reduction theorem) & \tP & Hypothesis (GEN$_Q$). Rewritten here without the filtration\\
\footnotesize Prop.\ \ref{prop:grQ} ((GEN$'$) $\Rightarrow$ (GEN$_Q$), $\mathrm{gr}\,Q$ for general $n$) & \tQ & General-$n$ local computation is a sketch. For $n=2$ see Lemma \ref{lem:Q}\\
\footnotesize Thm.\ \ref{thm:V} (Theorem V: $H^k(U;L_U)=0$) & \tP & SNC, $U$ affine, all $r_i\notin\Z$. Mixed case (some $r_i\in\Z$): \tK\\
\footnotesize Thm.\ \ref{thm:K} (RR-type Poincar\'e polynomial) & \tP \tV & Hypotheses of Thm.\ \ref{thm:R}, $U$ affine, $n$ even\\
\footnotesize Thm.\ \ref{thm:L}, Cor.\ \ref{cor:L} (regular type) & \tP \tQ & Hard Lefschetz. Identification of $\partial$ with Gauss--Manin is a sketch\\
\footnotesize Props.\ \ref{prop:one}, \ref{prop:leaf}, Thm.\ \ref{thm:K1leaf} (Kodaira type) & \tP & Impossibility and leaf-space version\\
\footnotesize Thm.\ \ref{thm:PKtoric} (toric Poisson--Kodaira) & \tP & Obs.\ \ref{obs:P2} ($\mathbb P^2$, $m=2$) is a numerical existence result \tV\\
\footnotesize Conj.\ \ref{conj:PK} (general Poisson--Kodaira) & \tK & $\mathbb P^2$, $m\ge3$ not computed\\
\footnotesize Conj.\ \ref{conj:index} (index conjecture) & \tK & Unproved\\
\end{longtable}

\subsection{Correspondence with the three treasures}
\begin{center}\small
\begin{tabular}{>{\raggedright\arraybackslash}p{2.4cm}>{\raggedright\arraybackslash}p{6.5cm}>{\raggedright\arraybackslash}p{6.0cm}}
\toprule Treasure & What holds (in some form) & What fails / is open\\\midrule
Riemann--Roch & Total Euler characteristic $(-1)^n\rk(E)\chi(X)$ (Thm.\ \ref{thm:E}), log version $\chi(U)$ (Thm.\ \ref{thm:K1}), Poincar\'e polynomial formulas (surfaces: Thm.\ \ref{thm:B}; general dimension: Thm.\ \ref{thm:K}), degree-wise polynomials for regular type (Cor.\ \ref{cor:L}) & A naive HRR depending on $c_1(E)$ fails. Unification by an index theorem is open (Conj.\ \ref{conj:index})\\\midrule
Serre vanishing & Vanishing in transversal degree $a>0$ ($m\gg0$, $X$ projective, $\sigma$ arbitrary, Thm.\ \ref{thm:serre}) & Vanishing in all degrees fails even for $\sigma=0$\\\midrule
Kodaira embedding & Leaf-space version (Thm.\ \ref{thm:K1leaf}), Poisson--Kodaira for toric varieties (Thm.\ \ref{thm:PKtoric}) & Embedding by generalized holomorphic sections is impossible (Prop.\ \ref{prop:one}). General Poisson--Kodaira is a conjecture (Conj.\ \ref{conj:PK}). Generalization to general generalized complex structures not attempted\\\bottomrule
\end{tabular}
\end{center}

\subsection{Three ideas running through the paper}
(i) \emph{Translation to flat connections.} On a $\beta$-transform structure the generalized Dolbeault complex is identified with $(\wedge^\bullet T\otimes E,\ \dbar+[\sigma,\cdot])$ \tC (Laurent-Gengoux--Sti\'enon--Xu \cite{LGSX1}), and generalized holomorphic bundles are Poisson modules. If $\sigma$ is log-symplectic along $D$, this becomes a flat logarithmic connection along $D$, and the residues are constant.
(ii) \emph{Separation of local and global.} Comparing the log de Rham complex with $Rj_*L_U$ (Deligne \cite{Del}) brings in the cohomology of the local system on $U=X\setminus D$; when residues are integers (resonance), the quotient complex contributes. The global part is $H^*(U;L_U)$ and $\chi(X)$; the local part is expressed by stalk formulas at the singularities of $D$.
(iii) \emph{Three behaviours.} For regular type the dimensions in each degree are polynomial in $m$; for log-symplectic type they are periodic; the total Euler characteristic is independent of $m$ in both cases (\S\ref{sec:regular}, \S\ref{sec:RR}).

\subsection{Meaning of the tags, symbol table, reading guide}\label{subsec:tags}
\paragraph{Status tags.} \tC Quoted from the literature (the extent of checking against the original is stated case by case). \tP Proved in this paper (dependence on quoted results is stated). \tV Computer verification by exact linear algebra (rational or finite field $\mathbb F_{32003}$, $\mathbb F_{2^{61}-1}$) and the like; not a proof. \tK Conjecture or outlook. \tQ Unverified point (quotation from memory, derivation only sketched, sign convention unchecked, etc.). When several tags are attached, they refer to the respective parts of the claim.

\paragraph{Reading guide.} \S\ref{sec:prelim}--\S\ref{sec:euler} are the foundations. \S\ref{sec:surface} is the complete result for surfaces; the toric case and local terms (\S\ref{sec:toric}) and general dimension (\S\ref{sec:reduction}) follow. \S\ref{sec:regular} (regular type) and \S\ref{sec:kodaira} (Kodaira type) are almost independent of each other; \S\ref{sec:RR} uses \S\ref{sec:reduction} and \S\ref{sec:euler}.

\paragraph{Symbol table.}
\begin{longtable}{>{\raggedright\arraybackslash}p{3.5cm}>{\raggedright\arraybackslash}p{11.3cm}}
\toprule Symbol & Meaning\\ \midrule\endhead
\bottomrule\endfoot
\footnotesize $X,\ n$ & Compact complex manifold, $\dim_\C X=n$ ($n$ even in the log-symplectic case)\\
\footnotesize $T,\ T^*,\ \Omega^p,\ K_X$ & Holomorphic tangent bundle, cotangent bundle, sheaf of holomorphic $p$-forms, $K_X=\Omega^n$\\
\footnotesize $\bT=T\oplus T^*$ & Generalized tangent bundle. Pairing $\langle X+\xi,Y+\eta\rangle=\tfrac12(\xi(Y)+\eta(X))$, Courant (Dorfman) bracket $[X+\xi,Y+\eta]=[X,Y]+\mathcal L_X\eta-\iota_Yd\xi$\\
\footnotesize $\Jc,\ L$ & Generalized complex structure and its $+i$-eigenbundle (a maximal isotropic subbundle of $\bT_\C$, a complex Lie algebroid). $L^*\cong\Lb$. $E$, $L_0$, $\mathcal O(m)$ etc.\ are also used for line bundles (distinguished by context)\\
\footnotesize $\sigma,\ \sigma^\sh,\ \Jc_\sigma$ & Holomorphic Poisson structure ($\sigma\in H^0(\wedge^2T)$, $[\sigma,\sigma]=0$), $\sigma^\sh:T^*\to T$, $\xi\mapsto\sigma(\xi,\cdot)$, and its $\beta$-transform $L_\sigma=T^{0,1}\oplus\{\sigma^\sh\xi+\xi:\xi\in T^{*1,0}\}$\\
\footnotesize $A=(T^*)_\sigma$ & Holomorphic Lie algebroid: anchor $\sigma^\sh$, bracket $[a,b]_\sigma=\mathcal L_{\sigma^\sh a}b-\mathcal L_{\sigma^\sh b}a-d\sigma(a,b)$\\
\footnotesize Poisson module & Holomorphic vector bundle $(E,\nabla)$ with a flat $A$-connection $\nabla$ ($=$ generalized holomorphic bundle, Prop.\ \ref{prop:ident})\\
\footnotesize $V$ & Holomorphic vector field (Poisson vector field) defined by $\nabla_\beta e=\beta(V)e$ for a local frame $e$\\
\footnotesize $C(A;E)$, $\HH^k(E)$ & $(\wedge^bT\otimes E,\ d_E=[\sigma,\cdot]+V\wedge)$ (Chevalley--Eilenberg complex of $A$), $\HH^k(E)=H^k(X;C(A;E))$\\
\footnotesize $E_1^{a,b}$ & $H^a(X,\wedge^bT\otimes E)\Rightarrow\HH^{a+b}$. We call $a$ the transversal degree\\
\footnotesize $P_E(t),\ [\cdots]$ & $P_E(t)=\sum_k\dim\HH^k(E)t^k$; in tables abbreviated as $[\dim\HH^0,\dots,\dim\HH^{2n}]$\\
\footnotesize $D=\sum D_i,\ U$ & Degeneracy divisor of $\sigma$ (reduced, SNC; for surfaces ordinary double points are allowed), $U=X\setminus D$. $D_T=\bigcap_{i\in T}D_i$, $D_{(n)}$ is the set of $n$-fold points\\
\footnotesize $T_D,\ B,\ \mathrm{DR}_{\log}(E)$ & $T_D=T_X(-\log D)$, $B=\Omega^1(\log D)=T_D^*$, log de Rham complex $(\Omega^\bullet(\log D)\otimes E,\nabla^{\log})$\\
\footnotesize log-symplectic & $\sigma\in H^0(\wedge^2T_D)$ and $\sigma^\sh:B\to T_D$ is an isomorphism\\
\footnotesize $r_i,\ L_U$ & Residue along $D_i$ (constant); the local system of flat sections on $U$ (rank 1, monodromy $e^{\mp2\pi\sqrt{-1}r_i}$ around $D_i$)\\
\footnotesize $Q$ & Quotient complex $C(A;E)/C(T_D;E)$, supported on $D$\\
\footnotesize $\mu,\ c,\ \rho_D,\ b_1$ & (surfaces) number of double points $\mu$, number of irreducible components $c$ of $D$, dimension $\rho_D$ of the span of the $[D_i]\in H^2(X;\C)$, $b_1=c-\rho_D$\\
\footnotesize $j_1,j_2,\ \mathcal F$ & Open embeddings $U\hookrightarrow X'\hookrightarrow X$, $\mathcal F=j_{2!}Rj_{1*}L_U$ (Theorem \ref{thm:main2})\\
\footnotesize $N,M,\Sigma,u_i$ & Toric: lattice $N\cong\Z^n$, dual $M$, fan and ray generators\\
\footnotesize $\rho,\ A,\ \chi,\ \Lambda$ & Toric: $\nabla=d+\rho$ ($\rho\in M_\C$), constant skew matrix $A$, character $\chi\in M$, set of characters in the weight formula\\
\footnotesize $h,\ (w_z,w_w),\ N_h,\ \delta$ & Defining function of a quasi-homogeneous isolated singularity, weights, degree, $\delta=N_h-w_z-w_w$ (\S\ref{subsec:local})\\
\footnotesize $\mu_p,\ \mathrm{Sp},\ \mathrm{lct}$ & Milnor number, spectrum, log canonical threshold $(w_z+w_w)/N_h$\\
\footnotesize $s$ & Residue parameter of the local model $\mathcal O h^s$\\
\footnotesize $Y\to M\to X$, $V^j$ & Flat symplectic fibre bundle; Gauss--Manin bundle of the fibre cohomology (\S\ref{sec:regular})\\
\footnotesize $\partial U$ & Boundary of a tubular neighbourhood of $D$ (closed real $(2n-1)$-manifold, $\chi(\partial U)=0$)\\
\end{longtable}

%==============================================================
\section{Preliminaries: generalized holomorphic line bundles and generalized Dolbeault cohomology}\label{sec:prelim}
%==============================================================
\subsection{The classical three treasures: statements and where the hypotheses enter}\label{subsec:classical}
To make the goal of the generalization precise, we recall the classical statements and the skeleton of their proofs. The content is standard and not new \tC.

\begin{itemize}
\item \textbf{HRR.} For a compact complex manifold $M$ and a holomorphic vector bundle $E$, $\chi(M,E)=\int_M\ch(E)\Td(T^{1,0}M)$. \src{Hirzebruch \cite{Hirz} (projective case); Atiyah--Singer \cite{AS} (general compact complex manifolds, index of the Dolbeault complex); Huybrechts \cite{Huy} \S5.1.}
\item \textbf{Serre vanishing.} If $X$ is projective, $\mathcal O(1)$ very ample and $\mathcal F$ coherent, there is $m_0$ with $H^q(X,\mathcal F(m))=0$ $(q>0)$ for $m\ge m_0$. \src{Serre \cite{Serre}; Hartshorne \cite{Hart} Thm.\ III.5.2.}
\item \textbf{Kodaira--Akizuki--Nakano vanishing.} If $M$ is compact K\"ahler and $E$ a positive holomorphic Hermitian line bundle, then $H^q(M,\Omega^p\otimes E)=0$ for $p+q>n$. \src{Kodaira \cite{Kod53}; Akizuki--Nakano \cite{AN}; Griffiths--Harris \cite{GH} Ch.~1 \S2.}
\item \textbf{Kodaira embedding theorem.} If $M$ is compact and $L$ a positive holomorphic line bundle, then $M$ is projective and for $m\gg0$ the map $M\to\mathbb P(H^0(M,L^m)^*)$ given by $|L^m|$ is an embedding. \src{Kodaira \cite{Kod54}; Griffiths--Harris Ch.~1 \S4; Huybrechts \S5.3.}
\end{itemize}

\begin{center}\small
\begin{tabular}{>{\raggedright\arraybackslash}p{2.3cm}>{\raggedright\arraybackslash}p{2.9cm}>{\raggedright\arraybackslash}p{4.4cm}>{\raggedright\arraybackslash}p{4.6cm}}
\toprule Classical theorem & Hypotheses & Why needed / where used in the proof & Fate in this paper\\\midrule
HRR & Compact, holomorphic structure & The Dolbeault complex is elliptic and Fredholm (index defined), $\dbar^2=0$. $\Td$ comes from the symbol, $E$ enters as $\ch(E)$ & The total Euler characteristic does not see $c_1$ (Thm.\ \ref{thm:E}). A refinement is needed (\S\ref{sec:RR})\\\midrule
Serre & Projective, ample, coherent & $\bigoplus\mathcal O(-N_i)\twoheadrightarrow\mathcal F$, \v Cech computation on $\mathbb P^N$, descending induction & Applied to the coherent sheaves $\wedge^bT\otimes E$: holds in transversal degree (Thm.\ \ref{thm:serre})\\\midrule
KAN & Compact, K\"ahler, positive & Hodge theory, K\"ahler identities $\Rightarrow$ BKN identity, positivity of the curvature term $[i\Theta_E,\Lambda]$ & The line-bundle-coefficient version is not written for generalized K\"ahler. Regular type: Pal--Poddar; Lie groupoids: Hu (\S\ref{subsec:compare})\\\midrule
Kodaira embedding & Compact, positive & Surjectivity of $H^0(L^m)\to L^m\otimes\mathcal O/\mathfrak m_x^2$ from KAN vanishing on the blow-up $\tilde M$ (many sections separate points and tangent vectors) & Generalized holomorphic sections are flat and cannot separate points (\S\ref{sec:kodaira})\\\bottomrule
\end{tabular}
\end{center}

\begin{rem}\tK (Suggestions for the generalization)
In the classical case four ingredients share the work: (1) Hodge theory (ellipticity), (2) K\"ahler identities $\Rightarrow$ BKN identity, (3) positivity of curvature, (4) abundance of global sections. In generalized complex geometry, (1) holds for the elliptic complex $(\wedge L^*\otimes E,d_L)$ (Gualtieri \cite{Gua11}). The line-bundle-coefficient versions of (2), (3) have not been written in the generalized K\"ahler framework (as far as we found). (4) fails in the classical form because generalized holomorphic sections are Casimir-type flat sections (Observation \ref{obs:casimir}).
\end{rem}

%--------------------------------------------------------------
\subsection{Generalized complex structures and the $\beta$-transform}
%--------------------------------------------------------------
\begin{defn}[Generalized complex structure, type]\tC\label{defn:GC}
A map $\Jc:\bT\to\bT$ with $\Jc^2=-1$ preserving the pairing is a generalized almost complex structure. It is integrable (a generalized complex structure) if its $+i$-eigenbundle $L\subset\bT_\C$ is closed under the Courant bracket. The number $\mathrm{type}(\Jc)_p=\mathrm{codim}_\C\,\pi_T(L_p)$ is the type; points where it is locally constant are called regular points. The type may jump. A complex structure $I$ has type $n$ ($L=T^{0,1}\oplus T^{*1,0}$); a symplectic structure has type $0$.
\src{Hitchin \cite{Hit03}; Gualtieri \cite{Gua11} \S3 (type Def.~3.5, canonical bundle Def.~3.7, generalized Calabi--Yau Def.~3.17, Picard group Def.~3.18; numbers to be checked against the original); Abouzaid--Boyarchenko \cite{AB}.}
\end{defn}
The generalized Dolbeault complex $(\wedge^\bullet L^*,d_L)$ is an elliptic complex, and its cohomology is finite-dimensional on a compact manifold (Gualtieri \cite{Gua11} Prop.~3.13, Cor.~3.14 \tC).

\begin{prop}[$\beta$-transform]\tC\label{prop:Jsigma}
Let $X$ be a complex manifold and $\sigma$ a holomorphic Poisson structure. (a) $L_\sigma=T^{0,1}\oplus\{\sigma^\sh\xi+\xi:\xi\in T^{*1,0}\}$ is an integrable maximal isotropic subbundle of $\bT_\C$, and the corresponding generalized complex structure $\Jc_\sigma$ has type $n-\mathrm{rk}_\C\sigma^\sh_p$ at $p$. In particular the type is $0$ where $\sigma$ is non-degenerate and $n$ where $\sigma=0$. (b) $(\Xi,\xi)\mapsto(\Xi+\sigma^\sh\xi,\xi)$ is an isomorphism of complex Lie algebroids $T^{0,1}\bowtie(T^{1,0})^*_\sigma\to L_\sigma$ (matched pair).
\src{(a) Gualtieri \cite{GuaBr} Prop.~5; Laurent-Gengoux--Sti\'enon--Xu (LGSX) \cite{LGSX1} Thm.~2.7, Rem.~2.8 ($\Jc_\pi$, $\pi=\pi_R+i\pi_I$). The type formula is our own organization and agrees with Gualtieri's examples \tQ. (b) LGSX Thm.~4.21, Lemma 4.22 (the factor 4 in $\Jc_{4\pi}$ comes from the normalization of $\pi=\pi_R+i\pi_I$). The choice of sign of the eigenvalue differs among references but they are exchanged by complex conjugation and nothing below is affected.}
\end{prop}

\begin{defn}[Generalized holomorphic bundle, Poisson module]\tC\label{defn:GH}
A complex vector bundle $E$ with a flat $L$-connection $\dbar_E:\Gamma(E)\to\Gamma(L^*\otimes E)$ ($\dbar_E^2=0$) is called a generalized holomorphic bundle. A holomorphic vector bundle $E$ together with a flat $A$-connection $\nabla:\Gamma(E)\to\Gamma(A^*\otimes E)$ ($\nabla_\xi(fs)=f\nabla_\xi s+(\sigma^\sh\xi)(f)s$, $[\nabla_a,\nabla_b]=\nabla_{[a,b]_\sigma}$, $a,b$ holomorphic 1-forms) is called a Poisson module.
\src{Gualtieri \cite{Gua11} \S3 (generalized holomorphic bundles are introduced in prose after Thm.~3.15, no numbered definition); Gualtieri \cite{GuaBr} \S4.1; LGSX Def.~4.15, Lemma 4.16.}
\end{defn}

\begin{defn}[Generalized Dolbeault cohomology $\HH$]\label{defn:HH}
For a Poisson module $E$, the hypercohomology $\HH^k(E):=H^k(X;C(A;E))$ of the Chevalley--Eilenberg complex $C(A;E)=(\wedge^\bullet T\otimes E,\ d_E)$ (the differential is $\C$-linear; for a local holomorphic frame $e$, $d_E=[\sigma,\cdot]+V\wedge$, where $V$ is determined by $\nabla_\beta e=\beta(V)e$) is called generalized Dolbeault cohomology. The filtration by $b$ gives a spectral sequence $E_1^{a,b}=H^a(X,\wedge^bT\otimes E)\Rightarrow\HH^{a+b}(E)$ (finite-dimensional if $X$ is compact). $P_E(t)=\sum_k\dim\HH^k(E)\,t^k$.
\end{defn}

\begin{prop}[Generalized holomorphic bundle $=$ Poisson module; Dolbeault model]\label{prop:ident}
(a) A generalized holomorphic bundle with respect to $\Jc_\sigma$ is the same as a holomorphic vector bundle with a Poisson module structure. \tC Gualtieri \cite{GuaBr} Prop.~8; LGSX Lemma 4.16. A sketch of a direct proof is given below \tP.
(b) From $L^*=T^{*0,1}\oplus T^{1,0}$ we get $\Gamma(\wedge^kL^*\otimes E)=\bigoplus_{a+b=k}A^{0,a}(\wedge^bT\otimes E)$, with $d_L=\dbar+(-1)^ad_A$. Hence the generalized Dolbeault cohomology $H^k(\Gamma(\wedge^\bullet L^*\otimes E),d_L)$ equals $\HH^k(E)$. \tC LGSX Thm.~4.19 ($H^*(A,E)=H^*(T^{0,1}\bowtie A^{1,0},E)$, where $H^*(A,E)$ is defined as sheaf cohomology in Def.~4.17). \tP The structure of the proof is shown below, but it depends on Mokri's theorem on matched pairs (\cite{Mok}) and LGSX Lemma 4.22, which we did not re-derive \tQ.
\end{prop}
\begin{proof}[Sketch of proof]
(a) ($\Leftarrow$) For a holomorphic frame $\{e_s\}$ put $\nabla_{\bar\Xi}\sum f_se_s=\sum(\bar\Xi f_s)e_s$ and $\nabla_\xi\sum f_se_s=\sum((\sigma^\sh\xi)f_s)e_s+f_s\nabla_\xi e_s$ (the unique extension forced by the Leibniz rule). Flatness: for the pair $(\bar\Xi,\bar\Upsilon)$ it is the integrability of the holomorphic structure, for $(\bar\Xi,\xi)$ it is that $\nabla_\xi$ maps holomorphic sections to holomorphic sections, and for $(\xi,\eta)$ it is the flatness of the Poisson module (the curvature is a tensor, so it extends to $C^\infty$). ($\Rightarrow$) The restriction to $T^{0,1}$ gives a holomorphic structure (Koszul--Malgrange); from flatness of mixed type, $\nabla_\xi$ maps holomorphic sections to holomorphic sections, and the Leibniz rule and flatness give the module condition.
(b) By the general theory of matched pairs (Mokri), the Chevalley--Eilenberg complex of $L\cong T^{0,1}\bowtie A^{1,0}$ is the total complex of a double complex in the $T^{0,1}$-direction ($\dbar$) and the $A^{1,0}$-direction ($d_A$). The two differentials anticommute under the Koszul sign $(-1)^a$, and the total complex is the Dolbeault resolution of $(\wedge^\bullet T\otimes E,d_A)$, so its cohomology is the hypercohomology. \qedhere
\end{proof}
Below, all results are statements about $\HH^*(E)=H^*(X;C(A;E))$; Proposition \ref{prop:ident}(b) is used wherever the identification with generalized Dolbeault cohomology is needed.

\begin{obs}[Generalized holomorphic sections are flat sections]\tP\label{obs:casimir}
$\HH^0(E)=\ker(d_E:H^0(E)\to H^0(T\otimes E))=\{s\in H^0(X,E):\nabla_\xi s=0\ \forall\xi\}$, i.e.\ flat (Casimir-type) sections. Near a point where $\sigma$ is non-degenerate, $\nabla$ corresponds to an ordinary flat connection via $\sigma^\sh$, and flat sections are locally constant.
\end{obs}
This is why generalized holomorphic sections cannot give an embedding of $X$ into projective space (\S\ref{sec:kodaira}).

%--------------------------------------------------------------
\subsection{The log-symplectic case: Poisson modules $=$ flat logarithmic connections}\label{subsec:logsymp}
%--------------------------------------------------------------
Let $X$ be a compact complex manifold ($\dim_\C X=n=2k$), $D\subset X$ a reduced divisor (simple normal crossings, SNC; for surfaces ordinary double points only are also allowed), $T_D=T_X(-\log D)$, $B=\Omega^1(\log D)$.
\begin{defn}[Log-symplectic]\tC\label{defn:logsymp}
If $\sigma\in H^0(\wedge^2T_D)$ and $\sigma^\sh:B\to T_D$ is an isomorphism, $\sigma$ is called log-symplectic along $D$ (equivalently, $\omega:=\sigma^{-1}\in H^0(\Omega^2(\log D))$ is non-degenerate). Then $\mathrm{div}(\sigma^k)=D$, i.e.\ $-K_X=\mathcal O(D)$. For surfaces this means that $\sigma\in H^0(-K_X)$ has a reduced zero divisor with only ordinary double points. Locally $\sigma=h\,\partial_1\wedge\partial_2$, $h=zu$ (smooth point) or $z_1z_2u$ (double point), $u$ a unit.
\src{Goto \cite{Goto}; Gualtieri--Li \cite{GL}; Pym \cite{Pym}, Lett.\ Math.\ Phys.\ 108 (2018) 573--632 (bibliographic data checked). Publication data of Goto and Gualtieri--Li: \tQ.}
\end{defn}

\begin{thm}[Poisson modules $=$ flat logarithmic connections]\tP\label{thm:main1}
Let $(X,D,\sigma)$ be log-symplectic (any even dimension, $D$ SNC) and $E$ a holomorphic line bundle with a Poisson module structure.
\begin{enumerate}
\item[(a)] For a local frame $e$, $\nabla_\beta e=\beta(V)e$ defines a holomorphic vector field $V$ with $[\sigma,V]=0$ (a Poisson vector field). Under the change of frame $e\mapsto ge$, $V$ changes by the Hamiltonian term $X_{\log g}$.
\item[(b)] $V$ is automatically logarithmic: $V\in T_D$. Hence the $A$-module structure extends uniquely to a $B$-module structure and, through $\sigma^\sh:B\cong T_D$, is the same as a flat connection with logarithmic poles $\nabla^{\log}:E\to\Omega^1(\log D)\otimes E$.
\item[(c)] The residue of $\nabla^{\log}$ along $D_i$ is a constant $r_i\in\C$, and $c_1(E)=-\sum_ir_i[D_i]$ in $H^2(X;\C)$.
\item[(d)] Near a point of $D$ there is a frame with $\nabla^{\log}e=\sum_ir_i\,\dlog z_i\,e$.
\item[(e)] The subcomplex $C(T_D;E)=(\wedge^bT_D\otimes E,d_E)\subset C(A;E)$ is isomorphic via $\wedge^b\sigma^\sh$ to the log de Rham complex $\mathrm{DR}_{\log}(E)$.
\end{enumerate}
\end{thm}
\begin{hyptab}
Log-symplectic ($\sigma^\sh:B\cong T_D$) & To identify the Lie algebroid structure $(T_D^*)_\sigma$ on $B$ with $T_D$ and treat Poisson modules as flat connections (with log poles) & Lemma \ref{lem:iso}, (b), (e)\\\midrule
$D$ reduced, SNC & To estimate the pole order from each factor of the local equation $f=z_1\cdots z_s$ independently & Lemma \ref{lem:log} ($V(f)\in(f)\Rightarrow V(z_i)\in(z_i)$)\\\midrule
$X$ compact & The relation between $c_1(E)$ and the residues is an equality of global cohomology classes & (c). That residues are constant is local\\\midrule
$[\sigma,\sigma]=0$ & $A$ is a Lie algebroid and $d_E^2=0$ & Lemma \ref{lem:flat}
\end{hyptab}

\begin{lem}[Flatness $\Leftrightarrow$ Poisson vector field]\tP\label{lem:flat}
For a local frame $e$ put $\Theta(\beta):=\nabla_\beta e/e$. Then $\nabla_a\nabla_be=(\sigma^\sh a)(\Theta b)e+\Theta(a)\Theta(b)e$. Hence flatness $[\nabla_a,\nabla_b]=\nabla_{[a,b]_\sigma}$ is equivalent to $(\sigma^\sh a)(\Theta b)-(\sigma^\sh b)(\Theta a)=\Theta([a,b]_\sigma)$, i.e.\ $\Theta$ is a 1-cocycle of the Lichnerowicz complex ($\Theta(\beta)=\beta(V)$, $[\sigma,V]=0$). For a surface with $\sigma=h\partial_1\wedge\partial_2$, $a=\Theta(dz_1)$, $b=\Theta(dz_2)$, the condition reads $h(a_1+b_2)=h_1a+h_2b$, i.e.\ $V(h)=h\,\mathrm{div}V$. The differential is $d_E(fe)=(-\sigma^\sh(df)+fV)e$ in degree $0$, and for $c=X\otimes e$, $X=p\partial_1+q\partial_2$,
\begin{equation}\label{eq:dE}
(d_Ec)(dz_1,dz_2)=h(p_1+q_2)-h_1p-h_2q+qa-pb.
\end{equation}
\end{lem}
\begin{lem}[$V$ is logarithmic]\tP\label{lem:log}
Let $f=z_1\cdots z_s$ be a local equation of $D$. From $L_V\sigma=0$ we get $L_V\sigma^k=0$. Since $\sigma^k=f'\partial_1\wedge\cdots\wedge\partial_n$ ($f'=f\cdot$unit), $V(f')-f'\,\mathrm{div}V=0$, so $V(f)\in(f)$. In $V(f)/f=\sum_iV(z_i)/z_i$ the $i$-th term has a pole only along $\{z_i=0\}$; since the sum is holomorphic, each $V(z_i)/z_i$ is holomorphic, i.e.\ $V(z_i)\in(z_i)$.
\end{lem}
\begin{lem}[$\sigma^\sh:B\to T_D$ is an isomorphism of Lie algebroids]\tP\label{lem:iso}
$\sigma\in\wedge^2T_D$ is a non-degenerate Poisson structure on the Lie algebroid $T_D$. That $\omega=\sigma^{-1}$ is $d_{T_D}$-closed follows from: $\omega$ is closed on $U$ (inverse of a Poisson structure), and a section of $\Omega^3(\log D)$ that vanishes on a dense open set vanishes. For a non-degenerate Poisson structure given by a closed 2-form, $\sigma^\sh$ preserves brackets, so $(T_D^*)_\sigma=B\cong T_D$.
\end{lem}
\begin{lem}[Residues, local normal form, $c_1$]\tP \tC\label{lem:res}
Write $\nabla^{\log}e=\theta e$, $\theta\in\Omega^1(\log D)$; $\theta$ is closed by flatness. The residue along $D_i$ is constant, as seen by restricting $d\theta$ to $D_i$. $\theta-\sum r_i\dlog z_i$ is a closed holomorphic 1-form $=dg$, and for the frame $e'=e^{-g}e$ we get $\nabla^{\log}e'=\sum r_i\dlog z_i\,e'$. $c_1(E)=-\sum r_i[D_i]$ follows from the Esnault--Viehweg residue formula (\tC \cite{EV} \S2; \cite{Del} Ch.~II. Sign check: for $\mathcal O(D)$ with the canonical section flat, the residue is $-1$ and $c_1=[D]$). The set of module structures with the same $c_1$ is a torsor under $H^0(\Omega^1(\log D))=\ker(\C^c\to H^2(X;\C))$ (Deligne), of dimension $c-\rho_D$. If $D$ is irreducible it is unique; if reducible, it is an affine family of residue vectors.
\end{lem}
\begin{proof}[Proof of Theorem \ref{thm:main1}]
(a) Lemma \ref{lem:flat}. (b) By Lemma \ref{lem:log}, $\Theta$ extends to $B$; flatness is an equality on a dense open set and so persists after the extension. Combine with Lemma \ref{lem:iso}. (c), (d) Lemma \ref{lem:res}. (e) Lemma \ref{lem:iso}.
\end{proof}

%==============================================================
\section{Serre-type vanishing and duality}\label{sec:serre}
%==============================================================
\subsection{Serre-type vanishing (transversal degree)}
The correspondence with classical Serre vanishing (coherent sheaf $\mathcal F$ and ample $\mathcal O(1)$) is as follows \tK (the correspondence is our own organization). The coherent sheaf $\mathcal F$ corresponds to each term $K^b=\wedge^bT\otimes E$ of the Lie algebroid complex; the ample $\mathcal O(1)$ corresponds to an ample holomorphic line bundle $L_0$ which is a Poisson module (the module structure is needed only so that $C(A;\cdot)$ is a complex, and plays no role in positivity); the conclusion $H^q(\mathcal F(m))=0\ (q>0)$ corresponds to $E_1^{a,b}=H^a(X,\wedge^bT\otimes E\otimes L_0^m)=0$ ($a>0$: the \emph{transversal (Dolbeault) degree}).

\begin{thm}[Serre-type theorem ($\beta$-transform)]\tP\label{thm:serre}
Let $X$ be a compact complex projective manifold ($\dim_\C X=n$), $\sigma$ an \emph{arbitrary} holomorphic Poisson structure (no non-degeneracy or log condition needed), $E$ a Poisson module (holomorphic vector bundle), and $L_0$ an ample holomorphic line bundle with a Poisson module structure. Then there is $m_0$ depending only on $X,\sigma,E,L_0$ such that for $m\ge m_0$:
\begin{enumerate}
\item[(a)] $H^a(X,\wedge^bT\otimes E\otimes L_0^m)=0$ for all $a>0$ and all $b$,
\item[(b)] $\HH^k(E\otimes L_0^m)\cong H^k\bigl(\Gamma(X,\wedge^\bullet T\otimes E\otimes L_0^m),\,d\bigr)$. In particular $\HH^k=0$ for $k>n$,
\item[(c)] $\HH^0(E\otimes L_0^m)=$ the space of all flat (Casimir) holomorphic sections.
\end{enumerate}
\end{thm}
\begin{hyptab}
$X$ projective, $L_0$ ample & The hypotheses of classical Serre vanishing & (a): apply Hartshorne III.5.2 to the finitely many ($b=0,\dots,n$) coherent sheaves $\wedge^bT\otimes E$\\\midrule
$L_0$, $E$ Poisson modules & $L_0^m$ gets the tensor product module structure so that $C(A;E\otimes L_0^m)$ is a complex & Definition of the complex $K^\bullet$. Unrelated to positivity\\\midrule
$[\sigma,\sigma]=0$ & $d^2=0$ & Same\\\midrule
(not needed) non-degeneracy of $\sigma$, normal crossings & --- & Not used in this theorem
\end{hyptab}
\begin{proof}
$L_0^m$ has the tensor product module structure (Leibniz rule). Apply Serre's theorem (Hartshorne III.5.2; if $L_0$ is merely ample, take $m_1$ with $L_0^{m_1}$ very ample and apply it also to the finitely many twists $\mathcal F\otimes L_0^j$, $0\le j<m_1$) to the $n+1$ coherent sheaves $\wedge^bT\otimes E$, and let $m_0$ be the maximum of the thresholds; this gives (a). Each $K^b$ is $\Gamma$-acyclic, so for the bounded complex $K^\bullet=C(A;E\otimes L_0^m)$ (the differential is $\C$-linear) $E_1^{a,b}=H^a(K^b)$ is concentrated in $a=0$ and $R\Gamma(K^\bullet)=\Gamma(K^\bullet)$ (lemma on acyclic resolutions). This is (b). Since $b\le n$, it vanishes for $k>n$. (c) is $H^0$ of the $\Gamma$-complex, by Observation \ref{obs:casimir}. \qedhere
\end{proof}
\begin{rem}[Vanishing in all degrees is false]\tP\label{rem:total}
For $\sigma=0$ (complex structure) $\HH^k(L)=\bigoplus_{a+b=k}H^a(X,\wedge^bT\otimes L)$, and even for $m\gg0$ we have $\HH^b\ne0$ $(b=0,\dots,n)$ because $H^0(\wedge^bT\otimes L^m)\ne0$. For $\mathbb{CP}^2$ at $m=0$ one gets $[1,8,10,0,0]$. Only transversal degree $a>0$ and total degree $>n$ vanish; ``Serre vanishing in all degrees'' does not hold.
\end{rem}
\begin{rem}[Effective thresholds \tV]
From Bott's formula, the condition $E_1^{a>0}=0$ is $m\ge-2$ for $\mathbb{CP}^2$ ($E=\mathcal O(m)$; at $m=-3$ the three terms $(2,0),(1,1),(0,2)$ are all non-zero) and $m\ge-1$ for $\mathbb{CP}^1\times\mathbb{CP}^1$ ($L_0=\mathcal O(1,1)$; $m=-2$ is the exceptional case $K_X$). Both agree with the range in the computed tables (\S\ref{sec:surface}) where degrees $\ge3$ vanish.
\end{rem}
\begin{rem}[Relation to Pal--Poddar]
Thm.~11.12(2) of Pal--Poddar \cite{PP} is an analogous vanishing in regular type, but the groups are not the same (e.g.\ presence of the leafwise de Rham factor). The comparison is made in \S\ref{subsec:PPcompare} (Proposition \ref{prop:PPcompare}). This theorem does not assume regular type or an orbifold leaf space, and includes the type-jumping case, but requires the underlying complex manifold to be projective (only for the $\beta$-transform case). For a general generalized complex manifold there is no way to define ``the $E_1$ term of a coherent filtration $K^b$'', and this is the main obstacle to generalization \tK.
\end{rem}

%--------------------------------------------------------------
\subsection{Serre duality}
%--------------------------------------------------------------
\begin{thm}[Serre duality for generalized holomorphic line bundles]\tP\label{thm:dual}
Let $X$ be any compact complex manifold ($\dim_\C X=n$), $\sigma$ any holomorphic Poisson structure (neither K\"ahler, projective, log-symplectic nor non-degenerate is needed; $\sigma=0$ is included), and $L$ a line bundle which is a Poisson module. Then for all $k$
\[\HH^k(L)^*\;\cong\;\HH^{2n-k}\bigl(L^\vee\otimes K_X^{2}\bigr).\]
$L^\vee\otimes K_X^2$ has the tensor product module structure: for the dual frame $e^\vee$ of a local frame $e$ and the frame $\omega^2$ for $\omega=dz_1\wedge\cdots\wedge dz_n$, $V'=-V+2V_K=-V-2M$. Here the canonical module structure of $K_X$ is $V_K^k=\sum_i\partial_i\sigma^{ki}$, and $M=-V_K$ ($M^j=\sum_i\partial_i\sigma^{ij}$) is the modular vector field with respect to $\omega$.
\end{thm}
\begin{hyptab}
$X$ compact & Closed range and finite-dimensionality for elliptic complexes; no boundary terms in Stokes & Step (4) (identification of the dual), step (3)\\\midrule
$[\sigma,\sigma]=0$, $L$ a Poisson module & $D^2=0$, $D'^2=0$ ($L'$ is also a module) & To treat them as complexes\\\midrule
(not needed) K\"ahler / non-degenerate / log & Ellipticity is checked directly from the principal symbol, and the transpose is an equality of local operators, so no global primitive is needed & Steps (1), (3)
\end{hyptab}
\begin{proof}
(0) Dolbeault model: $\mathrm{Tot}^k(L)=\bigoplus_{a+b=k}A^{0,a}(\wedge^bT\otimes L)$, $D=\dbar+d_K$ ($d_K(\phi\otimes Pe)=(-1)^a\phi\otimes d(Pe)$, $d=[\sigma,\cdot]+V\wedge$). $D^2=0$ follows from $\dbar^2=0$, $d^2=0$, $\dbar d_K+d_K\dbar=0$ (Proposition \ref{prop:ident}(b)).

(1) Ellipticity: for a real cotangent vector $\xi\ne0$ the principal symbol of $D$ is the exterior product with $v(\xi)=(\xi^{0,1},\sigma^\sh\xi^{1,0})\in L^*=T^{*0,1}\oplus T^{1,0}$, and $(\wedge L^*,v\wedge)$ is exact if and only if $v\ne0$. If $v=0$ then $\xi^{0,1}=0$, and $\xi$ real implies $\xi=0$. Hence the complex is elliptic at all points, including those where $\sigma$ degenerates (a direct check of Gualtieri \cite{Gua11} Prop.~3.13 in this case). The same holds for $L'=L^\vee\otimes K^2$.

(2) Perfect pairing: integration of $E^b(L)\times E^{n-b}(L')\to\wedge^nT\otimes K^2=K$ and $A^{0,a}\times A^{0,n-a}\to A^{0,n}(K)=A^{n,n}$ gives pointwise perfect pairings $\mathrm{Tot}^k(L)\times\mathrm{Tot}^{2n-k}(L')\to\C$.

(3) Transpose: by Lemma \ref{lem:E} below, the formal transpose of the $d$-part equals $d'$ (up to sign). The right-hand side of Lemma \ref{lem:E} is a total derivative of the form $\partial_iF^i\,\omega\wedge d\bar z^{1\cdots n}=\pm d(F^i\iota_{\partial_i}\omega\wedge d\bar z^{1\cdots n})$ (the $\dbar$-component vanishes after the product with $d\bar z^{1\cdots n}$), and its integral vanishes by Stokes. For the $\dbar$-part, the classical Serre duality computation gives $\langle\dbar\alpha,\beta\rangle=-(-1)^{a+b}\langle\alpha,\dbar\beta\rangle$, for the $d$-part $\langle d_K\alpha,\beta\rangle=-(-1)^{a+b}\langle\alpha,d'_K\beta\rangle$, and we checked by hand that the two signs agree (not machine-checked \tQ). Since $D^t$ is a local operator, equalities in coordinate charts give globally $\langle D\alpha,\beta\rangle=-(-1)^{\deg\alpha}\langle\alpha,D'\beta\rangle$.

(4) Conclusion: $D,D'$ are elliptic complexes on a closed manifold and are transposes of each other (up to sign) with respect to the perfect pairing. By Hodge theory for elliptic complexes, $\HH^k$ is finite-dimensional with closed range, $\HH^k(L)^*$ is the cohomology of the transpose complex (with distributional coefficients), and by elliptic regularity it equals the one with smooth coefficients, $\HH^{2n-k}(L')$. \qedhere
\end{proof}
\begin{lem}[Lemma D: top degree of the Schouten bracket]\tP \tV\label{lem:D}
For $T=\sum_j(-1)^ju_j\,\partial_1\wedge\cdots\widehat{\partial_j}\cdots\wedge\partial_n\in\wedge^{n-1}T$, $[\sigma,T]_{\rm top}=-\sum_{i,j}\partial_i(\sigma^{ij}u_j)+2\sum_jM^ju_j$ and $(M\wedge T)_{\rm top}=\sum_jM^ju_j$.
\end{lem}
\begin{lem}[Lemma E: transpose identity]\tP \tV\label{lem:E}
For $P\in\wedge^bT$, $Q\in\wedge^{n-b-1}T$ (smooth coefficients), $d=[\sigma,\cdot]+V\wedge$, $d'=[\sigma,\cdot]+V'\wedge$, $V'=-V-2M$, and $\langle P,Q\rangle:=(P\wedge Q)_{\rm top}$,
\[\langle dP,Q\rangle+(-1)^b\langle P,d'Q\rangle=-\sum_i\partial_i\Bigl(\sum_j\sigma^{ij}u_j\Bigr),\qquad T=P\wedge Q=\sum_j(-1)^ju_j\partial_{\hat j}.\]
Proof: the Leibniz rule $[\sigma,P\wedge Q]=[\sigma,P]\wedge Q+(-1)^bP\wedge[\sigma,Q]$ ($\sigma$ has even degree); the $V$-terms give $V\wedge P\wedge Q+(-1)^bP\wedge(-V)\wedge Q=0$; the $M$-term is $-2M\wedge T$. The claim follows from Lemma D.
\tV For $n=2,3,4$ the identity was checked in exact arithmetic with random polynomial and Laurent monomial data and arbitrary bivector fields $\sigma$ (the identity itself holds without the Poisson condition) (script \texttt{duality\_ndim\_exact.py}). The constant $c=-2$ and the sign $(-1)^b$ were found independently of the derivation by solving a system of equations (\texttt{duality\_ndim.py}). $[\sigma,M]=0$ and $d^2=0$ ($V=M$) were also checked.
\end{lem}
\begin{rem}[Consistency checks \tV]
(i) $\sigma=0$: classical Serre duality ($\wedge^bT^*=\wedge^{n-b}T\otimes K$). (ii) $\mathbb{CP}^2$, $K^2=\mathcal O(-6)$: $m\leftrightarrow-m-6$; $\mathbb{CP}^1\times\mathbb{CP}^1$, $K^2=\mathcal O(-4,-4)$: $m\leftrightarrow-m-4$ (agrees with the tables of \S\ref{sec:surface}). (iii) In the log-symplectic case the residue of the canonical module structure of $K_X$ is $1$ (because the section $\sigma^k$ of $\wedge^nT=K_X^{-1}$ is flat and vanishes to first order along $D$), so the residue of $L'$ is $2-r_i$. This is the origin of the symmetry $r\leftrightarrow2-r$ in the formula of the main theorem for surfaces (in any even dimension).
\end{rem}

%==============================================================
\section{Euler characteristic}\label{sec:euler}
%==============================================================
\subsection{The total Euler characteristic is a topological invariant}
\begin{thm}[Euler characteristic theorem]\tP\label{thm:E}
Let $X$ be a compact complex manifold ($\dim_\C X=n$), $\sigma$ any holomorphic Poisson structure, and $E$ a holomorphic vector bundle of rank $\rk E$ with a flat $A$-connection. The hypercohomology $\HH^k(E)$ of the complex $K^\bullet=(\wedge^\bullet T\otimes E,\ d_E)$ (the differential is $\C$-linear) is finite-dimensional and
\[\sum_k(-1)^k\dim\HH^k(E)=\sum_b(-1)^b\chi(X,\wedge^bT\otimes E)=(-1)^n\,\rk(E)\,\chi(X).\]
\end{thm}
\begin{hyptab}
$X$ compact & Each $H^a(X,K^b)$ is finite-dimensional and HRR can be used & (i), (ii)\\\midrule
$K^b$ coherent & The Euler characteristic can be computed by HRR and the descending induction works & (i): induction over truncation triangles\\\midrule
Flat $A$-connection & $d_E^2=0$ (a complex) & Definition of the hypercohomology. Not used in computing the Euler characteristic itself\\\midrule
(not used) non-degeneracy of $\sigma$, positivity of $E$ & --- & Do not appear in the proof. Hence no dependence on $c_1(E)$
\end{hyptab}
\begin{proof}
(i) The differential is $\C$-linear but not $\mathcal O$-linear; still the truncations $K^{\ge b}$ are subcomplexes and there are triangles $K^{\ge b+1}\to K^{\ge b}\to K^b[-b]\to$. Each $K^b$ is coherent and $X$ is compact, so $H^*(X;K^b)$ is finite-dimensional. By descending induction on $b$, all $H^*(X;K^{\ge b})$ are finite-dimensional and $\chi(K^{\ge b})=\chi(K^{\ge b+1})+(-1)^b\chi(X,K^b)$. Hence $\sum_k(-1)^k\dim\HH^k=\sum_b(-1)^b\chi(X,\wedge^bT\otimes E)$.
(ii) By HRR, $\chi(X,\wedge^bT\otimes E)=\int\ch(E)\ch(\wedge^bT)\Td(T)$. With Chern roots $x_i$ of $T$, $\sum_b(-1)^b\ch(\wedge^bT)=\prod(1-e^{x_i})$, $\Td(T)=\prod x_i/(1-e^{-x_i})$, and $(1-e^x)/(1-e^{-x})=-e^x$, so the integrand is $(-1)^n\ch(E)e^{c_1(T)}c_n(T)$. Since $c_n(T)$ already has top degree, only the degree-0 part $\rk E$ of $\ch(E)e^{c_1(T)}$ survives, and by Gauss--Bonnet we get $(-1)^n\rk E\int c_n(T)=(-1)^n\rk E\,\chi(X)$. \qedhere
\end{proof}
\tV For $\mathbb P^n$ ($n\le6$), $E=\mathcal O(m)$ ($m=-12,\dots,12$) we recomputed $\sum_b(-1)^b\chi(\mathbb P^n,\wedge^bT(m))=(-1)^n(n+1)$ from Bott's formula and confirmed it (script \texttt{roundD\_check.py}).

\begin{rem}[Meaning \tP]\label{rem:Emeaning}
The left-hand side does not depend on $c_1(E)$. Example: $\sigma=0$, $\mathbb P^1$, $E=\mathcal O(m)$: $\chi(\mathcal O(m))-\chi(\mathcal O(m+2))=(m+1)-(m+3)=-2$. The classical $\chi(X,E)=\int\ch(E)\Td(X)$ corresponds to the Euler characteristic of the ``column'' $b=0$, $H^\bullet(X,E)$, not of the total complex. Hence a Riemann--Roch formula for the \emph{total} Euler characteristic of the generalized Dolbeault complex is topological and is not a generalization of classical HRR. Riemann--Roch as a ``polynomial in $m$'' becomes constant, and the dependence on $m$ appears only in the Poincar\'e polynomial (dimensions degree by degree) (\S\ref{sec:surface}, \S\ref{sec:regular}, \S\ref{sec:RR}).
\end{rem}
\begin{rem}[General generalized complex structures \tK]
For a general $\Jc$ and a flat $L$-connection $E$, the Atiyah--Singer theorem gives $\mathrm{ind}(\wedge L^*\otimes E)=\int\ch(E)\cdot S$ ($S$ is the class of the symbol of $(\wedge L^*,d_L)$). By Chern--Weil and flatness of the $L$-connection, $\ch_j(E)$ should lie in the kernel of $H^{2j}_{dR}(M)\to H^{2j}(L)$, so only the rank should matter. Not proved.
\end{rem}

%--------------------------------------------------------------
\subsection{Euler characteristic of the log de Rham complex and wall-crossing}
%--------------------------------------------------------------
\begin{thm}[Log version of the Euler characteristic]\tP\label{thm:K1}
Let $X$ be a compact complex manifold, $D\subset X$ an SNC divisor, and $L$ a line bundle with a flat logarithmic connection $\nabla:L\to L\otimes\Omega^1_X(\log D)$. Then the hypercohomology $\mathbb H^k(X;\mathrm{DR}_{\log}(L))$ is finite-dimensional and
\[\sum_k(-1)^k\dim\mathbb H^k(X;\mathrm{DR}_{\log}(L))\;=\;\int_Xc_n\bigl(T_X(-\log D)\bigr)\;=\;\chi(U)\;=\;\sum_{T}(-1)^{|T|}\chi(D_T),\]
independent of $L$ and of the residues $r_i$.
\end{thm}
\begin{hyptab}
$X$ compact & $H^q(X;\Omega^p(\log D)\otimes L)$ is finite-dimensional, HRR is available & (1), (2)\\\midrule
$D$ SNC & $\Omega^p(\log D)$ is locally free, and $\Omega^\bullet(\log D)\simeq Rj_*\C_U$ (log Poincar\'e lemma) can be used & (3) (to get the topological value for $L=\mathcal O$)\\\midrule
$\nabla$ flat log connection & $\mathrm{DR}_{\log}(L)$ is a complex (residues are constant since the rank is 1) & (1)\\\midrule
(not needed) residues $r_i\notin\Z_{\ge1}$ & --- & Needed only for the identification $\mathrm{DR}_{\log}(L)\simeq Rj_*L_U$, not for the Euler characteristic\\\midrule
(not needed) log-symplectic $\sigma$ & --- & Used only for the relation with the generalized Dolbeault complex (Theorem \ref{thm:main1}(e))
\end{hyptab}
\begin{proof}
(1) The $E_1$ term of the filtration by $\Omega^p(\log D)\otimes L$ is $H^q(X;\Omega^p(\log D)\otimes L)$, which is finite-dimensional. Hence $\chi(\mathrm{DR}_{\log}(L))=\sum_p(-1)^p\chi(X;\Omega^p(\log D)\otimes L)$.
(2) By HRR this equals $\int_X\bigl[\sum_p(-1)^p\ch(\Omega^p(\log D))\bigr]\ch(L)\Td(T_X)$. For the rank $n$ bundle $F=\Omega^1(\log D)$, $\sum_p(-1)^p\ch(\wedge^pF)=c_n(F^\vee)\,\Td(F^\vee)^{-1}$ starts in degree $n$. Hence only the constant term $1$ of $\ch(L)\Td(T_X)\Td(F^\vee)^{-1}$ matters, and the result is $\int_Xc_n(T_X(-\log D))$, independent of $L$.
(3) For $L=\mathcal O$ ($\nabla=d$), the log Poincar\'e lemma (Deligne \cite{Del}) $\Omega^\bullet(\log D)\simeq Rj_*\C_U$ shows that the hypercohomology is $H^*(U;\C)$, and the Euler characteristic of (1) is $\chi(U)$. By (2) this equals $\int c_n(T_X(-\log D))$. Finally $\chi(U)=\sum_T(-1)^{|T|}\chi(D_T)$ follows from $\chi_c=\chi$ on the open manifold $U$ (Poincar\'e duality, real even dimension) and additivity of $\chi_c$. \qedhere
\end{proof}
\tV For $X=\mathbb P^2$, $D$ a smooth curve of degree $d=1,\dots,7$, and $L=\mathcal O(a)$ ($a=-12,\dots,12$) we computed $\sum_p(-1)^p\chi(\Omega^p(\log D)\otimes L)$ directly, and in all 175 cases it equals $3-3d+d^2=\chi(U)$ (script \texttt{roundK\_check.py}; re-run and confirmed).

\begin{cor}[Wall-crossing identity]\tP\label{cor:wall}
If $L_U$ is a rank one local system, then $\chi_c(U;L_U)=\chi(U;L_U)=\chi(U)$. This follows from the Euler characteristic of the triangle $j_!L_U\to Rj_*L_U\to i_*i^*Rj_*L_U$: the stalk of $Rj_*L_U$ over $D_T$ ($|T|\ge1$) is $H^*((S^1)^{|T|};\cdot)$, which has Euler characteristic 0 (the same reason as $\chi(\partial U)=0$). Hence the Euler characteristic does not change when $j_!$ and $Rj_*$ are exchanged as the residues cross integers.
\end{cor}
\begin{rem}[Relation to Theorem \ref{thm:E}]\tP
Theorem \ref{thm:E} and this theorem are two faces of the same phenomenon ``the Euler characteristic does not see the higher terms of $\ch(L)$''. The Euler characteristic of the generalized Dolbeault complex is $(-1)^n\chi(X)$, that of its subcomplex $B\simeq\mathrm{DR}_{\log}(L)$ is $\chi(U)$, and the difference $(-1)^n\chi(X)-\chi(U)$ is taken up by the Euler characteristic of the quotient complex $Q$ (\S\ref{sec:reduction}).
\end{rem}

%==============================================================
\section{Closed formula for log-symplectic surfaces}\label{sec:surface}
%==============================================================
\subsection{Setting and main theorem}
Let $X$ be a compact complex surface, $\sigma\in H^0(-K_X)$ with reduced zero divisor $D=\bigcup D_i$ (irreducible components $D_i$, singularities only ordinary double points, their number being $\mu$), and $E$ a line bundle with a Poisson module structure. The residues $r_i$ are given by Theorem \ref{thm:main1}(c) (constant along each branch; the two branches at a double point are regarded as different branches, and for an irreducible nodal curve the residues of the two branches are equal). Put
\[I_*=\{i:r_i\in\Z_{\le0}\},\quad I_!=\{i:r_i\notin\Z_{\le0}\},\quad X'=X\setminus\bigcup_{i\in I_!}D_i,\quad j_1:U\hookrightarrow X',\ j_2:X'\hookrightarrow X,\]
\[\mathcal F:=j_{2!}\,Rj_{1*}L_U\]
($Rj_*$ along the components in $I_*$, extension by zero along those in $I_!$).

\begin{thm}[Main theorem for surfaces]\tP\label{thm:main2}
In the above setting assume $r_i\ne1$ for all branches. Then
\[\HH^k(E)\;\cong\;H^k(X;\mathcal F)\ \oplus\ \C^{\mu}\quad(\text{the second term only for }k=2).\]
In particular (i) if all $r_i\notin\Z_{>0}$ then $\HH^*(E)=H^*(U;L_U)\oplus\C^\mu[-2]$; (ii) if all $r_i\notin\Z_{\le0}$ then $\HH^*(E)=H^*_c(U;L_U)\oplus\C^\mu[-2]$; (iii) the total Euler characteristic is $\chi(U)+\mu=\chi(X)$ (independent of $E$, consistent with Theorem \ref{thm:E}).
\end{thm}
\begin{hyptab}
Surface, $\sigma\in H^0(-K_X)$, $D$ reduced with only ordinary double points & To use the local normal forms $h=zu$ or $z_1z_2u$ and compute the quotient complex $Q$ explicitly & Lemma \ref{lem:Q}, Prop.\ \ref{prop:moser}, Lemma \ref{lem:conn}\\\midrule
$r_i\ne1$ (all branches) & If $r_i=1$ the differential $-u(1-r)p$ of $Q$ vanishes, the cohomology sheaves of $Q$ become infinite-dimensional and the reduction breaks down & Lemma \ref{lem:Q} ($r_i=1$ is treated separately for $E=K_X$, Prop.\ \ref{prop:KX})\\\midrule
Residues not positive integers / not non-positive integers & Condition for the comparison of the log complex with $Rj_*L_U$, $j_!L_U$ (Deligne) to be a quasi-isomorphism & Lemma \ref{lem:stalk}\\\midrule
$X$ compact & $R\Gamma(X;j_!L_U)=R\Gamma_c(U;L_U)$ and the Euler characteristic computation & (ii), (iii)\\\midrule
$-K_X=\mathcal O(D)$ & Gives $\chi(D)=\mu$ ($\omega_D=\mathcal O_D$, $p_a(D)=1$) & (iii)
\end{hyptab}

\subsection{Proof}
\begin{lem}[Stalks of the log complex and Deligne's comparison]\tP \tC\label{lem:stalk}
In the local frame of Theorem \ref{thm:main1}(d), the stalk of $\mathrm{DR}_{\log}$ at a point of $D_I$ ($|I|=s$) is the tensor product of the Koszul complex of the commuting operators $z_i\partial_i+r_i$ ($i\in I$) with the ordinary de Rham complex in the remaining variables. On a monomial $z^a$, $z_i\partial_i+r_i=a_i+r_i$ (a scalar), so the stalk cohomology is $\wedge^*\C^s$ (dimension $2^s$, Euler characteristic 0) if and only if $r_i\in\Z_{\le0}$ for all $i\in I$, and 0 otherwise. On the other hand the stalk of $Rj_*L_U$ at the same point is $\wedge^*\C^s$ if and only if $e^{2\pi\sqrt{-1}r_i}=1$ for all $i\in I$ ($r_i\in\Z$). Hence
\begin{itemize}
\item The natural map $\mathrm{DR}_{\log}\to Rj_*L_U$ is a quasi-isomorphism when no $r_i$ is a positive integer (Deligne \cite{Del} Ch.~II \S3, number \tQ; here confirmed by direct computation). Conversely, if $r_i\in\Z_{>0}$ for some $i$, at a smooth point of $D_i$ the stalks are $0$ versus $\wedge^*\C$, so it is not a quasi-isomorphism.
\item The extension by zero $j_!L_U\to\mathrm{DR}_{\log}$ (extension by zero of a flat section is holomorphic) is a quasi-isomorphism when no $r_i$ lies in $\Z_{\le0}$.
\item In the mixed case, $\mathrm{DR}_{\log}\simeq Rj_{1*}L_U$ on $X'$, and the stalks vanish along $\bigcup_{I_!}D_i$, so $\mathrm{DR}_{\log}\simeq j_{2!}j_2^*\mathrm{DR}_{\log}\simeq\mathcal F$.
\end{itemize}
Division by the coefficients $1/(a_i+r_i)$ does not spoil convergence, since $|a_i+r_i|\ge c>0$ except for finitely many $a$.
\end{lem}
\begin{lem}[Quotient complex $Q$]\tP \tV\label{lem:Q}
$Q:=C(A;E)/C(T_D;E)$ is supported on $D$ and lives only in degrees 1, 2. $Q^1=(T/T_D)\otimes E$, $Q^2=(\wedge^2T/\wedge^2T_D)\otimes E$.
\begin{itemize}
\item Smooth point ($h=zu$): $Q^1=\partial_z\otimes\mathcal O_D$, $Q^2=\partial_z\wedge\partial_w\otimes\mathcal O_D$, $d_Q(p)=-u(1-r)\,p|_D$. Hence acyclic if $r\ne1$.
\item Double point ($h=z_1z_2u$): $Q^1=\mathcal O_{D_1}\partial_1\oplus\mathcal O_{D_2}\partial_2$, $Q^2=\mathcal O_D$, $d_Q(p,q)=-u[(1-r_1)z_2p+(1-r_2)z_1q]\bmod z_1z_2$. If $r_1,r_2\ne1$ it is injective with image the maximal ideal $\mathfrak m_D$, so $H^1(Q)_p=0$ and $H^2(Q)_p=\mathcal O_D/\mathfrak m_D=\C$.
\end{itemize}
Hence for $r_i\ne1$, $Q\simeq\bigoplus_{\text{double points }p}\C_p[-2]$ in the derived category.
\end{lem}
\begin{proof}
Restrict formula \eqref{eq:dE} to $z=0$ (to $z_1=0$ or $z_2=0$ at a double point). The sign convention for residues is $V=-\sigma^\sh\theta$ ($\nabla^{\log}e=\theta e$); at a smooth point $a|_D=\theta_2h$, $b|_D=-ru$ ($\theta=r\,\dlog z+\theta_2\,dw$), and at a double point $a=r_2z_1u$, $b=-r_1z_2u$. For $p\partial_z$, the value of $h(p_1+q_2)-h_1p-h_2q+qa-pb$ at $q=0$, $z=0$ is $-up+rup=-u(1-r)p$. The double point is computed similarly, and injectivity of $d_Q(p,q)$ follows by restricting to $z_1=0$: $(1-r_1)z_2p|_{z_1=0}=0\Rightarrow p=0$. The image is $(z_1,z_2)/(z_1z_2)=\mathfrak m_D$, so the cokernel is $\C$. \qedhere
\end{proof}
\tV With random polynomials $u,g,p,q$ and residues, we checked $d_E^2=0$ and the formula for $d_Q$ in exact arithmetic, 20 times each at smooth and double points (\texttt{ncd\_cochain\_check.py}, re-run and confirmed).

\begin{prop}[Logarithmic Moser normal form at a double point]\tP\label{prop:moser}
Near a double point $p$ let $\omega=\sigma^{-1}=u^{-1}\dlog z_1\wedge\dlog z_2$ and $c:=u(0)$. There is a biholomorphism $\phi$ preserving each branch of $D$ and fixing $p$ with $\phi^*\omega=c^{-1}\dlog z_1\wedge\dlog z_2$. The constant $c$ is a local invariant and cannot in general be made 1.
\end{prop}
\begin{proof}
Put $\omega_0=\dlog z_1\wedge\dlog z_2$, $k=1/u-1/c$ ($k(0)=0$), $\omega_t=(c^{-1}+tk)\omega_0$. Then $\omega_t$ is closed and non-degenerate near $p$ for each $t\in[0,1]$. For $k=\sum k_{ab}z_1^az_2^b$,
\[\alpha=\sum_{a\ge1,b\ge0}\tfrac{k_{ab}}{a}z_1^az_2^b\,\dlog z_2-\sum_{b\ge1}\tfrac{k_{0b}}{b}z_2^b\,\dlog z_1\]
is a convergent log 1-form with $d\alpha=k\omega_0$. Since $X\mapsto\iota_X\omega_t$ is an isomorphism $T_D\to\Omega^1(\log D)$, there is a unique $X_t\in T_D$ with $\iota_{X_t}\omega_t=-\alpha$; it is tangent to each branch and vanishes at $p$. Its flow $\phi_t$ preserves $D$, and $\tfrac d{dt}\phi_t^*\omega_t=\phi_t^*(d\iota_{X_t}\omega_t+k\omega_0)=\phi_t^*(-d\alpha+k\omega_0)=0$. Hence $\phi_1^*\omega=\phi_1^*\omega_1=\omega_{t=0}=c^{-1}\omega_0$. \qedhere
\end{proof}
\begin{lem}[Vanishing of the connecting morphism]\tP \tV\label{lem:conn}
The component $e_p\in H^3_p(X;\mathrm{DR}_{\log})$ at a double point $p$ of the connecting morphism of the exact triangle $C(T_D;E)\to C(A;E)\to Q\to C(T_D;E)[1]$ vanishes when $(r_1,r_2)\ne(1,1)$.
\end{lem}
\begin{proof}[Sketch]
By Proposition \ref{prop:moser} and Lemma \ref{lem:res} we may assume $\sigma=c\,z_1z_2\partial_1\wedge\partial_2$ and $\nabla^{\log}e=(r_1\dlog z_1+r_2\dlog z_2)e$. All complexes are graded by the torus weight $\gamma\in\Z^2$ and the differentials preserve weights ($\lambda_i=\gamma_i+r_i$). Local cohomology can be computed from the Laurent tail $H^2_p(\mathcal O)=\mathrm{span}\{z^{-a}:a\ge1\}$. (a) $H^3_p(\mathrm{DR}_{\log})$ occurs only in the weight $\lambda_1=\lambda_2=0$ ($\gamma=-r$), and is non-zero only when $r_1,r_2$ are positive integers. (b) $H^2_p(Q)=\C$ sits in the single weight $\gamma=(-1,-1)$ (coming from the extra class $H^1_p(\mathcal O_D)$ in the local cohomology of $\mathcal O_D$), and $(-1,-1)=-r$ only if $r_1=r_2=1$. Since the connecting morphism preserves weights, $e_p=0$ for $(r_1,r_2)\ne(1,1)$. All maps are diagonal on monomials, so there is no convergence issue.
\tV \texttt{ncd\_local\_weights.py}: for weights $|\gamma_i|\le8$ and 78 pairs of residues $r_i\in\{-3,\dots,3,\pm\tfrac12,\tfrac13,-\tfrac23,\tfrac52\}$ (including resonance $r_i=1$), we checked in exact arithmetic the stalk cohomology, the cone of local cohomology, and the injectivity of $H(\mathrm{LOG})\to H(\mathrm{FULL})$ (re-run and confirmed). There were 0 violations for $r_i\neq1$, and for $r=(1,1)$ the connecting morphism is non-zero. The weight-by-weight hand computations behind (a), (b) are omitted for reasons of space; their basis is this computation. \qedhere
\end{proof}
\begin{proof}[Proof of Theorem \ref{thm:main2}]
By Theorem \ref{thm:main1}(e) and Lemma \ref{lem:stalk}, $C(T_D;E)\simeq\mathrm{DR}_{\log}\simeq\mathcal F$. By Lemmas \ref{lem:Q}, \ref{lem:conn}, $C(A;E)\simeq\mathrm{DR}_{\log}\oplus\bigoplus_p\C_p[-2]$ (in $D^b(\C_X)$). Apply $R\Gamma(X;\cdot)$. (i), (ii): $\mathcal F=Rj_*L_U$ or $j_!L_U$, and since $X$ is compact $R\Gamma(X;j_!L_U)=R\Gamma_c(U;L_U)$. (iii): $\chi(X;\mathcal F)=\sum_S\chi_c(S)\chi(\mathcal F_x)$, and the stalks over $D$ have Euler characteristic 0, so this is $\chi_c(U)=\chi(U)=\chi(X)-\chi(D)$. Since $D$ is anticanonical, $\omega_D=\mathcal O_D$ and $\chi(\mathcal O_D)=0$ gives $\sum_j(2-2\tilde g_j)=2\mu$ ($\tilde D_j$ the normalizations), so $\chi_{\rm top}(D)=\mu$ and $\chi(U)+\mu=\chi(X)$. (Alternative proof in the SNC case: $\chi(U)=\int c_2(T_X(-\log D))$, Theorem \ref{thm:K1}.) \qedhere
\end{proof}

\begin{prop}[$r_i=1$: $E=K_X$]\tP\label{prop:KX}
For $E=K_X$ (canonical module structure, $r_i=1$ on all branches) we have $\wedge^bT\otimes K_X=\Omega^{2-b}$ and $E_1^{a,b}=H^{2-b,a}(X)$ (total degree $2+a-p$, $p=2-b$). If $h^{p,q}(X)=0$ for $p\ne q$, then $E_1$ is concentrated in total degree 2, the differential raises total degree by 1 and so degenerates, hence $\HH^k=0$ ($k\ne2$) and $\dim\HH^2=\sum_ph^{p,p}=\chi(X)$ (3 for $\mathbb{CP}^2$, 4 for $\mathbb{CP}^1\times\mathbb{CP}^1$).
\end{prop}

%--------------------------------------------------------------
\subsection{Closed formula for the Poincar\'e polynomial}
%--------------------------------------------------------------
In the setting of Theorem \ref{thm:main2} we assume in addition: (H1) $H^1(X;\C)=0$; (H2) $U$ is affine; (H3) all components have the same residue $r_i=r\ne1$; (H4) if $r\notin\Z$ then $H^0(U;L_U)=H^1(U;L_U)=0$. Put $b_1=c-\rho_D$ ($c$ the number of components, $\rho_D$ the dimension of the span of the $[D_i]\in H^2(X;\C)$), $\chi_U=\chi(X)-\mu$, $b_2=\chi_U-1+b_1$.

\begin{thm}[Riemann--Roch-type closed formula (surfaces)]\tP\label{thm:B}
Under the above assumptions
\begin{align*}
P_E(t)&=1+b_1t+(b_2+\mu)t^2 &&(r\in\Z_{\le0}),\\
P_E(t)&=(b_2+\mu)t^2+b_1t^3+t^4 &&(r\in\Z_{\ge2}),\\
P_E(t)&=\chi(X)\,t^2 &&(r\notin\Z).
\end{align*}
Equivalently, $P_E(t)=\chi(X)t^2+[r\le0]\,\Delta(t)+[r\ge2]\,t^4\Delta(1/t)$ with $\Delta(t)=(1+t)(1+(b_1-1)t)$.
\end{thm}
\begin{hyptab}
(H1) $H^1(X;\C)=0$ & From the Gysin sequence, $H^1(U)=\ker(\C^c\to H^2(X))$ & $b_1(U)=c-\rho_D$\\\midrule
(H2) $U$ affine & Artin vanishing (\cite{SGA4}) & $b_3=b_4=0$, hence $b_2=\chi_U-1+b_1$\\\midrule
(H3) equal residues & $\mathcal F$ is either $Rj_*$ or $j_!$ & Theorem \ref{thm:main2}(i), (ii)\\\midrule
(H4) $r\notin\Z$ and $H^0(U;L)=H^1(U;L)=0$ & So that $H^*(U;L)=[0,0,\chi_U]$ & The row $r\notin\Z$. For nodal curves in $\mathbb{CP}^2$, $\pi_1(U)$ is abelian (Fulton \cite{Ful}, Deligne \cite{DelFG}) \tC\\\midrule
$r_i\ne1$ & Hypothesis of Theorem \ref{thm:main2} & Lemma \ref{lem:Q}
\end{hyptab}
\begin{proof}
Substitute into Theorem \ref{thm:main2}(i), (ii) the values $H^*(U)=[1,b_1,b_2,0,0]$ ($r\in\Z_{\le0}$, trivial local system), $H_c^*(U)=[0,0,b_2,b_1,1]$ (Poincar\'e duality, $r\in\Z_{\ge2}$), $H^*(U;L)=[0,0,\chi_U,0,0]$ ($r\notin\Z$: (H2), (H4), $\chi(U;L)=\chi(U)$) and add $\mu t^2$. The equivalent form follows from $\chi(X)=\chi_U+\mu=1-b_1+b_2+\mu$. \qedhere
\end{proof}
\begin{rem}[Properties \tP]
(i) $P_E(-1)=\chi(X)$ (consistent with Theorem \ref{thm:E}, $n=2$). (ii) The duality $P_{2-r}(t)=t^4P_r(1/t)$ is evident from the formulas (a formula-level consequence of Theorem \ref{thm:dual}). (iii) $\mu$ and $b_2$ are absorbed in $\chi(X)t^2$, and the jump term $\Delta$ depends only on $b_1$. (iv) For $E=L_0^m$, the quasi-polynomial in $m$ given by classical RR is the constant $\chi(X)$, and the dependence on $m$ appears only through integrality (period $N$) and the sign of the residue $r=-ma$.
\end{rem}
\begin{rem}[On (H4) \tQ]\label{rem:G2}
We verified (H4) in 8 examples (4 nodal curves in $\mathbb{CP}^2$ and 4 in $\mathbb{CP}^1\times\mathbb{CP}^1$). For smooth $(2,2)$ curves and two $(1,1)$ curves in $\mathbb{CP}^1\times\mathbb{CP}^1$ we cited Timmerscheidt \cite{Tim} and Arapura \cite{Ara} (the theorem on torsion translates) (applicability not fully checked \tQ). In general $H^1(U;L_\chi)=0\Leftrightarrow\chi\notin V^1(U)$ (the first characteristic variety), which can fail for non-nodal $D$ (e.g.\ three concurrent lines) (Libgober, Arapura).
\end{rem}

%--------------------------------------------------------------
\subsection{Examples and periodic stabilization}
%--------------------------------------------------------------
\begin{ex}[$\mathbb{CP}^2$ and a smooth cubic]\tP \tV\label{ex:cp2}
Let $F$ be a smooth homogeneous cubic on $\mathbb C^3$ and $\Pi^{ij}=\epsilon_{ijk}\partial_kF$. $\Pi$ is a homogeneous Poisson structure whose contraction with the Euler vector field vanishes, and defines a holomorphic Poisson structure $\sigma\in H^0(\wedge^2T)=H^0(\mathcal O(3))$ on $\mathbb{CP}^2$ vanishing to first order along $D=\{F=0\}$ (standard construction; identification of the source \tQ; the vanishing along $D$ is checked by local computation).
\emph{Modules:} for every $m\in\Z$, $\mathcal O(m)$ has exactly one Poisson module structure, with residue $r=-m/3$. Indeed, Theorem \ref{thm:main1}(c) gives $r=-m/3$ ($[D]=3H$). Existence: the Atiyah class of $\mathcal O(m)$ vanishes in $H^1(\Omega^1(\log D))$ iff $c_1\in\C[D]$ ($H^0(\mathcal O_D)\to H^1(\Omega^1)\to H^1(\Omega^1(\log D))$), and $m[H]\in\C\cdot3[H]$. The curvature is an element of $H^0(\Omega^2(\log D))=\C\omega$ with residue $dr=0$, hence flat. Uniqueness: the difference lies in $H^0(\Omega^1(\log D))=0$.
\emph{Result:} the dimensions $[\dim\HH^0,\dots,\dim\HH^4]$ of $\HH^*(\mathcal O(m))$ are
\begin{center}\small
\begin{tabular}{lll}\toprule
$m$ & dimensions & Interpretation\\\midrule
$m\in3\Z_{\ge0}$ & $[1,0,2,0,0]$ & $H^*(U;\C)$. $\HH^0=\C F^{m/3}$ is a Casimir section\\
$m\in3\Z_{\le-6}$ & $[0,0,2,0,1]$ & $H^*_c(U;\C)$\\
otherwise (including $m=-3$) & $[0,0,3,0,0]$ & $m\notin3\Z$: $H^*(U;L_m)$. $m=-3$ is $E=K_X$ (Prop.\ \ref{prop:KX})\\\bottomrule
\end{tabular}
\end{center}
The total Euler characteristic is always $3=\chi(\mathbb{CP}^2)$. This follows from Theorem \ref{thm:B}: $c=\rho_D=1$, $b_1=0$, $\mu=0$, $\chi_U=3$, $b_2=2$. (H4) holds since $U=\{F=1\}/\mu_3$ and $\{F=1\}$ (the Milnor fibre, a bouquet of 8 copies of $S^2$) is simply connected, so $\pi_1(U)=\Z/3$ and $H^0=H^1=0$ for non-trivial characters (\tC Milnor \cite{Milnor}, details \tQ). $r=1$ occurs only for $m=-3$.
\end{ex}
\tV The table above agrees, for $m=-15,\dots,15$ and Fermat and Hesse cubics, with the direct computation (for $m\ge-2$ the complex of global sections, for $m\le-4$ a \v Cech computation with Laurent monomials) of the differential $d_1$ of the spectral sequence $E_1^{a,b}=H^a(\wedge^bT(m))$ by exact linear algebra modulo the prime $32003$ (scripts \texttt{poisson.py, table.py, neg\_check.py, predict.py}). Even for $m\gg0$, $\HH^2$ remains of dimension 2 or 3, so Serre vanishing in all degrees fails.

\begin{ex}[$\mathbb{CP}^1\times\mathbb{CP}^1$]\tP \tV\label{ex:p1p1}
$\sigma=\Pi_F$, $F\in H^0(\mathcal O(2,2))$ ($-K=\mathcal O(2,2)$). When $D$ is smooth, $\mathcal O(a,b)$ has a Poisson module structure only if $a=b=m$, the residue is $-m/2$ and the structure is unique (same argument as above). Exact computation for $m=-9,\dots,8$ (\texttt{p1p1.py}):
\begin{center}\small
\begin{tabular}{lll}\toprule
$D$ & $m$ even $\ge0$ / even $\le-4$ & $m$ odd, $m=-2$\\\midrule
smooth $(2,2)$ curve & $[1,0,3,0,0]$ / $[0,0,3,0,1]$ & $[0,0,4,0,0]$\\
two transversal $(1,1)$ curves ($\mu=2$) & $[1,1,4,0,0]$ / $[0,0,4,1,1]$ & $[0,0,4,0,0]$\\
four lines $x_0x_1y_0y_1$ ($\mu=4$) & $[1,2,5,0,0]$ / $[0,0,5,2,1]$ & $[0,0,4,0,0]$\\
$(1,1)+(1,0)+(0,1)$ ($\mu=3$) & $[1,1,4,0,0]$ / $[0,0,4,1,1]$ & $[0,0,4,0,0]$\\\bottomrule
\end{tabular}
\end{center}
Period 2, Serre duality $m\leftrightarrow-m-4$ and $\chi=4$ always hold. For reducible $D$ the Poisson module structure is not unique ($H^0(\Omega^1(\log D))\ne0$), and the table is the computation for the symmetric structure with equal residues on all components.
\end{ex}
\tV We compared Theorems \ref{thm:main2}, \ref{thm:B} with exact computations for smooth cubic, nodal cubic ($\mu=1$), conic plus secant ($\mu=2$), triangle ($\mu=3$) in $\mathbb{CP}^2$ ($m=-15..15$) and the four examples in $\mathbb{CP}^1\times\mathbb{CP}^1$ above ($m=-9..8$), in total 196 values of $m$, with 0 discrepancies (\texttt{ncd\_global\_check.py}). The duality $r\leftrightarrow2-r$ of the closed formula was checked on a lattice of 73 rational values of $r$ (\texttt{serre\_rr\_check.py}, re-run and confirmed). Cuspidal cubic, conic plus tangent, three concurrent lines and non-reduced $D$ are outside the scope of this section and are treated in \S\ref{subsec:local}.

\begin{thm}[Periodic stabilization (log-symplectic surfaces)]\tP\label{thm:A2}
Let $X$ be a simply connected projective surface, $\sigma$ log-symplectic along $D$ (anticanonical, only ordinary double points), $L_0$ a line bundle with a Poisson module structure whose residues are $\rho_i=-a_i$ ($a_i\in\mathbb Q_{>0}$, constant along each component), and $U=X\setminus D$ affine ($D$ ample). Let $N$ be the least common multiple of the denominators of the $a_i$. Then for $m\ge1$
\[\HH^k(L_0^m)=0\ (k\ge3),\qquad\HH^*(L_0^m)=H^*(U;L_U^{(m)})\oplus\C^\mu[-2].\]
Here $L_U^{(m)}$ is the local system with monodromy $e^{-2\pi\sqrt{-1}ma_i}$ around $D_i$. In particular $\HH^*(L_0^m)$ depends only on $m\bmod N$.
\end{thm}
\begin{hyptab}
Residues $-ma_i<0$ & $r_i\ne1$ and not positive integers & Theorem \ref{thm:main2}(i): $\mathcal F=Rj_*L_U$\\\midrule
$U$ affine ($D$ ample) & Artin vanishing $H^k(U;L)=0\ (k>2)$ & Vanishing in degrees $\ge3$ (Theorem \ref{thm:serre}(b) gives the same conclusion from projectivity and ampleness, making affineness unnecessary)\\\midrule
$X$ simply connected & The abelianization of $\pi_1(U)$ is generated by meridians & The local system is determined by $m\bmod N$\\\midrule
$a_i\in\mathbb Q_{>0}$ & $c_1(L_0)$ is a positive rational combination of the boundary components & Periodicity
\end{hyptab}
\begin{proof}
Since $r_i=-ma_i<0$, Theorem \ref{thm:main2}(i) applies and $\HH^*(L_0^m)=H^*(U;L_U^{(m)})\oplus\C^\mu[-2]$. By Artin's vanishing theorem the terms in degree $k\ge3$ vanish. If $X$ is simply connected, $\pi_1(U)$ is generated by meridians, so a rank one local system is determined by the monodromies $e^{-2\pi\sqrt{-1}ma_i}$, which depend only on $m\bmod N$. \qedhere
\end{proof}
\tV Period 3 for $\mathbb{CP}^2$ ($L_0=\mathcal O(1)$, $a=1/3$) and period 2 for $\mathbb{CP}^1\times\mathbb{CP}^1$ ($L_0=\mathcal O(1,1)$, $a=1/2$), $m\ge1$. Confirmed in all 8 computed examples (\texttt{serre\_rr\_check.py}, re-run and confirmed).

%==============================================================
\section{Toric varieties and local terms}\label{sec:toric}
%==============================================================
\subsection{Toric log-symplectic structures and the weight formula}
Let $X$ be a smooth projective toric variety ($N=\Z^n$, $n=2k$, $M=\mathrm{Hom}(N,\Z)$), $u_i\in N$ the ray generators, and $D=\sum D_i$ the toric boundary (SNC). In the logarithmic frame $E_j=z_j\partial_j$ take a constant non-degenerate skew matrix $A$ and put $\sigma=\tfrac12\sum A_{ij}E_i\wedge E_j$ (torus invariant and log-symplectic along $D$). Give $L=\mathcal O(\sum a_iD_i)$ a flat log connection $d+\rho$ ($\rho\in M_\C$). The residues are $r_i=\langle\rho,u_i\rangle-a_i$ and $c_1(L)=-\sum r_i[D_i]$. Every module has this form ($H^0(\Omega^1(\log D))=M_\C$, Lemma \ref{lem:res}) \tP. For $\chi\in M$ put
\begin{gather*}c_i(\chi)=\langle\chi,u_i\rangle+a_i,\quad C(\chi)=\{i:c_i(\chi)=-1\},\\ W(\chi)=\mathrm{span}(u_i:i\in C(\chi)),\quad s(\chi)=\dim W(\chi).\end{gather*}

\begin{thm}[Weight formula]\tP \tV\label{thm:weight}
Assume $E_1^{a,b}=H^a(X,\wedge^bT\otimes L)=0\ (a>0)$ (e.g.\ $L=L_0^m$, $m\gg0$, $L_0$ an ample Poisson module; Theorem \ref{thm:serre}). Then
\[P_{\HH}(t)=\sum_{\chi\in\Lambda}t^{s(\chi)}(1+t)^{n-s(\chi)},\qquad\Lambda=\{\chi\in M:\ c_i(\chi)\ge-1\ \forall i,\ A(\chi+\rho)\in W(\chi)\}.\]
\end{thm}
\begin{hyptab}
Toric, $\sigma$ torus invariant, constant $A$ & The complex of global sections is graded by torus characters $\chi\in M$ and each weight is finite-dimensional linear algebra & Determination of the weight-$\chi$ subcomplex $G^b_\chi$\\\midrule
$E_1^{a>0}=0$ & $\HH^*$ equals the cohomology of the complex of global sections (argument of Theorem \ref{thm:serre}(b)) & The first equality\\\midrule
$A$ non-degenerate & $\sigma^\sh:B\cong T_D$ holds and the log de Rham differential $(\chi+\rho)\wedge$ goes to $v_\chi\wedge$ ($v_\chi=A(\chi+\rho)$) & Lemma \ref{lem:iso}, Theorem \ref{thm:main1}(e)
\end{hyptab}
\begin{proof}
$z^\chi w$ ($w\in\wedge^bN_\C$) is a holomorphic section of $\wedge^bT\otimes L$ iff $c_i(\chi)\ge-1$ along each $D_i$ and, if $c_i(\chi)=-1$, $u_i\wedge w=0$ ($w$ is divisible by $E_i=z_i\partial_i$). So the weight-$\chi$ part is $G^b_\chi=\{w\in\wedge^bN_\C:u_i\wedge w=0\ (i\in C(\chi))\}$ (empty if some $c_i<-1$), and with a generator $\omega_W$ of $\wedge^sW$ we have $G^\bullet_\chi=\omega_W\wedge\wedge^{\bullet-s}(N_\C/W)$. The differential corresponds to $v_\chi\wedge$ via $\wedge\sigma^\sh$. If $v_\chi\in W$ the differential is 0 and the Poincar\'e polynomial is $t^s(1+t)^{n-s}$; otherwise it is the Koszul complex of $v_\chi\bmod W\neq0$ and acyclic. \qedhere
\end{proof}
\tV The closed formula and direct linear algebra (\texttt{toric\_hh.py}) agree in 60 random cases ($\mathbb P^2$, $\mathbb P^1\times\mathbb P^1$, $\mathbb P^4$, $\mathbb P^2\times\mathbb P^2$, $(\mathbb P^1)^4$).

\begin{cor}\tP\label{cor:weight}
(a) If $E_1^{a>0}=0$, the total Euler characteristic is the number of weights with $s(\chi)=n$ (the vertices) $=\chi(X)$ (consistent with Theorem \ref{thm:E}).
(b) Origin of the contributions: put $x=\chi+\rho$. $s=0$ is $\chi=-\rho$, $\rho\in M$ and all $r_i\le0$, giving $(1+t)^n=H^*(U)$ (the $Rj_*$ case). $s=n$ always satisfies the condition and gives the vertex weights, one $t^n$ each. $0<s<n$ requires $x\in A^{-1}W$ and $\langle x,u_i\rangle=r_i-1$ ($i\in C$), and has solutions only at ``resonances'' on the integral lattice.
(c) For $s=n-1$ and all $r_i=1$ ($i\in C$) the condition holds automatically, and a class $t^{n-1}(1+t)$ appears for each of the $\deg(L|_C)+1$ lattice points on the toric curve $C$ (corresponding to the resonance $r_i=1$ in dimension 2).
\end{cor}

%--------------------------------------------------------------
\subsection{Toric surfaces: unequal residues}
%--------------------------------------------------------------
Let $\Delta_\Sigma$ be the simplicial complex whose vertices are the rays of the fan $\Sigma$ (simplices are the sets of rays of cones of $\Sigma$), $V_S$ the full subcomplex spanned by a vertex set $S$, and $\tilde h^q(V_S)$ the dimension of reduced cohomology (with the convention $\tilde h^{-1}=1$ and the others $0$ for $V_\emptyset=\emptyset$).

\begin{thm}[Toric cohomology of $\mathrm{DR}_{\log}$]\tP \tC\label{thm:C1}
Let $X=X_\Sigma$ be a complete smooth toric variety, $L=\mathcal O(\sum a_iD_i)$, $\nabla=d+\rho$, and $S=\{i:\langle\rho,u_i\rangle-a_i>0\}$. Then
\[\sum_k\dim H^k(X;\mathrm{DR}_{\log}(L))\,t^k=\begin{cases}(1+t)^n\sum_q\tilde h^{q-1}(V_S)\,t^q&(\rho\in M),\\ 0&(\rho\notin M).\end{cases}\]
\end{thm}
\begin{hyptab}
Torus-invariant structure & The complex is graded by weights $\chi\in M$ and each weight can be computed independently & First decomposition in the proof\\\midrule
Completeness & Hypercohomology is finite-dimensional and the \v Cech complex can be computed with finitely many affine opens & \v Cech complex\\\midrule
$\Omega^1(\log D)\cong\mathcal O\otimes M$ & Splitting as a product with a Koszul complex (K\"unneth) & The weight-$\chi$ part has the form $\check C^\bullet_\chi\otimes(\wedge M_\C,(\chi+\rho)\wedge)$\\\midrule
Cox--Little--Schenck weight formula & $H^a(X;L)_\chi$ is expressed by reduced cohomology of simplicial complexes & The last identification
\end{hyptab}
\begin{proof}
The \v Cech complex of a $T$-equivariant affine cover splits by weight $\chi$, and the weight-$\chi$ part is the tensor product complex $\check C^\bullet(L)_\chi\otimes\wedge^\bullet M_\C$ (differential $d_{\check C}\otimes1\pm1\otimes(\chi+\rho)\wedge$). By K\"unneth, $H^k(X;\mathrm{DR}_{\log})_\chi=\bigoplus_{a+b=k}H^a(X;L)_\chi\otimes H^b(\wedge M_\C,(\chi+\rho)\wedge)$. The Koszul complex is acyclic for $\chi+\rho\neq0$ and equals $\wedge^bM_\C$ of dimension $\binom nb$ for $\chi+\rho=0$ (which requires $\rho\in M$). Hence the total is 0 if $\rho\notin M$, and for $\rho\in M$ only $\chi=-\rho$ contributes, giving $(1+t)^n\sum_ah^a(X;L)_{-\rho}t^a$. Finally by Cox--Little--Schenck \cite{CLS} (Thm.~9.1.3 \tC \tQ): $H^a(X,\mathcal O(\sum a_iD_i))_m\cong\tilde H^{a-1}(V_{D,m};\C)$, where $V_{D,m}$ is the full subcomplex spanned by the rays with $\langle m,u_i\rangle<-a_i$. For $m=-\rho$ the condition is $\langle\rho,u_i\rangle>a_i$, i.e.\ $i\in S$. \qedhere
\end{proof}
\begin{thm}[Toric surfaces, unequal residues]\tP \tV\label{thm:C2}
Let $X$ be a complete smooth toric surface, $A$ a non-degenerate $2\times2$ skew matrix, and $r_i\ne1$ for all $i$. Then
\[P_{\HH}(t)=\chi(X)\,t^2+[\rho\in M]\,(1+t)^2\sum_q\tilde h^{q-1}(V_S)\,t^q,\qquad S=\{i:r_i\ge2\}.\]
In particular for $S=\emptyset$ (all $r_i\le0$) it is $(1+t)^2+\chi(X)t^2$, and when $S$ is all rays it is $t^2(1+t)^2+\chi(X)t^2$.
\end{thm}
\begin{hyptab}
Toric surface, $\sigma$ log-symplectic & Hypotheses of Theorem \ref{thm:main2}: $D$ reduced with only ordinary double points (double points are torus-fixed points), $\mu=\chi(X)$ & $\HH^*=H^*(X;\mathcal F)\oplus\C^\mu[-2]$\\\midrule
$r_i\ne1$ (all components) & Hypothesis of Lemma \ref{lem:Q} (acyclicity of the quotient complex $Q$) & Same\\\midrule
$\mathrm{DR}_{\log}\simeq\mathcal F$ & The mixed case of Lemma \ref{lem:stalk} & Reduction to $H^*(X;\mathcal F)=H^*(X;\mathrm{DR}_{\log})$ and application of Theorem \ref{thm:C1}
\end{hyptab}
\begin{proof}
By Theorem \ref{thm:main2}, $\HH^*=H^*(X;\mathcal F)\oplus\C^\mu[-2]$, and for a toric surface $\mu=$ number of maximal cones $=\chi(X)$. By Lemma \ref{lem:stalk}, $H^*(X;\mathcal F)=H^*(X;\mathrm{DR}_{\log})$, and Theorem \ref{thm:C1} gives, for $\rho\in M$ (i.e.\ all $r_i\in\Z$), $(1+t)^2\sum\tilde h^{q-1}(V_S)t^q$ with $S=\{r_i\ge1\}=\{r_i\ge2\}$ (as $r_i\ne1$). \qedhere
\end{proof}
\tV For $\mathbb{CP}^2$ (triangle, $m=-9..6$, residues $\{-2..3,\pm\tfrac12\}$) 755 rows, and for $\mathbb{CP}^1\times\mathbb{CP}^1$ (four lines) 729 rows, there are 0 discrepancies with exact linear algebra of the weight-decomposed \v Cech hypercohomology (\texttt{roundC\_toric\_cech.py, roundC\_p2\_sweep*.py, roundC\_p1p1\_sweepA.py}). For $\mathbb{CP}^1\times\mathbb{CP}^1$ with $S=$ two opposite lines ($V_S$ two points), the term $t(1+t)^2$ predicted by the formula does appear. For equal residues $r_i=r$, $S$ is empty or all rays, in agreement with the three cases of Theorem \ref{thm:B} (e.g.\ triangle with $r\le0$: $(1+t)^2+3t^2=[1,2,4]$). \tC These large tables are transcribed from the earlier records and were not re-run while writing this paper \tQ.

%--------------------------------------------------------------
\subsection{Toric varieties of any even dimension: general $\sigma$}
%--------------------------------------------------------------
\begin{thm}[Toric, general $\sigma$]\tP \tV\label{thm:T}
Let $X$ be a smooth projective toric variety ($\dim X=n$ even), $\sigma$ a torus-invariant log-symplectic structure (constant matrix $A$), $L=\mathcal O(\sum a_iD_i)$, $\nabla=d+\rho$, and assume
\begin{itemize}
\item[(V)] $H^a(X,\wedge^bT\otimes L)=0$ ($a>0$, all $b$).
\item[(A)] $L(D):=L\otimes\mathcal O(\sum D_i)$ is ample.
\item[(GEN)] $\mathrm{Pf}(A)\ne0$ and the numbers $A_{ij}\ (i<j)$ are linearly independent over $K=\mathbb Q(\rho_1,\dots,\rho_n)$.
\item[(R)] $r_i\ne1$ for all $i$.
\end{itemize}
Then $P_{\HH}(t)=[\ \rho\in M\ \text{and}\ r_i\le0\ (\forall i)\ ]\,(1+t)^n+\chi(X)\,t^n$. Equivalently: if $L$ has no positive integral residue, $\HH^*(L)\cong H^*(U;L_U)\oplus\C^{\chi(X)}[-n]$ (if $\rho\notin M$, $L_U$ is a non-trivial local system on $(\C^*)^n$ and $H^*(U;L_U)=0$).
\end{thm}
\begin{hyptab}
(V) & Hypercohomology equals the cohomology of the complex of global sections & Prerequisite of the weight formula (Theorem \ref{thm:weight}). For $L=L_0^m$ ($m\gg0$) follows from Theorem \ref{thm:serre}\\\midrule
(A) $L(D)$ ample & The polytope $\{\chi:c_i(\chi)\ge-1\}$ is simple, with normal fan $\Sigma$ and exactly $\chi(X)$ vertices & $s=n$: $\chi(X)$ vertex weights, each $t^n$\\\midrule
(GEN) & $A(\chi+\rho)\in W(\chi)$ degenerates over $K$ to $x\parallel m$, excluding arithmetic coincidences & Kills the contributions of intermediate strata $0<s<n$\\\midrule
(R) $r_i\ne1$ & For $r_i=1$ solutions with $x=0$ or $x\parallel m_0$ appear (resonance) & Kills resonance contributions in $s=n-1$ and $\chi=-\rho$\\\midrule
$\mathrm{Pf}(A)\ne0$ & $\sigma$ is log-symplectic & The log de Rham differential $(\chi+\rho)\wedge$ corresponds to $v_\chi\wedge$
\end{hyptab}
\begin{proof}
Use Theorem \ref{thm:weight}. Put $x=\chi+\rho$, $C=C(\chi)$, $W=W(\chi)$, $s=\dim W$. For $i\in C$, $\langle x,u_i\rangle=r_i-1$. The condition $Ax\in W$ is equivalent to $\langle m,Ax\rangle=\sum_{i<j}A_{ij}(m_ix_j-m_jx_i)=0$ for all $m\in\mathrm{Ann}_M(W)$.
\emph{$s=n$:} the condition is empty. By (A), the $\chi$ with $c_i\ge-1$ for which $C(\chi)$ spans the whole space are in bijection with vertices of the polytope (simple; lattice points since $X$ is smooth), one for each maximal cone, $\chi(X)$ in total; each gives $t^n$.
\emph{$1\le s\le n-2$:} $\mathrm{Ann}_M(W)$ has rank $\ge2$, so there are linearly independent integer vectors $m,m'$. The coefficients $m_ix_j-m_jx_i\in K$; by (GEN), $\sum A_{ij}(\cdots)=0$ forces all coefficients to vanish, i.e.\ $x\parallel m$. Likewise $x\parallel m'$, so $x=0$. For $C\ne\emptyset$, $x=0$ means $r_i-1=0$, contradicting (R). No solution.
\emph{$s=n-1$:} $\mathrm{Ann}_M(W)=\Z m_0$. Similarly $x=\tau m_0$, and for $i\in C$, $0=\tau\langle m_0,u_i\rangle=r_i-1$, contradicting (R). No solution.
\emph{$s=0$:} $W=0$ and the condition is $x=0$, i.e.\ $\chi=-\rho$ (requiring $\rho\in M$). Since $c_i(\chi)=-r_i\ge-1$ and $C=\emptyset$, we need $r_i\le0$. The contribution is $(1+t)^n$.
Add up. For the last statement: for $\rho\in M$, $L_U$ is trivial, and under (R) ``all $r_i\le0$'' is equivalent to ``no positive integral residue''. \qedhere
\end{proof}
\tV In 246 examples with all $r_i\ne1$ ($\mathbb{P}^2$ ($m=1,2,3$), $\mathbb{P}^1\times\mathbb{P}^1$ ($\mathcal O(1,1),\mathcal O(1,2),\mathcal O(2,3)$), $\mathbb{P}^4$ ($m=1$), $\mathbb{P}^2\times\mathbb{P}^2$ ($(1,1)$), $(\mathbb{P}^1)^4$ ($(1,1,1,1)$), lattice and rational points $\rho$, random integer $A$ modulo the prime $2^{61}-1$) the formula of the theorem agrees (0 discrepancies, \texttt{roundD\_check.py}, re-run and confirmed). In the wider set of 3626 examples including resonances the predictions of the weight formula also agree (\texttt{toric\_generic\_check2.py}; not re-run \tQ).

\begin{obs}[Necessity of (GEN) and dependence on the arithmetic of $\sigma$ \tV]\label{obs:gen}
In dimension 4, $\HH^*$ is not determined by $(X,D,L,\text{residues})$ alone. For $\mathbb P^2\times\mathbb P^2$ ($\mathcal O(1,1)$, $\rho=0$) one gets $[1,4,12,16,16]$ when $\sigma$ is block diagonal (product of the results for two surfaces, K\"unneth) and $[1,4,6,4,10]$ for general $A$. The difference $(0,0,6,12,6)=6t^2(1+t)^2$ comes from six codimension-2 strata $\{\text{double point of one factor}\}\times(\C^*)^2$. Mechanism (via the weight formula): for a weight with $|C|=2$ where $u_i,u_j$ are symplectic with respect to $\sigma$, $x\in A^{-1}W$ is uniquely determined by $\langle x,u_i\rangle=r_i-1$, and $t^2(1+t)^2$ contributes only if the arithmetic condition $x\in\rho+M$ holds. For the block structure this condition holds automatically. \tK $\HH$ jumps upward on countably many proper analytic subsets of the parameter space of log-symplectic structures (upper semicontinuity), and the general value is the minimum.
\end{obs}

%--------------------------------------------------------------
\subsection{Resonance $r_i=1$}
%--------------------------------------------------------------
If some component has $r_i=1$, Theorem \ref{thm:main2} does not apply ($d_Q=-u(1-r)p$ vanishes). The following is based on computation.
\begin{obs}[Single and double resonance ($\mathbb{CP}^2$, $\mathbb{CP}^1\times\mathbb{CP}^1$) \tV]\label{obs:res}
Let $R$ be the set of branches with $r_i=1$, $N_R$ the normal bundles of the resonant branches tensored with $L$, $K_R=(K_X^{-1}\otimes L)|_{D_R}$, and $\nu$ the number of double points not containing a resonant branch. Then
\[P_{\HH}(t)=t\,h^0(N_R)+t^2\bigl(h^1(N_R)+h^0(K_R)+\nu\bigr)+t^3\,h^1(K_R).\]
There are 0 discrepancies for single and double resonances in $\mathbb{CP}^2$ and for $\mathbb{CP}^1\times\mathbb{CP}^1$ (108 examples with non-integral residues). The exact definition of $N_R$ when there are several resonant branches is copied from the earlier records and not checked \tQ. When all residues are integers ($L_U$ trivial) and $m\ll0$, and when all branches are resonant ($L=K_X$), connecting maps of $H_c$-type act and the formula breaks (e.g.\ double resonance in $\mathbb{CP}^2$, $m=-5$: observed $[0,0,5,2]$, predicted $[0,0,6,3]$). $L=K_X$ is treated separately in Proposition \ref{prop:KX}. \tK The resonance formula is probably the general formula when $L_U$ is non-trivial and not all branches are resonant, but its proof requires a computation of connecting maps (an extension of Lemma \ref{lem:conn}) and is unfinished. For toric varieties Corollary \ref{cor:weight}(c) corresponds (proved \tP).
\end{obs}

%--------------------------------------------------------------
\subsection{Local terms at non-nodal singularities}\label{subsec:local}
%--------------------------------------------------------------
If $D$ has cusps, tacnodes, triple points etc., Theorem \ref{thm:main2} does not apply. The local model is $\sigma=h\,\partial_z\wedge\partial_w$, the module $\mathcal O\,h^s$ ($V=s\chi_h$, $\chi_h=h_w\partial_z-h_z\partial_w$), and the stalk complex $\mathcal C_s=(\mathcal O\xrightarrow{d_0}\mathcal O^2\xrightarrow{d_1}\mathcal O)$:
\[d_0f=(hf_w+sfh_w)\partial_z-(hf_z+sfh_z)\partial_w,\qquad d_1(p,q)=h(p_z+q_w)-(1-s)(h_zp+h_wq)\]
(formula \eqref{eq:dE} with $a=sh_w$, $b=-sh_z$). Globally $s=r$; for $\mathbb{CP}^2$, $s=-m/3$, and for $\mathbb{CP}^1\times\mathbb{CP}^1$, $-m/2$.

Let $h\in\C\{z,w\}$ be a quasi-homogeneous polynomial for positive integer weights $(w_z,w_w)$ and degree $N_h$, whose origin is an isolated singularity (hence $h$ is reduced) or a smooth point. Put $\delta=N_h-w_z-w_w$, $e=w_zz\partial_z+w_ww\partial_w$ ($eh=N_hh$), $J=(h_z,h_w)$, $T=\mathcal O/J$, $\mu=\dim T$, and let $\mathrm{Sp}$ be the multiset $\{\alpha_m=(\mathrm{wt}(m)+w_z+w_w)/N_h\}$ ($m$ running over a monomial basis of $T$).

\begin{thm}[Quasi-homogeneous local formula]\tP \tQ\label{thm:local}
For $s\ne1$ the Poincar\'e polynomial of the cohomology of $\mathcal C_s$ is
\[\mu\,t^2\;+\;[\,s\in\Z,\ s\le0\,](1+t)\;+\;(t+t^2)\cdot\#\{\alpha\in\mathrm{Sp}:\ s+\alpha\in\Z,\ s+\alpha\le1\}.\]
\end{thm}
\begin{hyptab}
Quasi-homogeneous & $e\,h=N_hh$ makes the complex split by weight, each weight finite-dimensional & Steps (2), (3)\\\midrule
Isolated singularity & $h_z,h_w$ is a regular sequence (Hilbert series formula), $h$ reduced (Saito's basis $e,\chi_h$) & Steps (1), (5), (4)\\\midrule
$s\ne1$ & For $s=1$ the differential of the quotient $Q$ is zero and the stalk is infinite-dimensional & Step (1)\\\midrule
Positive weights & Each weight component $\mathcal O_k$ is a finite-dimensional space of polynomials, and weightwise division is bounded & Step (2)
\end{hyptab}
\begin{proof}
\emph{(1) Decomposition.} $D=\{h=0\}$ is a plane curve, $e,\chi_h\in\Theta(-\log D)$ ($eh=N_hh$, $\chi_hh=0$) and $e\wedge\chi_h=-N_hh\,\partial_z\wedge\partial_w$, so since $h$ is reduced Saito's criterion (\cite{Saito} Thm.~1.8 \tC \tQ) gives $\Theta(-\log D)=\mathcal Oe\oplus\mathcal O\chi_h$. In the basis $(e,\chi_h)$ (the weight of $\chi_h$ is $\delta$, $[e,\chi_h]=\delta\chi_h$)
\begin{align}
d_0f&=-\tfrac1{N_h}\chi_h(f)\,e+\tfrac1{N_h}\bigl(ef+sN_hf\bigr)\chi_h,\label{eq:d0}\\
d_1(a\,e+b\,\chi_h)&=\bigl[(e+sN_h-\delta)a+\chi_h(b)\bigr]\,h\,\partial_z\wedge\partial_w.\label{eq:d1}
\end{align}
(\tV We checked these against the defining formulas with random polynomials $f,a,b$ and rational $s$ for 7 quasi-homogeneous $h$: 210 checks, 0 discrepancies, \texttt{local\_basis\_check.py}, written for this paper.) Hence $B_s:=(\mathcal O\to\Theta(-\log D)\to h\mathcal O\,\partial_z\wedge\partial_w)$ is a subcomplex of $\mathcal C_s$. The quotient $Q=\mathcal C_s/B_s$ is, via $X\mapsto X(h)\bmod h$, $Q^1\cong J\mathcal O_D\subset\mathcal O_D$, $Q^2=\mathcal O_D$, $Q^0=0$, with differential $-(1-s)\cdot$(inclusion), injective for $s\ne1$. $H^2(Q)=\mathcal O_D/J\mathcal O_D=\mathcal O/J$ ($h\in J$: $N_hh=w_zzh_z+w_wwh_w$). In the long exact sequence of the triangle $B_s\to\mathcal C_s\to Q$, $H^3(B_s)=0$ and $H^0(Q)=H^1(Q)=0$, so $H^k(B_s)=H^k(\mathcal C_s)\ (k\le1)$ and $0\to H^2(B_s)\to H^2(\mathcal C_s)\to\mathcal O/J\to0$. That is, the stalk is $\mu t^2+P(B_s)$.

\emph{(2) Weight decomposition.} $\mathcal O=\prod_k\mathcal O_k$ ($\mathcal O_k$ the quasi-homogeneous polynomials of weight $k$). The weight-$k$ part of $B_s$ is, with $a\in\mathcal O_{k+\delta}$, $b\in\mathcal O_k$, $B^0_k=\mathcal O_k$, $B^1_k=\{a\,e+b\,\chi_h\}$, $B^2_k=\mathcal O_{k+\delta}$; with $\kappa:=k+sN_h$, formulas \eqref{eq:d0}, \eqref{eq:d1} give
\[d_0f=\bigl(-\tfrac1{N_h}\chi_hf,\ \tfrac\kappa{N_h}f\bigr),\qquad d_1(a,b)=\kappa a+\chi_hb.\]
For $\kappa\ne0$ the homotopies $H^2(c)=(c/\kappa,0)$, $H^1(a,b)=N_hb/\kappa$ satisfy $d_1H^2=\mathrm{id}$, $H^1d_0=\mathrm{id}$, $H^2d_1+d_0H^1=\mathrm{id}$ (direct computation). Over all weights with $\kappa\ne0$ we have $|\kappa|\ge c>0$, so $H$ preserves convergent series, and $B_s$ splits into the finite-dimensional subcomplex of weight $k_0:=-sN_h$ (zero if $k_0\notin\Z$) and a contractible complex. The comparison problem between formal and convergent series does not arise thanks to this bounded homotopy.

\emph{(3) The weight-$k_0$ part.} Here $\kappa=0$, so $d_0f=(-\tfrac1{N_h}\chi_hf,0)$ and $d_1(a,b)=\chi_hb$. Put $K=\ker(\chi_h|\mathcal O_{k_0})$ and $\mathrm{Cok}=\mathrm{coker}(\chi_h:\mathcal O_{k_0}\to\mathcal O_{k_0+\delta})$. Then $H^0=K$, $H^1=K\oplus\mathrm{Cok}$, $H^2=\mathrm{Cok}$ (with $\mathcal O_j=0$ for $j<0$).

\emph{(4) Determination of $K$.} Let $f\in\mathcal O_{k_0}$ with $\chi_hf=0$. Since $\chi_hf=(df\wedge dh)/(dz\wedge dw)$ up to sign, $df\wedge dh=0$. Contracting with $e$, $k_0f\,dh-N_hh\,df=0$, i.e.\ $d(f^{N_h}/h^{k_0})=0$ ($k_0>0$). So $f^{N_h}=c\,h^{k_0}$. Since $h$ is reduced, for an irreducible factor $p\mid h$ we have $N_hv_p(f)=k_0$. If $f\ne0$ then $N_h\mid k_0$, and $f/h^{k_0/N_h}$ is a rational function whose $N_h$-th power is constant, i.e.\ a constant. Hence $K=\C\,h^{k_0/N_h}$ ($k_0\in N_h\Z_{\ge0}$) and $K=0$ otherwise (for $k_0=0$, $f$ is constant).

\emph{(5) $\dim\mathrm{Cok}$.} Let $d=k_0+\delta$, $a=w_z$, $b=w_w$. $\sum_k\dim\mathcal O_kx^k=1/((1-x^a)(1-x^b))$, and since $h_z,h_w$ is a regular sequence of degrees $N_h-a$, $N_h-b$, $\sum\dim T_kx^k=(1-x^{N_h-a})(1-x^{N_h-b})/((1-x^a)(1-x^b))$. From the identity $x^\delta+x^{N_h}-x^{N_h-a}-x^{N_h-b}=x^\delta(1-x^a)(1-x^b)$,
\[\frac{(1-x^{N_h-a})(1-x^{N_h-b})}{(1-x^a)(1-x^b)(1-x^{N_h})}-\frac{1-x^\delta}{(1-x^a)(1-x^b)}=\frac{x^\delta}{1-x^{N_h}}.\]
Taking the coefficient of $x^d$ gives $\sum_{r\ge0}\dim T_{d-rN_h}-(\dim\mathcal O_d-\dim\mathcal O_{k_0})=[\,k_0\in N_h\Z_{\ge0}\,]=\dim K$. Since $\dim\mathrm{Cok}=\dim\mathcal O_d-\dim\mathcal O_{k_0}+\dim K$, we get $\dim\mathrm{Cok}=\sum_{r\ge0}\dim T_{d-rN_h}$.

\emph{(6) Translation into the spectrum.} $s=-k_0/N_h$. For a basis monomial $m$ of $T_{d-rN_h}$ ($\mathrm{wt}(m)=d-rN_h$), $\alpha_m=(\mathrm{wt}(m)+w_z+w_w)/N_h=1-r-s$, i.e.\ $s+\alpha_m=1-r\in\Z_{\le1}$. Conversely each $\alpha$ with $s+\alpha\in\Z_{\le1}$ corresponds to exactly one $r\ge0$. Hence $\dim\mathrm{Cok}=\#\{\alpha\in\mathrm{Sp}:s+\alpha\in\Z,\ s+\alpha\le1\}$ and $\dim K=[\,s\in\Z_{\le0}\,]$. Steps (3) and (1) give the theorem. \qedhere
\end{proof}
\begin{rem}[Status and remaining unverified points]
\tV The dimensions of steps (3)--(6) ($K$, $\mathrm{Cok}$ and the formula above) were checked by exact rational computation for 14 types ($A_1,A_2,A_3,A_6,D_4,D_5,E_6,E_8$, $z^3+w^6$, $z^4+w^5$, $z^5+w^5$, $z^4+w^4$, $z^3+w^7$, $z^5+w^6$) and about 700 values of $k_0$, with 0 discrepancies (\texttt{roundH\_proof\_check.py}, re-run and confirmed). \tQ Remaining unverified: (a) that the local model $\mathcal C_s$ ($V=s\chi_h$) appears as the stalk of a global Poisson module (including germs that are not quasi-homogeneous; the sign convention of the residue $s$ was fixed to agree with $V=-\sigma^\sh\theta$ in formula \eqref{eq:dE}), (b) the number of Saito's criterion. The range of applicability is only quasi-homogeneous isolated singularities (or smooth points) and $s\ne1$; non-quasi-homogeneous isolated singularities, non-isolated singularities and $s=1$ are \tK (unproved).
\end{rem}
\begin{rem}[Interpretation \tC \tK]
From the Wang sequence of the Milnor fibration $h:B\setminus D\to\C^*$ (the fibre is a bouquet of $\mu$ circles, and the monodromy eigenvalues on $H^1$ are $e^{-2\pi\sqrt{-1}\alpha}$) one gets $(Rj_*L_s)_p=[s\in\Z](1+t)+(t+t^2)\#\{\alpha:s+\alpha\in\Z\}$. The formula of the theorem is this with the truncation $s+\alpha\le1$ (coming from the subcomplex $B_s$ computing the part of weight $\ge-\delta$). The rule $[s<\mathrm{lct}]\cdot(\cdots)$ of the earlier version (Part II v1) is wrong for $\mathrm{lct}<1/2$ (example: $z^4-w^5$, $\mathrm{lct}=9/20$, where for $s=11/20=1-\mathrm{lct}$ we have $s>\mathrm{lct}$ but the class exists) and agrees for $\mathrm{lct}\ge1/2$. Calder\'on-Moreno--Narv\'aez-Macarro \cite{CMNM} give an analogous logarithmic comparison theorem for free divisors (pages \tQ).
\end{rem}
\begin{obs}[Stalk dimensions for representative types \tP]\label{obs:localtable}
From Theorem \ref{thm:local}, $(\dim H^0,\dim H^1,\dim H^2)$ is as follows ($\mu$ is the Milnor number; the column ``$\alpha$'' is for the case that exactly one $\alpha\in\mathrm{Sp}$ with $s+\alpha\in\Z_{\le1}$ exists and $\alpha\notin\Z$):
\begin{center}\small
\begin{tabular}{llcccc}\toprule
Singularity & $\mathrm{Sp}$ & $\mu$ & generic $s$ & $s\in\Z_{\le0}$ & $\alpha$\\\midrule
smooth $(z)$ & $\emptyset$ & 0 & $(0,0,0)$ & $(1,1,0)$ & ---\\
double point $(zw)$ & $\{1\}$ & 1 & $(0,0,1)$ & $(1,2,2)$ & ---\\
cusp $(z^2-w^3)$ & $\{\frac56,\frac76\}$ & 2 & $(0,0,2)$ & $(1,1,2)$ & $(0,1,3)$\\
tacnode $(z^2-w^4)$ & $\{\frac34,1,\frac54\}$ & 3 & $(0,0,3)$ & $(1,2,4)$ & $(0,1,4)$\\
$A_4$ $(z^2-w^5)$ & $\{\frac7{10},\frac9{10},\frac{11}{10},\frac{13}{10}\}$ & 4 & $(0,0,4)$ & $(1,1,4)$ & $(0,1,5)$\\
$D_4$ $(zw(z-w))$ & $\{\frac23,1,1,\frac43\}$ & 4 & $(0,0,4)$ & $(1,3,6)$ & $(0,1,5)$\\
$D_5$ $(w(z^2-w^3))$ & $\{\frac58,\frac78,1,\frac98,\frac{11}8\}$ & 5 & $(0,0,5)$ & $(1,2,6)$ & $(0,1,6)$\\\bottomrule
\end{tabular}
\end{center}
Of the values in the table of the earlier version (Part I \S I.6.6), the last column for $A_4$ had been left blank.
\end{obs}
\noindent \tP Correction: for $D_5=z^2w-w^4$ the weights are $(3,2)$, degree 8, $\mathrm{Sp}=\{5/8,7/8,1,9/8,11/8\}$, $\mathrm{lct}=5/8$, $\mu=5$. The weights $(3,1,7)$ in an internal memo were wrong, and values derived from them are invalid (the table and text of the earlier version were written with the correct weights).

\begin{obs}[Triple point with unequal residues \tV]\label{obs:triple}
For $h=zw(z-w)$ ($D_4$) and the module $\prod h_i^{s_i}$ (a different $s_i\ne1$ on each branch), the stalk is $\C^4[-2]\oplus R$, with $R=H^*(U_{\rm loc};L_s)$ ($s_1+s_2+s_3<2$) and $R=0$ (otherwise). $U_{\rm loc}=\C^2\setminus\{3\text{ lines}\}\simeq\C^*\times(\mathbb P^1\setminus3\text{ points})$. Concretely, the stalk is $(1,3,6)$ if all $s_i\in\Z_{\le0}$, $(0,1,5)$ if $\sum s_i\in\Z_{\le1}$ and not all $s_i$ are integers, and $(0,0,4)$ otherwise. There are 0 discrepancies for 533 random triples $s_i\in\tfrac16\Z\cap[-2,2]\setminus\{1\}$ (\texttt{local\_unequal.py, local\_unequal\_rule.py}; not re-run \tQ). At a node the per-branch threshold $s_i<1$ decides the presence of the $Rj_*$ part, whereas at a triple point it is decided by whether the \emph{sum} of the residues exceeds the threshold $2=w_z+w_w$. For equal residues $s_1=s_2=s_3=s$ it agrees with the values of Theorem \ref{thm:local} for $D_4$ at $s\in\tfrac13\Z$, $-10\le3s\le10$ ($s\ne1$) (checked \tV). \tK General principle: the threshold is determined by comparing the weight of the module under the symmetry torus of the germ with the log canonical weight.
\end{obs}
\begin{conj}[Refined Riemann--Roch (equal residues)]\tK\label{conj:refined}
Let $X$ be a compact surface, the zero curve $D$ of $\sigma\in H^0(-K_X)$ reduced and all its singularities isolated, quasi-homogeneous with $\mathrm{lct}_p\ge1/2$ (types $A_k,D_k,E_{6,7,8}$ etc.). Let $L$ be a Poisson module with $c_1(L)=sK_X$, and $\mathcal F_s$ the constructible complex equal to $L_s$ on $U$ and at $p\in D$ to the stalk of $Rj_*L_s$ if $s<\mathrm{lct}_p$ and $0$ if $s>\mathrm{lct}_p$. (a) For $s\notin\mathrm{Sp}(D)\cup\{1\}$, $\HH^*=H^*(X;\mathcal F_s)\oplus\bigoplus_p\C^{\mu_p}[-2]$. (b) The Euler characteristic is always $\chi(X)$ (from $\chi(D)=\sum\mu_p$ and $\delta$-invariants \tP). (c) For $s\ll0$, $\mathcal F_s=Rj_*L_s$ and $\HH^*=H^*(U;L_s)\oplus\bigoplus\C^{\mu_p}[-2]$ depends only on $s\bmod\Z$ and vanishes in degrees $\ge3$. (d) For $s\in\mathrm{Sp}(D)$, rank-one cancellations may occur.
\end{conj}
\tV For $\mathbb{CP}^2$ (smooth, nodal, cuspidal, conic plus secant, conic plus tangent, triangle, triple point) and $\mathbb{CP}^1\times\mathbb{CP}^1$ (4 examples) there is no contradiction except in (d), for $m=-15..15$ ($-9..8$). Triple point $xy(x-y)$ ($\mathbb{CP}^2$): $m=1,2,4,5,\dots$ and $m=-1$ give $[0,1,4,0,0]$, $m=-5,-7,-8$ give $[0,0,4,1,0]$, and $m=-2,-3,-4$ give $[0,0,3,0,0]$. We conjecture that the latter anomaly (at $m=-2,-4$ one class of $\C^4$ in degree 2 cancels against one class of $H_c$ in degree 3) is due to higher differentials of the spectral sequence, but there is no proof and no explanation by Theorem \ref{thm:local} \tK. Not proved: the global identification is a fit, and the stalk result concerns only quasi-homogeneous germs.
\begin{obs}[Non-reduced boundary \tV]\label{obs:nr}
When $\sigma=\Pi_F$ with $F$ having multiple factors, the dimension of $\HH^*$ grows linearly in $m$ (exact linear algebra, $m=-15,\dots,15$, \texttt{roundF\_nonreduced.py}; not re-run \tQ). (a) $F=x^2y$ (double line plus a line, one node, $U\simeq S^1$): for $m\ge-2$, $P(t)=[3\mid m](1+t)+3t^2+(m+2)(t+t^2)$; $m=-3$: $3t^2$; $m\le-4$: image under Serre duality $m\leftrightarrow-m-6$. (b) $F=x^3$ (triple line, $U=\C^2$): $[1,2m+5,2m+7]$ for $m\ge0$, $[0,2,5]$ for $m=-1$, $[0,0,3]$ for $m=-2,-3,-4$. (c) $\mathbb P^1\times\mathbb P^1$, $D=2C$ ($C$ a smooth $(1,1)$ curve), $L=\mathcal O(m,m)$: $H^0=1$, $H^1=2m+3$, $H^2=2m+6$ for $m\ge0$.
\end{obs}
\begin{conj}[Non-reduced boundary]\tK\label{conj:NR}
For $X=\mathbb P^2$ or $\mathbb P^1\times\mathbb P^1$, $D=\sum k_iD_i$ ($D_i$ smooth, $D_{\mathrm{red}}$ SNC), and $m\ge0$ (or in the range where higher sheaf cohomology vanishes),
\[\HH^*(L)=H^*(U;L_U)\oplus\C^{\chi(D_{\mathrm{red}})}[-2]\oplus\bigoplus_i\bigoplus_{j=1}^{k_i-1}\bigl(H^0(D_i,L|_{D_i}\otimes N_{D_i}^{j})[-1]\oplus H^0(D_i,L|_{D_i}\otimes N_{D_i}^{j})[-2]\bigr).\]
\end{conj}
\noindent It is confirmed only for the three families of Observation \ref{obs:nr}. The total Euler characteristic is $\chi(U)+\chi(D_{\mathrm{red}})=\chi(X)$, consistent with Theorem \ref{thm:E}. Heuristics: the Poisson cohomology of $z^k\partial_z\wedge\partial_w$ has an infinite-dimensional part (results of Monnier, Roger--Vanhaecke type \tC \tQ, unchecked). It fails for small $m$ ($m=-1$ for $x^3$ does not match), so the exact boundary of its range of validity is unknown.

%==============================================================
\section{General dimension: the reduction theorem and Theorem V}\label{sec:reduction}
%==============================================================
\subsection{Setting, the quotient complex, the condition (GEN$_Q$)}
Let $X$ be a compact complex manifold ($\dim X=n$ even), $D=\sum D_i$ an SNC divisor, $\sigma$ log-symplectic, and $L$ a Poisson module (a flat log connection along $D$ with residues $r_i$) (Definition \ref{defn:logsymp}, Theorem \ref{thm:main1}). We use the following sheaves of complexes.
\[C=(\wedge^\bullet T\otimes L,\ d_L=[\sigma,\cdot]+V\wedge),\quad B=\wedge^\bullet T_D\otimes L\subset C,\quad Q=C/B.\]
$\HH^k=H^k(X;C)$. By $\wedge\sigma^\sh$, $B\cong\mathrm{DR}_{\log}(L)$ (Lemma \ref{lem:iso}) \tC. There is a triangle $B\to C\to Q\xrightarrow{\ \partial\ }B[1]$.

\paragraph{Structure of $Q$.} In logarithmic coordinates $(z_1,\dots,z_k,w)$ an $\mathcal O$-basis of $\wedge T$ is $\partial_S\wedge\partial_w^{Y}$ ($S\subset\{1..k\}$) and one of $\wedge T_D$ is $E_S\wedge\partial_w^Y=z^S\partial_S\wedge\partial_w^Y$. Hence $Q$ is spanned by the monomials $z^a\partial_S\wedge\partial_w^Y$ for which the set $T(s)$ of $i\in S$ with $a_i=0$ is non-empty; these are read as ``multivector fields with poles along each branch in $T(s)$'', $z^{a-S}E_S\wedge\cdots$. Since $\sigma\in\wedge^2T_D$ and $V\in T_D$, $d_L$ creates no new poles ($[X,z_i^{-1}]\in z_i^{-1}\mathcal O$ for $X\in T_D$). Hence $F_kQ=\{|T(s)|\le k\}$ is an increasing sequence of subcomplexes of $Q$, and $\mathrm{gr}_0Q=0$. In the top piece, $|T(s)|=n$ corresponds only to the monomial $\partial_1\wedge\cdots\wedge\partial_n$ at an $n$-fold point $p$. In degree $n$, $Q^n=(\wedge^nT/\wedge^nT_D)\otimes L=K_X^{-1}\otimes L|_D$ ($\wedge^nT=K_X^{-1}$, $\wedge^nT_D=K_X^{-1}(-D)$). This observation defines a map from $Q$ to points:
\[\varepsilon^n:\ Q^n=K_X^{-1}\otimes L|_D\longrightarrow\bigoplus_{p\in D_{(n)}}(K_X^{-1}\otimes L)_p\otimes\C,\quad(\text{value at the }n\text{-fold points}),\qquad \varepsilon^k=0\ (k\ne n).\]

\begin{prop}[$\mathrm{gr}\,Q$ and $\varepsilon$ \tQ]\label{prop:grQ}
(a) $\varepsilon:Q\to\bigoplus_{p\in D_{(n)}}\C_p[-n]$ is a morphism of complexes (i.e.\ $d_L(Q^{n-1})$ takes value 0 at each $n$-fold point).
(b) $\mathrm{gr}_kQ=\bigoplus_{|T|=k}Q_T$, $Q_T\cong\bigl(\wedge^\bullet T_{D_T}(-\log D'_T)\otimes L|_{D_T}\bigr)\cdot z^{-T}E_T$ ($D'_T$ the divisor cut out on $D_T$ by the other components), with differential $d_T=[\sigma_T,\cdot]+\nu_T\wedge$. Here $\sigma_T$ is the Poisson structure induced on $D_T$, $\nu_T=\sum_{i\in T}(r_i-1)H_i$, $H_i=\sigma^\sh(\dlog z_i)$.
\end{prop}
\begin{proof}[Status and basis]
For $n=2$ this is Lemma \ref{lem:Q} (explicit formula for $d_Q$): at a smooth point $d_Q=-u(1-r)\cdot$ and at a double point the image is $\mathfrak m_D$, so (a) holds, and (b) follows from the definitions \tP \tV. For general $n$, the ``$-1$'' in (b) comes from $[\sigma,z_i^{-1}]=-z_i^{-1}H_i$, the ``$r_i$'' from the module structure $V=\sigma^\sh(\sum r_i\dlog z_i)+\cdots$, and $[\sigma_T,\nu_T]=0$ from the Jacobi identity for $\sigma$ and the fact that $\nu_T$ is a sum of Poisson vector fields. We have not written out the Schouten bracket computation line by line, and there is no computation checking (a), (b) for general $n$. This paper uses the proposition only to rephrase the \emph{hypothesis} (GEN$_Q$) of Theorem \ref{thm:R} in local terms, not in the proof of Theorem \ref{thm:R}. \tQ
\end{proof}

\begin{defn}[Condition (GEN$_Q$); distinct from (GEN) of Theorem \ref{thm:T}]\label{defn:GEN}
$\varepsilon:Q\to\bigoplus_{p\in D_{(n)}}\C_p[-n]$ is a morphism of complexes and a quasi-isomorphism.
\end{defn}
\begin{defn}[Condition (GEN$'$)]
For every stratum $T$ with $1\le|T|\le n-1$, $Q_T$ is acyclic stalkwise. If Proposition \ref{prop:grQ} holds, (GEN$'$) $\Rightarrow$ (GEN$_Q$) (all graded pieces of $F_{n-1}Q$ are acyclic so $F_{n-1}Q$ is acyclic, and $Q/F_{n-1}Q=\mathrm{gr}_nQ\cong\bigoplus\C_p[-n]$) \tQ. For $n=2$, (GEN$_Q$) $\iff$ all $r_i\ne1$ (Lemma \ref{lem:Q}) \tP.
\end{defn}

\begin{prop}[Necessary conditions for (GEN$'$) and examples]\label{prop:genNec}
(a) If $\nu_T\equiv0$ then $Q_T$ is the Lichnerowicz complex of $\sigma_T$ and locally constant Casimirs enter $H^0$. In particular $r_i\ne1$ is necessary for $|T|=1$. \tP \tQ (depends on Proposition \ref{prop:grQ}(b))
(b) For $|T|=1$ at a smooth point of $D$, if in logarithmic Darboux coordinates $\sigma=E_z\wedge\partial_{w_1}+\sum_j\partial_{w_{2j}}\wedge\partial_{w_{2j+1}}$ then $\nu=(r-1)\partial_{w_1}$, and $Q$ is the mapping cone of $(r-1)\cdot\mathrm{id}$ into the Lichnerowicz complex of the symplectic part, which is acyclic $\iff r\ne1$. \tC \tQ Logarithmic Darboux is known for real $b$-symplectic structures by Guillemin--Miranda--Pires, Adv.\ Math.\ 264 (2014); the holomorphic version presumably follows from the same Moser argument (the logarithmic version of Proposition \ref{prop:moser}), but we have not checked it.
\end{prop}

\subsection{The reduction theorem}
\begin{thm}[Reduction theorem]\tP\label{thm:R}
In the above setting assume (i) no $r_i$ lies in $\Z_{\ge1}$, and (ii) (GEN$_Q$) (Definition \ref{defn:GEN}). Then
\[\HH^*(L)\ \cong\ H^*(U;L_U)\ \oplus\ \bigoplus_{p\in D_{(n)}}\C[-n].\]
In particular $\chi(X)=\chi(U)+\#D_{(n)}$ ($n$ even).
\end{thm}
\begin{hyptab}
(i) $r_i\notin\Z_{\ge1}$ & Avoids resonance. Condition for $\mathrm{DR}_{\log}(L)\simeq Rj_*L_U$ (Deligne) & Step (1)\\\midrule
(ii) (GEN$_Q$) & $Q$ contributes only at $n$-fold points. If it fails, extra classes appear (Observation \ref{obs:gen}) & Step (2)\\\midrule
$D$ SNC, $\sigma$ log-symplectic & $B\cong\mathrm{DR}_{\log}(L)$ & Step (1)\\\midrule
$X$ compact & $H^*(X;\cdot)$ is finite-dimensional and Theorem \ref{thm:E} can be used & Step (4)\\\midrule
$n$ even & Sign $(-1)^n=1$. Automatic for log-symplectic & Step (4)\\\midrule
(not needed) projectivity, ampleness of $L$, vanishing (V) & --- & Used only in the weight computation of Theorem \ref{thm:T}
\end{hyptab}
\begin{proof}
(1) By (i), Deligne \cite{Del} (II.3.13; Esnault--Viehweg \cite{EV}) \tC gives $B\simeq\mathrm{DR}_{\log}(L)\simeq Rj_*L_U$.
(2) By (ii), $\varepsilon:Q\xrightarrow{\sim}\bigoplus_p\C_p[-n]$. Hence the connecting morphism of the triangle $B\to C\to Q\xrightarrow\partial B[1]$ is $\partial\in\mathrm{Hom}_{D(X)}(Q,Rj_*L_U[1])=\mathrm{Hom}(\bigoplus_pi_{p*}\C[-n],Rj_*L_U[1])$. By the adjunction $j^*\dashv Rj_*$ this equals $\mathrm{Hom}(j^*\bigoplus i_{p*}\C[-n],L_U[1])=0$ (since $p\in D$, $j^*i_{p*}=0$).
(3) So the triangle splits and $C\simeq B\oplus Q\simeq Rj_*L_U\oplus\bigoplus_p\C_p[-n]$. Take $H^*(X;\cdot)$ and use $H^*(X;Rj_*L_U)=H^*(U;L_U)$ (by definition).
(4) Euler characteristic: by Theorem \ref{thm:E}, $\chi(\HH)=(-1)^n\chi(X)=\chi(X)$, and the right-hand side is $\chi(U)+(-1)^n\#D_{(n)}=\chi(U)+\#D_{(n)}$ ($L_U$ has rank 1). \qedhere
\end{proof}
\begin{rem}[Difference from the earlier version \tP]
The proof of the reduction theorem in the earlier version used the filtration of $Q$ and the graded pieces (GEN$'$), and the derivation of the graded pieces for general $n$ was a sketch. The proof above assumes the condition directly in the form (GEN$_Q$) ($\varepsilon$ is a quasi-isomorphism), so it does not depend on Proposition \ref{prop:grQ}. The step (GEN$'$) $\Rightarrow$ (GEN$_Q$) remains \tQ.
\end{rem}
\tV For $n=2$ it agrees with Theorem \ref{thm:main2}, and in 61 toric examples with $n=2,4$ (\texttt{roundK\_check.py}) the predictions of the reduction theorem and the closed formula (Theorem \ref{thm:K}) agree with the measured values (0 discrepancies). In the toric case, (GEN$_Q$) holds because the weight-$\chi$ part of $Q_T$ is the Koszul complex of $\wedge(N_\C/W_T)$ with differential $v_\chi\wedge$, acyclic when $v_\chi\notin W_T$, i.e.\ by the same argument as the exclusion of $x=0$ in the proof of Theorem \ref{thm:T} under the hypotheses (GEN) of Theorem \ref{thm:T} (independence of $A_{ij}$ over $\mathbb Q(\rho)$) and (R) \tP \tQ (depends on Proposition \ref{prop:grQ}). Theorem \ref{thm:T} itself is proved directly from the weight formula and does not depend on this argument.

\begin{conj}[When (GEN$'$) fails]\tK\label{conj:nonacyc}
$E_1^T=H^*(D_T;Q_T)\Rightarrow H^*(X;Q)$. If all $Q_T$ with $\nu_T=0$ are of ``product type'', the extra part is $\sum_Tt^{|T|}H^*(U_T;L_T)$. The $6t^2(1+t)^2$ of Observation \ref{obs:gen} can be read as follows: the six strata $T=(\text{double point of one factor})\times Y$ ($\nu_T=0$) each contribute $t^2(1+t)^2$ through $Q_T\simeq Rj_*\C_{U_Y}$ (\tV numerical agreement only).
\end{conj}

%--------------------------------------------------------------
\subsection{Theorem V: vanishing of $H^k(U;L_U)$}
%--------------------------------------------------------------
The reduction theorem leaves $H^*(U;L_U)$ as an unknown. If all residues are non-integers it is determined as follows.

\begin{thm}[Theorem V]\tP\label{thm:V}
Let $X$ be a compact complex manifold ($\dim_\C X=n\ge1$), $D=\sum_{i=1}^sD_i$ a reduced SNC divisor, and $U=X\setminus D$ \emph{affine}. Let $L_U$ be a rank one local system on $U$ whose simple monodromy around $D_i$ is $e^{-2\pi\sqrt{-1}\,r_i}$, with \emph{all $r_i\notin\Z$}. Then $H^k(U;L_U)=0$ for $k\ne n$, hence $\dim H^n(U;L_U)=(-1)^n\chi(U)$.
\end{thm}
\begin{hyptab}
$r_i\notin\Z$ (all $i$) & Stalks vanish on each stratum & (a). If all $r_i\in\Z$ on some stratum, the stalk is $H^*(T^k)\ne0$, which is the source of the jumps in \S\ref{sec:surface}, \S\ref{sec:toric}\\\midrule
$U$ affine & Artin vanishing & (c). Fails for non-affine $U$, e.g.\ $D=\emptyset$\\\midrule
SNC & Local fundamental group is $\Z^k$ & (a). For non-SNC (triple points etc.) the local fundamental group is not $\Z^k$ and $\sum r_i$ enters the stalk rule (Observation \ref{obs:triple})\\\midrule
$X$ compact & Global sections of $Rj_*$ give the cohomology of $U$, and if $j_!=Rj_*$ then $H^k=H^k_c$ & (b)
\end{hyptab}
\begin{proof}
(a) The stalk of $Rj_*L_U$ at $p\in D_I$ ($|I|=k\ge1$) is $H^*(U\cap B;L_U)$ for a small polydisc $B$. $U\cap B\simeq(\Delta^*)^k\times\Delta^{n-k}\simeq(S^1)^k$, and $L_U$ is the external tensor product of local systems with monodromy $\mu_i=e^{-2\pi\sqrt{-1}r_i}\ne1$ on each circle. By K\"unneth $H^*((S^1)^k;L)=\bigotimes_{i\in I}H^*(S^1;L_{\mu_i})$, and for $\mu\ne1$, $H^0=\ker(\mu-1)=0$ and $H^1=\mathrm{coker}(\mu-1)=0$. Hence the stalks vanish over $D$ and $Rj_*L_U=j_!L_U$.
(b) Since $X$ is compact, $H^k(U;L_U)=H^k(X;Rj_*L_U)=H^k(X;j_!L_U)=H^k_c(U;L_U)$.
(c) Since $U$ is affine, Artin's vanishing theorem (Andreotti--Frankel) \tC \cite{SGA4} gives $H^q(U;F)=0\ (q>n)$ for every local system $F$. Hence $H^k(U;L_U)=0$ for $k>n$. By Poincar\'e duality $H^k_c(U;L_U)\cong H^{2n-k}(U;L_U^\vee)^*$ on the oriented $2n$-manifold $U$ and (b), $H^k(U;L_U)=0$ when $2n-k>n$, i.e.\ $k<n$. Finally $\chi(U)=(-1)^n\dim H^n$. \qedhere
\end{proof}
\begin{cor}\tP\label{cor:Vgen}
If the hypotheses of Theorem \ref{thm:R} hold, all $r_i\notin\Z$, and $U$ is affine, then $P_{\HH}(t)=\chi(X)\,t^n$.
\end{cor}
\begin{proof}
By Theorems \ref{thm:R} and \ref{thm:V}, $P=\dim H^n(U;L_U)t^n+\#D_{(n)}t^n=(\chi(U)+\#D_{(n)})t^n=\chi(X)t^n$ ($n$ even). \qedhere
\end{proof}
For the ``mixed'' case where only some $r_i$ are integers: the stalk survives over $D_I$ only if all $r_i\ (i\in I)$ are integers, but vanishing for $k<n$ is unproved \tK.
\paragraph{Computational confirmation \tV.} Exact toric computations with $n=4$ (linear algebra modulo $2^{61}-1$): $\mathbb P^4$ ($\mathcal O(1)$, $r=(1/2,1/3,1/5,1/7,-457/210)$) gives $(0,0,0,0,5)$, $\mathbb P^4$ ($\mathcal O(2)$, $r=(1/3,1/3,1/3,1/3,-10/3)$) gives $(0,0,0,0,5)$, $\mathbb P^2\times\mathbb P^2$ gives $(0,0,0,0,9)$, and $(\mathbb P^1)^4$ gives $(0,0,0,0,16)$. All are $\chi(X)t^4$, in agreement with the prediction of Corollary \ref{cor:Vgen}. A mixed example ($\mathbb P^1\times\mathbb P^3$ with one integral residue $-2$) also gave $(0,0,0,0,8)$, but this is an observation, not a theorem. The scope of the verification is toric, constant $\sigma$ in the log frame and $L$ sufficiently ample; it is not a test of Theorem V for non-toric $U$. These $n=4$ computations were re-run while writing this paper with \texttt{roundJ2\_vanish\_check.py} ($\mathbb P^2\times\mathbb P^2$ with $\mathcal O(1,2)$ and all residues non-integral gives $(0,0,0,0,9)$; the mixed example $\mathcal O(1,1)$ with the integral residue $-1$ also gives $(0,0,0,0,9)$, which is also only an observation).

\begin{ex}[Products \tP \tV]\label{ex:prod}
Let $(X_1,D_1,L_1)$, $(X_2,D_2,L_2)$ be objects of Theorem \ref{thm:B} (or \ref{thm:main2}) and $\sigma=\sigma_1\boxplus\sigma_2$ (block diagonal). The complex is an external tensor product, so by K\"unneth $P_{\HH}=P_1\cdot P_2$ (as $\HH$ is finite-dimensional hypercohomology). For $\mathbb P^2\times\mathbb P^2$ ($\mathcal O(1,1)$, $\rho=0$) in dimension 4, when $P_1=P_2=[1,2,4]$ the product is $[1,4,12,16,16]$, which differs from the value $[1,4,6,4,10]$ of the weight formula for general $\sigma$ (Observation \ref{obs:gen}). A $\sigma$ with product structure is a special case violating (GEN$_Q$).
\end{ex}

%==============================================================
\section{Comparison with regular type: flat symplectic fibre bundles}\label{sec:regular}
%==============================================================
We treat the ``regular type'', in which the type is constant, in contrast to our earlier objects (type jumping), and show that the growth of dimensions is qualitatively different.

\subsection{Setting and independence of $\omega$}
\begin{defn}[Setting (S)]\label{defn:S}
Let $X$ be a compact complex manifold, $\tilde X\to X$ the universal cover, $\Gamma=\pi_1(X)$, and $(Y,\omega_Y)$ a compact connected symplectic manifold. For a homomorphism $\rho:\Gamma\to\mathrm{Symp}(Y,\omega_Y)$ put $M=(Y\times\tilde X)/\Gamma$. The generalized complex structure ``symplectic structure on $Y$ $\times$ complex structure on $X$'' is $\Gamma$-invariant and so descends to $M$, giving a regular-type generalized complex structure of type $k=\dim_\C X$. The leaves are the fibres $Y_x$ and the leaf space is $X$. Let $E=\pi^*G$ ($G$ a holomorphic line bundle on $X$) be an SGH (strong generalized holomorphic) line bundle. \tC Definitions of regular type and SGH: Pal--Poddar \cite{PP} Def.\ 3.1 etc.\ (numbers \tQ).
\end{defn}
This setting includes direct products $Y\times X$ ($\Gamma$ trivial), $S^2$-bundles (holonomy in $SO(3)$), Kodaira--Thurston type $T^2$-bundles, etc. $\HH(E)=H^*(M;C(L;E))$ is defined as in \S\ref{sec:prelim}.

\begin{prop}\tP\label{prop:omega}
Let $\tilde\omega$ be the closed 2-form on $M$ obtained by descending $\mathrm{pr}_Y^*\omega_Y$. Then $e^{i\tilde\omega}$ is a Courant automorphism and $L=e^{i\tilde\omega}L_0$ ($L_0=T_{F,\C}\oplus T^{0,1}_{\mathrm{hor}}\oplus N^*$, $N^*=\pi^*T^{*1,0}_X$). Hence $L\cong L_0$ as Lie algebroids, and the generalized Dolbeault complex is identified with
\[C(L;E)=\Omega^\bullet_{F_\C}(M;E\otimes\wedge^\bullet T^{1,0}_X),\qquad F_\C=T_{F,\C}\oplus T^{0,1}_{\mathrm{hor}}\ (\text{involutive}),\]
and $N$ acts trivially on SGH bundles. The cohomology does not see $\omega$ (it sees only the foliated complex structure).
\end{prop}
\begin{proof}
Locally this is the statement for products, and all objects are $\Gamma$-invariant. Since $d\tilde\omega=0$, $e^{i\tilde\omega}$ preserves the Courant bracket. That $N$ acts trivially is the definition of SGH (holomorphic in the transversal direction, flat along the leaves). \qedhere
\end{proof}

\subsection{Main theorem: Leray and Lefschetz}
\begin{thm}[$\HH^*$ for regular type]\tP \tQ\label{thm:L}
Assume setting (S) and (HL): $(Y,\omega_Y)$ satisfies hard Lefschetz (e.g.\ $Y$ K\"ahler). Let $V^j$ be the flat bundle with fibre $H^j(Y)$ (Gauss--Manin) and $E=\pi^*G$. Then
\[\HH^k(M,E)\ \cong\ \bigoplus_{a+b+j=k}H^a\bigl(X,\ \wedge^bT_X\otimes V^j\otimes G\bigr).\]
$X$ need not be K\"ahler or projective.
\end{thm}
\begin{hyptab}
(S) flat fibre bundle & The decomposition $L\cong L_0$ from Proposition \ref{prop:omega}, and Ehresmann-type local triviality (the fibre cohomology becomes the local systems $V^j$) & Step (i)\\\midrule
$Y$ compact connected & $H^j(Y)$ is finite-dimensional and the leafwise cohomology sheaf is the sheaf of smooth sections of $V^j$ & Step (i)\\\midrule
(HL) hard Lefschetz & $\Lambda=\omega_F\wedge$ gives isomorphisms $V^{m_Y-i}\to V^{m_Y+i}$ (of flat bundles) & Step (iii): Deligne's lemma, $d_r=0\ (r\ge2)$\\\midrule
$X$ compact & Each $H^a(X,\cdot)$ is finite-dimensional & Sum of dimensions\\\midrule
$E=\pi^*G$ (SGH) & $N$ acts trivially, so the $\wedge^bT_X$-degree $b$ is a direct summand (no Poisson-type $d_1$) & Step (i): form of $E_2$
\end{hyptab}
\begin{proof}
(i) (Double complex) By Proposition \ref{prop:omega}, $C(L;E)=\Omega^\bullet_{F_\C}(M;E\otimes\wedge^\bullet T_X)$. This forms a double complex of ``leafwise de Rham'' and ``Dolbeault on $X$ (coefficients $\wedge^bT_X\otimes G$)''. Since the fibres are compact, the leafwise cohomology sheaf is the sheaf of $C^\infty$ sections of $V^j$, and $\bar\partial$ on it is the $(0,1)$-part of the Gauss--Manin connection, independent of the choice of lifts (a difference of lifts is a leafwise Lie derivative, which is exact on closed forms). Hence
\[E_2^{a,j}=\bigoplus_bH^a(X,\wedge^bT_X\otimes V^j\otimes G).\]
(ii) The mixed components of $\omega_F:=\tilde\omega|_F$ vanish ($\tilde\omega$ is locally $\mathrm{pr}_Y^*\omega_Y$). Hence $\Lambda=\omega_F\wedge$ is a chain map of the double complex (commuting with $d_F$ and $\bar\partial$), acts on each term of the spectral sequence and commutes with $d_r$. On $E_2$ it is the fibrewise Lefschetz operator $V^j\to V^{j+2}$.
(iii) (Deligne's lemma) By hard Lefschetz $V^{m_Y-i}\to V^{m_Y+i}$ is an isomorphism. If $x$ is a primitive element of weight $s$, then $\Lambda^{s+1}d_2x=0$, but $d_2x$ lies on the side of fibre degree $m_Y-s-1$ where $\Lambda^{s+1}$ is injective. So $d_2=0$; inductively $d_r=0\ (r\ge2)$ and the dimensions add up. \qedhere
\src{Deligne \cite{Del68} (Publ.\ Math.\ IH\'ES 35). Leray spectral sequence: Bredon \cite{Bre} (location \tQ).}
\end{proof}
\tQ Points that remain a sketch: in (i) the identification of $\bar\partial$ with the $(0,1)$-part of Gauss--Manin and of $E_2$ of the leafwise Dolbeault-type double complex are not written in detail. Without (HL), $d_2\ne0$ may occur \tK (no concrete example computed).

\begin{cor}[Degree-wise Riemann--Roch]\tP\label{cor:L}
Under the hypotheses of Theorem \ref{thm:L}, assume further that $X$ is projective and $G$ ample. There is $m_0$ such that for $m\ge m_0$, $H^a(X,\wedge^bT_X\otimes V^j\otimes G^m)=0\ (a>0)$, hence
\begin{gather*}\HH^k(M,E^m)=\bigoplus_{b+j=k}H^0\bigl(X,\wedge^bT_X\otimes V^j\otimes G^m\bigr),\\ \dim\HH^k(M,E^m)=\sum_{b+j=k}b_j(Y)\int_X\ch(\wedge^bT_X)\,e^{m\,c_1(G)}\,\Td(X).\end{gather*}
The right-hand side is a polynomial in $m$ for each $k$.
\end{cor}
\begin{hyptab}
$G$ ample, $X$ projective & Serre vanishing is applied to the coherent sheaves $\wedge^bT_X\otimes V^j$ & Vanishing of all terms of degree $a>0$\\\midrule
$V^j$ flat & $\ch(V^j)=\mathrm{rk}\,V^j=b_j(Y)$ & $b_j(Y)$ comes out of the integral\\\midrule
Theorem \ref{thm:L} & Decomposition of $\HH^k$ into $\bigoplus H^a$ & First equality\\\midrule
HRR & When $a>0$ vanishes, $\dim H^0$ equals $\chi$ & Last equality
\end{hyptab}
\begin{proof}
Apply Serre vanishing \tC \cite{Hart} (III.5.2) to the locally free sheaves $\wedge^bT_X\otimes V^j$ and take the largest $m_0$ over the finitely many $(b,j)$. The rest is HRR and $\ch(V^j)=b_j(Y)$. \qedhere
\end{proof}
The total Euler characteristic is $\sum_k(-1)^k\dim\HH^k=\chi(Y)\cdot(-1)^{\dim X}\chi(X)$, consistent with (the extension to this setting of) Theorem \ref{thm:E} \tV (checked in examples). The dimension in each degree grows as a polynomial in $m$, while the total alternating sum is independent of $m$.

\subsection{Relation to the theorem of Pal--Poddar}\label{subsec:PPcompare}
\begin{thm}[Pal--Poddar, Thm.\ 11.12]\tC
Let $M$ be a compact regular-type generalized complex manifold, $M/\mathcal S$ a K\"ahler orbifold, and $E$ a positive SGH line bundle. (1) If $p+q>k$ then $H^q(M,(\mathcal G^*M)^p\otimes E)=0$. (2) For any SGH bundle $E'$, $H^q(M,E'\otimes E^m)=0$ for $m\ge m_0$, $q>0$.
\src{\cite{PP}, Thm.\ 11.12 (hypotheses and statement from a re-reading recorded in \texttt{pp\_recheck.md}).}
\end{thm}
\begin{prop}\tP \tQ\label{prop:PPcompare}
In setting (S), the groups $H^q(M,(\mathcal G^*M)^p\otimes E)$ of Pal--Poddar can be read as $H^q(X,\Omega^p_X\otimes G)$ (forms constant along the leaves). Then (1) is the classical Kodaira--Nakano vanishing on $X$, and (2) is the classical Serre vanishing for $E'=\pi^*(\wedge^bT_X\otimes V^j)$ (an SGH bundle), which is the same as the vanishing $H^a(X,\wedge^bT_X\otimes V^j\otimes G^m)=0\ (a>0)$ of Corollary \ref{cor:L}. Hence the vanishing in transversal degree coincides with Pal--Poddar Thm.\ 11.12(2) applied to $E'=\wedge^bT_X\otimes V^j$, and Corollary \ref{cor:L} is ``the Serre-type vanishing of Pal--Poddar $+$ Leray $+$ Lefschetz''.
\end{prop}
\begin{rem}[\tQ]
If one takes literal sheaf cohomology on $M$, an extra leafwise factor $H^j(Y)$ appears because $R^j\pi_*\pi^{-1}\ne0$. Since Pal--Poddar Thm.\ 8.9, Cor.\ 8.10 identify $H^q(M,(\mathcal G^*M)^p)$ with transversal Dolbeault cohomology, their $H^q$ is presumably to be read on the leaf space, but we have not checked whether the original text says so explicitly. Moreover, vanishing of $\HH^k$ holds only for transversal degree $a>0$; the factor $H^j_{dR}(Y)$ remains in each degree $j$.
\end{rem}

\subsection{Examples}
\begin{ex}[Direct product $Y\times\mathbb P^1$]\tP \tV\label{ex:prodY}
$X=\mathbb P^1$, $G=\mathcal O(m)$, $m\ge-1$: by Theorem \ref{thm:L} and $H^0(\mathbb P^1,\mathcal O(m))=m+1$, $H^0(\mathbb P^1,T\otimes\mathcal O(m))=m+3$, $P_{\HH}(t)=P_Y(t)\bigl((m+1)+(m+3)t\bigr)$. For $Y=S^2$: $(m+1)(1+t^2)+(m+3)(t+t^3)$ ($m=0$: $1,3,1,3$). For $Y=T^2$: $1,5,7,3$ at $m=0$ and in general $(m+1,\ 3m+5,\ 3m+7,\ m+3)$. In both cases the dimensions grow linearly in $m$.
\end{ex}
\begin{ex}[Kodaira--Thurston: $T^2\to M\to E$]\tP \tV \tQ
When the monodromy is the shear $\psi=\begin{pmatrix}1&1\\0&1\end{pmatrix}$, $V^0=V^2=\mathcal O$ and $V^1=F_2$ (Atiyah's non-trivial unipotent flat bundle \cite{Atiyah}, $h^0=h^1=1$ \tC). For $G=\mathcal O$, $P_{\HH}(t)=(1+t)^2(1+t+t^2)=1+3t+4t^2+3t^3+t^4$. Independent check: the Betti numbers of the Lie algebra complex obtained from the leafwise complex of the Heisenberg algebra $\mathfrak h_3$ times $\wedge N$ are $(1,3,4,3,1)$, in agreement \tV. This check assumes a Nomizu-type theorem saying that the model of invariant forms computes $\HH^*$, which we did not prove \tQ (Theorem \ref{thm:L} itself does not use it). For $\deg G=d>0$, $P_{\HH}(t)=d(1+t)^3$, and for $\deg G=0$, $G\not\cong\mathcal O$, $P_{\HH}=0$.
\end{ex}
\begin{ex}[Bieberbach quotient, $S^2$-bundles]\tP \tV
$(T^2\times E)/\mathbb Z_2$ ($(y,x)\mapsto(-y,x+\tau/2)$): $V^1=L_{-1}\oplus L_{-1}$ (non-trivial sign representation) is acyclic for $G=\mathcal O$, and $P_{\HH}(t)=(1+t)^2(1+t^2)$ (agrees with the $\mathbb Z_2$-invariant part of K\"unneth). For $S^2$-bundles with holonomy in $SO(3)$ the $V^j$ are trivial bundles, so the numbers are the same as for the direct product.
\end{ex}
\begin{rem}\tP
An SGH bundle with a non-trivial flat holonomy character $\chi$ along the leaves has $\chi$ locally constant over a compact base (maximum principle). If $Y$ is a torus and $\chi\ne1$ then $H^*(Y;L_\chi)=0$ and $\HH=0$.
\end{rem}
\tC The numerical values of these examples (Kodaira--Thurston, Bieberbach) are transcribed from the earlier records and were not recomputed while writing this paper \tQ.

%--------------------------------------------------------------
\subsection{Contrast with the type-jumping case: polynomial versus periodic}\label{subsec:contrast}
%--------------------------------------------------------------
\begin{center}\small
\begin{tabular}{>{\raggedright\arraybackslash}p{3.4cm}>{\raggedright\arraybackslash}p{4.4cm}>{\raggedright\arraybackslash}p{5.8cm}}\toprule
 & Regular type ($Y=S^2$, $X=\mathbb P^1$) & Type jumping ($\mathbb{CP}^2$, smooth cubic)\\\midrule
$\dim\HH^k$ for large $m$ & Linear in $m$: $(m+1),(m+3),(m+1),(m+3)$ & Bounded: $(1,0,2,0,0)$ for $m\in3\Z_{\ge0}$, $(0,0,3,0,0)$ otherwise\\\midrule
$m=0$ & $(1,3,1,3)$ & $(1,0,2,0,0)$\\\midrule
$m=1$ & $(2,4,2,4)$ & $(0,0,3,0,0)$\\\midrule
$m=3$ & $(4,6,4,6)$ & $(1,0,2,0,0)$\\\midrule
Total Euler characteristic & $\chi(Y)\cdot(-1)^{k}\chi(X)=-4$ ($k=\dim X=1$) & $\chi(\mathbb{CP}^2)=3$ (independent of $m$)\\\bottomrule
\end{tabular}
\end{center}
\tV The right column is Example \ref{ex:cp2} (agrees with direct computation for $m=-15,\dots,15$), the left column is Corollary \ref{cor:L}.
\begin{obs}\tP\label{obs:contrast}
Both start from $E_1^{a,b}=H^a(\wedge^bT\otimes E)$. For regular type with SGH bundles $N$ acts trivially, so $d_1=0$ (Theorem \ref{thm:L}). In the type-jumping case $d_1=[\sigma,\cdot]+V\wedge$ is \emph{almost exact} between the $H^0$ terms, and although $H^0(\wedge^bT\otimes L^m)$ grows with $m$, the remaining cohomology is a bounded quantity determined by the topology of $U=X\setminus D$. The difference between ``polynomial in $m$'' and ``periodic'' reduces to whether $d_1$ vanishes or is (almost) exact.
\end{obs}
\begin{conj}[Unification (imprecise)]\tK\label{conj:U}
For a compact generalized complex manifold with finitely many type strata and a module $E$ on which $N$ acts trivially along the holomorphic strata, $\HH^*(E)$ is filtered by stratumwise contributions of ``leafwise de Rham $\otimes$ transversal Dolbeault on the leaf space'', and for each stratum whether it is of $(c)$-type (compact support) or not is determined by the indicial roots (residues).
\end{conj}
\noindent This conjecture is not formulated precisely, and there is no verification beyond the two families above (regular type and log-symplectic). The dictionary is as follows.
\begin{center}\small
\begin{tabular}{>{\raggedright\arraybackslash}p{3.0cm}>{\raggedright\arraybackslash}p{5.4cm}>{\raggedright\arraybackslash}p{5.6cm}}\toprule
 & Regular type & Type jumping (log-symplectic)\\\midrule
Leafwise de Rham & $H^j(Y)$ of the compact leaf $Y$ (local systems $V^j$) & $H^j(U;L_U)$ of the open leaf $U$\\\midrule
Transversal & Dolbeault on $X$ & Complexes $Q_T$ on the strata of $D$\\\midrule
Treatment of the boundary & None ($Y$ compact) & Choice of $Rj_*$ / $j_!$ by the residues (Deligne)\\\midrule
Monodromy & Gauss--Manin & $e^{2\pi\sqrt{-1}r}$ (around $D$)\\\midrule
Degeneration & $d_r=0$ by Lefschetz & $U$ is non-compact with no analogue, $d_1\ne0$\\\bottomrule
\end{tabular}
\end{center}

%==============================================================
\section{The Kodaira-type problem}\label{sec:kodaira}
%==============================================================
\subsection{The classical Kodaira embedding theorem: where what is used}
\begin{thm}[Kodaira embedding theorem]\tC\label{thm:Kod}
Let $X$ be a compact complex manifold and $L$ a positive holomorphic line bundle. Then $X$ is projective, and for $m$ sufficiently large the map $X\to\mathbb P(H^0(X,L^m)^*)$ given by $|L^m|$ is an embedding.
\src{Kodaira \cite{Kod54}. Textbooks: Griffiths--Harris \cite{GH} Ch.~1 \S4, Huybrechts \cite{Huy} Prop.\ 5.3.1 (number \tQ).}
\end{thm}
\begin{hyptab}
$X$ compact & $H^0(X,L^m)$ is finite-dimensional and the map to $\mathbb P^N$ can be defined. Closed embedding & First and last steps of the proof\\\midrule
$L$ positive ($i\Theta_L>0$) & The curvature form is a K\"ahler form and Kodaira vanishing is available & The step showing $H^1(\tilde X,\pi^*L^m\otimes\mathcal O(-2E_x))=0$ on the blow-up $\tilde X$ at a point $x$\\\midrule
Enough holomorphic sections & They separate points and tangent vectors & Surjectivity of $H^0(L^m)\to L^m\otimes\mathcal O/\mathfrak m_x^2$ follows from the vanishing
\end{hyptab}
It matters at which step of the proof it is used that sections \emph{distinguish points}. This step breaks for generalized holomorphic (GH) sections.

\subsection{Embedding by generalized holomorphic sections is impossible}
\begin{prop}\tP\label{prop:one}
Let $X$ be a connected compact complex manifold, $\sigma$ a holomorphic Poisson structure non-degenerate on a dense open set $U_0$, and $L$ a Poisson module. Then $\dim_\C\HH^0(L)\le1$. A GH section is determined by its value at one point of $U_0$, and the Kodaira map by GH sections is constant (or undefined).
\end{prop}
\begin{hyptab}
Connected, $U_0$ dense & The identity theorem can be used & (iii)\\\midrule
$\sigma$ non-degenerate on $U_0$ & $\sigma^\sh:A|_{U_0}\to T$ is an isomorphism and the $A$-connection becomes an ordinary flat connection & (i)\\\midrule
$L$ a line bundle & A flat section is determined by its value at one point & (ii)
\end{hyptab}
\begin{proof}
(i) On $U_0$ put $\nabla_vs:=D_{(\sigma^\sh)^{-1}v}s$; by the Leibniz rule $\nabla$ is an ordinary flat holomorphic connection. GH sections satisfy $\nabla s=0$ (Observation \ref{obs:casimir}). (ii) $U_0$ is connected (a connected manifold minus a proper analytic subset) and $L$ has rank 1, so the space of flat sections has dimension $\le1$ and a section is determined by its value at one point. (iii) $s$ is determined by $s|_{U_0}$ (identity theorem). \qedhere
\end{proof}
\begin{ex}\tV
$\mathbb{CP}^2$ with $\Pi_F$ ($F$ a smooth cubic): $\HH^0(\mathcal O(m))=\C F^k$ ($m=3k\ge0$) and $0$ otherwise (Example \ref{ex:cp2}). $\mathbb P^1\times\mathbb P^1$ with a smooth $(2,2)$ curve: $\HH^0(\mathcal O(m,m))=\C F^{m/2}$ ($m\in2\Z_{\ge0}$, Example \ref{ex:p1p1}). For no $m$ is there an embedding by GH sections.
\end{ex}
\begin{prop}\tP\label{prop:leaf}
Let $M$ be a compact generalized complex manifold, $\Jc$ of regular type $k<n$ on a dense open set $W$ (leaf dimension $2(n-k)>0$), and $L$ a generalized holomorphic line bundle. For any GH sections $s_0,\dots,s_N$, the map $\varphi=[s_0:\cdots:s_N]$ (defined outside the base locus) is constant on each leaf of $W$ and hence not an immersion on $W$.
\end{prop}
\begin{hyptab}
Regular type (local model is a product) & $\pi_T(L)\supset T_F\otimes\C$ can be used & Local Darboux (Gualtieri \cite{Gua11} Thm.\ 4.3 \tC)\\\midrule
$k<n$ & Leaves have positive dimension & Can move along leaves\\\midrule
$L$ generalized holomorphic & The $L$-connection $D$ satisfies $D_\xi(fs)=\pi_T(\xi)(f)s+fD_\xi s$ & GH section $\Rightarrow$ parallel along the leaves
\end{hyptab}
\begin{proof}
In the product model $L=\mathrm{graph}(\omega_Y)\oplus(T^{0,1}\oplus T^{*1,0})$ we have $\pi_T(L)=T_Y\otimes\C\oplus T^{0,1}$. When $\xi\in L$ has a leafwise component, $D_\xi s=0$ means $s$ is parallel for the leafwise flat connection $\nabla^F$. For parallel transport $P:L_x\to L_y$ along a path in a leaf, $s_i(y)=Ps_i(x)$ for all GH sections. Hence the ratio $[s_0:\cdots:s_N]$ agrees at $x$ and $y$. \qedhere
\end{proof}
\begin{rem}[What breaks in the classical proof]\tP
The classical proof uses that sections separate points and tangent vectors. In Proposition \ref{prop:one} there is at most one separating section (up to scalars), and in Proposition \ref{prop:leaf} points of the same leaf cannot be separated. The hypothesis (positive curvature) guarantees only the \emph{existence} of sections, not that they are GH (flat). Pal--Poddar's ``positive SGH line bundle'' (Def.\ 11.10, 11.11 \tC \tQ) requires positivity only in the transversal direction, which corresponds to this situation.
\end{rem}

\subsection{Leaf-space version of Kodaira}
\begin{thm}[Leaf-space version of Kodaira]\tP\label{thm:K1leaf}
Let $M$ be a compact regular-type generalized complex manifold satisfying:
\begin{itemize}
\item[(K1)] The projection $\pi:M\to X$ to the symplectic leaves is a proper submersion (or orbifold submersion) with connected compact fibres ($=$ leaves), and $X$ is a compact complex manifold (orbifold) of dimension $k$.
\item[(K2)] $L$ is an SGH line bundle and there is a holomorphic line bundle $L_X\to X$ with $L\cong\pi^*L_X$ (as SGH bundles) (holonomy along the leaves is trivial).
\item[(K3)] $L_X$ is ample.
\end{itemize}
Then there is $m_0$ such that for $m\ge m_0$, $\HH^0(M,L^m)=\pi^*H^0(X,L_X^m)$, and the map $\varphi_m:M\to\mathbb P^N$ by GH sections is $\varphi_m=\iota_m\circ\pi$, where $\iota_m:X\hookrightarrow\mathbb P^N$ is the classical Kodaira embedding. Hence each fibre of $\varphi_m$ is exactly one leaf. In the orbifold case $\iota_m$ is an embedding of the coarse space ($m$ a multiple of the orders of the isotropy groups; Baily \cite{Bai}, quoted from memory \tQ).
\end{thm}
\begin{hyptab}
(K1) connected compact fibres & Flat sections on a leaf are constant (converse of Prop.\ \ref{prop:leaf}). Since fibres are connected the preimage under $\varphi_m$ is exactly one leaf & (i), (iii)\\\midrule
(K2) trivial holonomy & $L^m$ is a trivial flat bundle on leaves and GH sections exist. If the holonomy is non-trivial, $\HH^0=0$ along the leaves & (i)\\\midrule
(K3) $L_X$ ample & Classical Kodaira (Theorem \ref{thm:Kod}) is applied on $X$ & (ii)\\\midrule
Orbifold case & Kodaira for the coarse space (Baily) & Orbifold version of (ii)
\end{hyptab}
\begin{proof}
(i) $\HH^0(M,L^m)$ consists of the sections of $L^m$ that are flat along leaves and holomorphic transversally (definition of SGH bundles; degree-0 part of Pal--Poddar Cor.\ 8.10 \tC \tQ). By (K1), (K2), $L^m=\pi^*L_X^m$ is a trivial flat bundle on each fibre, so flat sections are constant on fibres, i.e.\ of the form $\pi^*t$ ($t$ a section of $L_X^m$). Transversal holomorphy is equivalent to holomorphy of $t$. Hence $\HH^0(M,L^m)=\pi^*H^0(X,L_X^m)$. (ii) By (K3) and Theorem \ref{thm:Kod}, $\iota_m$ is an embedding for $m\ge m_0$. (iii) $\varphi_m=\iota_m\circ\pi$ and $\iota_m$ is injective, so $\varphi_m(x)=\varphi_m(y)\iff\pi(x)=\pi(y)$, i.e.\ $x,y$ are in the same leaf. \qedhere
\end{proof}
\begin{rem}[Position]
\tP This is a direct consequence and not a deep result. \tC As far as we found, no reference states a Kodaira-type claim in regular type. It corresponds, on the embedding side, to Pal--Poddar Thm.\ 11.12. \tQ Equivalence with their definition of ``positive'' is unchecked. We found no Kodaira-type statements for the type-jumping case or for general generalized complex structures.
\end{rem}

%--------------------------------------------------------------
\subsection{Poisson--Kodaira: embedding of Poisson pairs}
%--------------------------------------------------------------
Since GH sections are (almost) constant, we change the problem and ask whether a Kodaira embedding $\nu:X\hookrightarrow\mathbb P^N$ by \emph{ordinary} holomorphic sections is a Poisson embedding, i.e.\ whether $\sigma$ is the restriction of a Poisson structure $\sigma_N$ on $\mathbb P^N$. (This is the author's proposal, not a formulation in the literature.)

\begin{lem}[$\C^*$-invariant lift of a rank one Poisson module]\tP \tQ\label{lem:lift}
Let $L$ be a Poisson module. On the total space $C=L^*\setminus0$ there is a unique $\C^*$-invariant Poisson structure $\hat\pi=\sigma+u\,\theta\wedge\partial_u$ with $p_*\hat\pi=\sigma$ (local frame $e$, $D_\alpha e=\theta(\alpha)e$, $u$ the fibre coordinate), and sections of $L^k$ correspond to homogeneous functions of degree $k$. The Jacobi identity for $\hat\pi$ is equivalent to flatness of $D$ and the Leibniz rule.
\end{lem}
The local computation is checked only as a sketch \tQ. Example: $\Pi_F$ on $\mathbb P^2$ (a quadratic Poisson structure on $\mathbb C^3$), $L=\mathcal O(m)$ (any $m$).
\begin{prop}\tP\label{prop:graded}
If $X$ is projective and $L$ an ample Poisson module, then $R=\bigoplus_{k\ge0}H^0(X,L^k)$ is a graded Poisson algebra whose bracket has degree 0 ($\{R_a,R_b\}\subset R_{a+b}$).
\end{prop}
\begin{proof}
By Lemma \ref{lem:lift}, $R$ is the ring of homogeneous functions on $C$ and the bracket of $\hat\pi$ preserves degree. \qedhere
\end{proof}
\begin{prop}[Almost-Poisson quadratic lift]\tP \tC\label{prop:almost}
For $m\ge m_1$, $R^{(m)}=\bigoplus_kR_{mk}=S/I$ ($S=\mathrm{Sym}\,R_m$, $I$ generated by quadrics) (Mumford \cite{Mum}; quoted from memory, pages \tQ). Then there is a degree-0 skew biderivation $\{\ ,\ \}'$ on $S$ (a quadratic almost-Poisson structure) inducing the bracket of $R$ on the quotient $S/I$.
\end{prop}
\begin{proof}
Choose a linear section $s$ of $S_2\to R_{2m}$, set $\{y_a,y_b\}:=s(\{\bar y_a,\bar y_b\})$ and extend by the Leibniz rule. For $q\in I_2$, $g\in S_1$, $\{q,g\}'\equiv\{\bar q,\bar g\}=0\bmod I_3$, so $\{I_2,S_1\}'\subset I_3$. Since $I$ is generated by $I_2$, the Leibniz rule gives $\{I,S\}'\subset I$. \qedhere
\end{proof}
The Jacobiator $\mathcal J\in\mathrm{Hom}(\wedge^3S_1,I_3)$ of the lift is in general non-zero, and the difference between choices of lift is $\epsilon\in\mathrm{Hom}(\wedge^2S_1,I_2)$. Extension to a Poisson structure is equivalent to the existence of a solution $\epsilon$ of the quadratic equation $\mathcal J_0+\mathcal L(\epsilon)+\mathcal J(\epsilon)=0$. The first-order obstruction class is $\mathcal J_0\bmod\mathrm{im}\,\mathcal L$.

\begin{thm}[Poisson--Kodaira, toric]\tP\label{thm:PKtoric}
Let $X$ be a smooth projective toric variety (torus $T=N\otimes\C^*$) and $\sigma=\sum\omega^{ij}\xi_i\wedge\xi_j$ ($\xi_i=x_i\partial_{x_i}$, $\omega\in\wedge^2N_\C$). Let $L=\mathcal O(D)$ be ample with polytope $P$ and $L^m$ very ample (sections $\chi^a$, $a\in mP\cap M$). Then $\nu:X\hookrightarrow\mathbb P^N$, $\nu^*y_a=\chi^a$, is a Poisson embedding for the diagonal quadratic Poisson structure on $\mathbb P^N$
\[\sigma_N=\sum_{a<b}\omega(a,b)\,y_ay_b\,\partial_{y_a}\wedge\partial_{y_b},\qquad\omega(a,b)=\sum\omega^{ij}a_ib_j.\]
This holds for any very ample power $L^m$.
\end{thm}
\begin{hyptab}
Toric, $\sigma$ torus invariant & Brackets of characters are monomials: $\{\chi^a,\chi^b\}=\omega(a,b)\chi^{a+b}$ & (2). The lift can be chosen diagonal\\\midrule
$L^m$ very ample & $\nu$ is an embedding & The ``embedding'' part of the conclusion\\\midrule
$X$ smooth, projective & The setting of classical Kodaira & Same
\end{hyptab}
\begin{proof}
(1) $\sigma_N$ is Poisson: for skew $\lambda_{ij}$ put $\{y_i,y_j\}=\lambda_{ij}y_iy_j$; then $\{y_i,\{y_j,y_k\}\}=\lambda_{jk}(\lambda_{ij}+\lambda_{ik})y_iy_jy_k$ and the cyclic sum vanishes because the six terms cancel in pairs by skew-symmetry. The scaling weight is 0, so it descends to $\mathbb P^N$. (2) On the torus, $\{\chi^a,\chi^b\}=\sigma(d\chi^a,d\chi^b)=\omega(a,b)\chi^{a+b}=\nu^*(\omega(a,b)y_ay_b)$. Since $\sigma$ and $\nu^*\sigma_N$ are both holomorphic, $\nu$ is a Poisson map also on the closure $X$ of the torus. \qedhere
\end{proof}

\begin{obs}[$\mathbb P^2$, cubic, $m=2$: numerical confirmation \tV]\label{obs:P2}
$X=\mathbb P^2$, $\sigma=\Pi_F$, Veronese $m=2$ ($\mathbb C^6$, $\dim I_2=6$, 90 unknowns $\epsilon$, 940 equations). The Jacobiator of the natural lift is non-zero for 18 triples for Fermat, 19 for Hesse, 1 for $xyz$, and also for the cusp. Solving the full quadratic equation $\mathcal J_0+\mathcal L(\epsilon)+\mathcal J(\epsilon)=0$ by the Levenberg--Marquardt method, solutions with residual $\le10^{-12}$ were found for all of $F=xyz$, Fermat, Hesse ($x^3+y^3+z^3-21xyz$) and cusp ($x^3-y^2z$). For Fermat we confirmed also independently by finite differences (Jacobiator $\sim10^{-9}$, error of restriction to the Veronese cone $\sim10^{-10}$). For Fermat there is no solution of the linear part $\mathcal J_0+\mathcal L(\epsilon)=0$ alone, and the quadratic term is essential. $m=1$ is trivial; for $m\ge3$ there are 1215 unknowns and the computation is too heavy by this method and was not done. Scripts \texttt{roundE\_*.py} (transcribed from the earlier records, not re-run \tQ). This is numerical existence, not an exact algebraic proof (certificate).
\end{obs}
\begin{obs}[Segre embedding of $\mathbb P^1\times\mathbb P^1$ ($m=1$): Sklyanin-type structure \tV]\label{obs:Segre}
For the quadratic Poisson structure on $\mathbb P^3$, $\{x_0,x_i\}=(J_j-J_k)x_jx_k$, $\{x_j,x_k\}=x_0x_i$ ($(i,j,k)$ cyclic in $1,2,3$, $J=(2,3,7)$), we confirmed in exact rational arithmetic the Jacobi identity (4 of 4 triples) and that there are exactly two quadratic Casimirs and the generic rank is 2. Hence each quadric surface of the Casimir pencil is a Poisson submanifold, the smooth ones are $\mathbb P^1\times\mathbb P^1$, and the zero locus of $\sigma$ is a smooth elliptic quartic curve $E=Q_1\cap Q_2$ (``the simple zeros are exactly $E$'' is an identification by rank and leaf arguments, not computed). This is classical and not new (work of Feigin--Odesskii, Odesskii, Hua--Polishchuk \tC \tQ). Whether all smooth $(2,2)$ curves arise in this form, and the case of singular $E$, are unchecked \tK.
\end{obs}
\begin{conj}[Poisson--Kodaira]\tK\label{conj:PK}
Let $(X,\sigma)$ be a projective holomorphic Poisson manifold and $L$ an ample Poisson module. Then there is $m_0$ such that for $m\ge m_0$ the Kodaira embedding by $|L^m|$ is a Poisson embedding for some quadratic Poisson structure $\sigma_N$ on $\mathbb P^N$.
\end{conj}
\noindent Evidence: Theorem \ref{thm:PKtoric}, Observations \ref{obs:P2}, \ref{obs:Segre}. Open: a general obstruction theory ($\mathcal J_0\bmod\mathrm{im}\,\mathcal L$ is a class of Poisson cohomology type, expected to vanish for $m\gg0$ by Serre-type vanishing for $H^1$ of a Koszul complex, but not proved). For the Fermat type of $\mathbb P^2$ at $m=3$ there is no computation and no evidence either for or against PK (the case $xyz$ is proved by the toric theorem and is not an independent test). No counterexample has been found, and the possibility of failure for small $m$ remains. For $\sigma=0$ it agrees with classical Kodaira. For regular type $M\to X$ with leaf $Y$, GH sections factor through $X$ (Proposition \ref{prop:leaf}), so to embed $(M,\sigma_Y)$ as a Poisson pair one must use non-GH ordinary holomorphic sections; this is the formulation of this subsection.

%==============================================================
\section{Riemann--Roch-type formulas and the index conjecture}\label{sec:RR}
%==============================================================
\subsection{The Poincar\'e polynomial formula}
Theorem \ref{thm:E} says that the total Euler characteristic is topological. The dimensions in each degree take the following form.

\begin{thm}[RR-type Poincar\'e polynomial]\tP \tV\label{thm:K}
In addition to the hypotheses of Theorem \ref{thm:R} (SNC, log-symplectic, $r_i\notin\Z_{\ge1}$, (GEN$_Q$)), assume (i) $U$ is affine and (ii) $n$ is even (automatic for log-symplectic). Then
\[P_{\HH}(t)\;=\;\chi(X)\,t^n\;+\;\sum_{k=0}^{n-1}h^k(U;L_U)\,\bigl(t^k-(-1)^kt^n\bigr).\]
\end{thm}
\begin{hyptab}
Hypotheses of the reduction theorem & Give $\HH^*\cong H^*(U;L_U)\oplus\C^{\#D_{(n)}}[-n]$ & (1)\\\midrule
$U$ affine & Artin vanishing: $H^k(U;L_U)=0\ (k>n)$ for local systems & (2). Without it the formula does not reduce to a sum over $k<n$\\\midrule
$n$ even & $\sum_k(-1)^kh^k+\#D_{(n)}=\chi(X)$ can be used without signs (Theorem \ref{thm:E}) & (3)\\\midrule
Projectivity, ampleness & Not needed & ---
\end{hyptab}
\begin{proof}
(1) By Theorem \ref{thm:R}, $P_{\HH}=\sum_{k\le n}h^kt^k+\#D_{(n)}t^n$ (the terms with $k>n$ vanish by (2)). (2) By Artin's vanishing theorem \cite{SGA4}, $h^k=0$ for $k>n$. (3) By Theorem \ref{thm:E} ($n$ even), $\sum_{k\le n}(-1)^kh^k+\#D_{(n)}=\chi(X)$, i.e.\ $h^n+\#D_{(n)}=\chi(X)-\sum_{k<n}(-1)^kh^k$. Substitute into the coefficient of $t^n$. \qedhere
\end{proof}
\begin{cor}\tP
If all $r_i\notin\Z$ then (Theorem \ref{thm:V}) $P_{\HH}(t)=\chi(X)t^n$. If all $r_i\in\Z_{\le0}$ and $L_U$ is trivial then $h^k=b_k(U)$. For $n=2$, $h^0=1$ and $h^1=b_1$, and $P=\chi(X)t^2+(1-t^2)+b_1(t+t^2)$.
\end{cor}
\paragraph{Verification.} \tV (a) 188 examples in the tables for $n=2$, (b) 61 toric examples with $n=2,4$ ($h^k=\binom nk$ for $\rho\in M$, 0 otherwise): the measured Poincar\'e polynomials agree with the predictions of the formula in all cases (0 discrepancies; earlier records, not re-run \tQ). For toric $\mathbb P^2$ with $\rho\in M$, $P=3t^2+(1-t^2)+2(t+t^2)=1+2t+4t^2$. This agrees with the case $[r\le0]$ of Theorem \ref{thm:B}. The side $r\ge2$ follows from duality (Theorem \ref{thm:dual}).

\subsection{Three behaviours of the dimensions}
Taking $L=L_0^m$ and varying $m$, the degree-wise dimensions behave in qualitatively different ways depending on the setting.
\begin{center}\small
\begin{tabular}{>{\raggedright\arraybackslash}p{3.0cm}>{\raggedright\arraybackslash}p{4.4cm}>{\raggedright\arraybackslash}p{6.4cm}}
\toprule Behaviour & Dimension in degree $k$ & Example\\\midrule
(A) Regular type & \emph{Polynomial} in $m$ ($m\gg0$) \tP & $X=\mathbb P^1$, $Y=T^2$ (product), $G=\mathcal O(m)$: $(m+1,\ 3m+5,\ 3m+7,\ m+3)$ ($m\ge-1$)\\\midrule
(B) Log-symplectic & \emph{Periodic} jumps according to integrality of the residue $r=-m/3$ \tP \tV & $\mathbb P^2$, smooth cubic: $(1,0,2,0,0)$ for $m=3k\ge0$, $(0,0,3,0,0)$ otherwise for $m\ge-3$ (Example \ref{ex:cp2})\\\midrule
(C) Total Euler characteristic & \emph{Independent} of $m$ \tP & In (A) $\sum(-1)^kh^k=0$, in (B) $3=\chi(\mathbb P^2)$\\\bottomrule
\end{tabular}
\end{center}
In (A), $(m+1)-(3m+5)+(3m+7)-(m+3)=0=\chi(Y)\chi(X)$, consistent. How to read this: by (C), the naive form of Riemann--Roch, ``$\chi(E)=\int\ch(E)\Td$ is a polynomial in $c_1(E)$'', degenerates in generalized geometry (it does not see $c_1$), and the information lies in the graded dimensions. In (A) these are polynomials degree by degree; in (B), as in Theorem \ref{thm:K}, they decompose into the topological quantities $h^k(U;L_U)$ (depending only on the integrality of residues) and $\chi(X)$. Only in regular type does polynomial behaviour hold in all degrees.

\subsection{Comparison with earlier work}\label{subsec:compare}
\paragraph{Hodge theory (formal coefficients).} \tC Gualtieri proved the generalized Hodge star, the Born--Infeld pairing, equality of Laplacians, Hodge decomposition and the $dd^{\Jc}$-lemma on generalized K\"ahler manifolds (\cite{GuaHodge}, \cite{GuaK}, Cavalcanti \cite{Cav}; theorem numbers not read closely \tQ). These are theories with trivial or $H$-twisted coefficients, and a Bochner--Kodaira--Nakano identity with (curved) line bundle coefficients is not written, as far as we found. The results of this paper can be seen as filling this gap for the $\beta$-transform case, without Hodge theory, by comparing the log complex with constructible sheaves.

\paragraph{Pal--Poddar.} \tC \cite{PP} treat regular-type (type $k$ constant) compact generalized complex manifolds whose leaf space is a K\"ahler orbifold, and give Hodge theory for SGH bundles, the transversal Dolbeault isomorphism (Thm.\ 8.9, Cor.\ 8.10), Hodge decomposition and Serre duality (Thm.\ 11.2--11.7), K\"ahler identities (Prop.\ 11.8, Lemma 11.9), and vanishing theorems for positive SGH line bundles (Def.\ 11.10, 11.11; Thm.\ 11.12) (numbers \tQ). Differences: (i) the type-jumping case is outside Pal--Poddar, and examples such as $\mathbb{CP}^2$ in this paper are of this kind. (ii) SGH bundles are a proper subclass of generalized holomorphic bundles, while Poisson modules in this paper are general GH bundles for the $\beta$-transform structure. (iii) The precise statement of Thm.\ 11.12 is in \S\ref{subsec:PPcompare}. Doubts about the original text (probably misprints or omissions, unverified \tQ): $(i/2\pi)$ in the statement versus $1/2\pi$ in the proof; the sign of $gc_1$; that ``a positive class yields a metric of positive curvature'' is not proved from Def.\ 11.11 and is assumed in the proof of Thm.\ 11.12; that compactness and the K\"ahler property are not stated in Def.\ 11.11.

\paragraph{Hu's work.} \tC Hu \cite{Hu1,Hu2}. The setting is a real Lie algebroid $\mathfrak g$ with a fibrewise integrable complex structure (different from the complex Lie algebroid $L\subset\bT_\C$ of this paper). Riemann--Roch (2403.13424 Thm.\ 4) asserts an index formula under ellipticity of $\dbar_{\mathfrak g}$ and unimodularity, and the vanishing (Thm.\ 7) is a statement about the kernel of the Laplacian, from which vanishing of cohomology cannot be derived. \tP The generalized Dolbeault complex $(T^{0,1}\bowtie A^{1,0})$ of this paper is not in Hu's setting ($L$ is a complex Lie algebroid without $J$, and $\wedge L^*$ in general has no $(p,q)$ bigrading). \tK Hu's index is expected to depend polynomially on $m$, whereas the total Euler characteristic here is independent of $m$, and the residue $-m/3$ that determines the periodic jumps is not part of Hu's data. Hence Hu's theorem is not ``Riemann--Roch for generalized Dolbeault cohomology''. Inconsistencies in the original (undefined symbols, sign discrepancies): noticed in the summary and part of the text, but the appendix could not be verified \tQ.

\paragraph{Integrability.} \tC Gualtieri--Li \cite{GL} classify Hausdorff symplectic groupoids of log-symplectic manifolds and show that the Lie algebroid $T_X(-\log D)$ of log vector fields tangent to a smooth $D$ is integrable. Bailey--Gualtieri \cite{BG}, Laurent-Gengoux--Sti\'enon--Xu \cite{LGSX2}, and Crainic--Fernandes \cite{CF} give general theory of integrability. These supply groups as input for index theory but do not treat the integration of the complex Lie algebroid $L$ itself (theorem numbers \tQ).

\subsection{The index conjecture}
The Hirzebruch--Riemann--Roch theorem is a special case of the Atiyah--Singer index theorem \cite{AS}, and Riemann--Roch in generalized geometry should be formulated as ``an index theorem for the generalized Dolbeault complex''. What follows is \emph{entirely conjectural and unproved} (\tK).

\paragraph{Why existing index theorems do not apply.} (1) The complex of Hu \cite{Hu2} is attached to a real Lie algebroid, whereas our complex $(\wedge^\bullet T\otimes L,\dbar+[\sigma,\cdot])$ is the complex of the complex Lie algebroid $T^{0,1}\ltimes A^{1,0}$ (no real structure). (2) The index of Pflaum--Posthuma--Tang \cite{PPT} on Lie groupoids requires a compact base, unimodularity, etc.\ \tQ. (3) We found no index theorem in the literature for log Dolbeault complexes weighted by the residues of $L$. Related literature: geometric quantization of $b$-symplectic manifolds (Braverman--Loizides--Song \cite{BLS}), index formulas for boundary groupoids (Qiao--So \cite{QS}), pseudodifferential operators on manifolds with Lie structure (Ammann--Lauter--Nistor \cite{ALN}), fibred boundaries (Mazzeo--Melrose \cite{MM}), complex logarithmic Hodge theory (Hausel--Hunsicker--Mazzeo \cite{HHM}). All have different settings and do not apply directly (only titles, authors and abstracts were checked \tQ).

\paragraph{Framework.} Regard $A=T_X(-\log D)$ as a Lie algebroid (Gualtieri--Li \cite{GL}: integrable for log-symplectic, checked for smooth $D$ \tC \tQ), and the complex $\mathcal D_L=(\wedge^\bullet T\otimes L,\ \dbar+[\sigma,\cdot])$ on $U$ as a closed operator on weighted Sobolev spaces determined by $A$ (the weight is the residue $r$). The principal symbol in the sense of $A$ is invertible (Lemma \ref{lem:iso} and a local computation \tV), so non-Fredholmness arises from indicial roots at the boundary. In the local model (a smooth point of $D$) $\nabla z^k=(k-r)z^k\,dz/z$, so the indicial root is $k=r$, which exists only for $r\in\Z$.

\begin{conj}[Index conjecture]\tK\label{conj:index}
Let $X,\sigma,D,L$ be as above (log-symplectic, SNC).
\begin{itemize}
\item[(I1)] The closed extension of $\mathcal D_L$ on Sobolev spaces of weight $r$ is Fredholm if and only if all $r_i\notin\Z$.
\item[(I2)] When Fredholm, the index (alternating sum) equals $\chi(U)+(-1)^n\#D_{(n)}=(-1)^n\chi(X)$ ($\int c_n(T(-\log D))$ plus the corner contribution $\#D_{(n)}$).
\item[(I3)] For $r_i\in\Z$, the difference of kernels and cokernels of the two closed extensions of $j_!$-type and $Rj_*$-type is given by a subquotient of $H^*(\partial U;L)$ (an APS-type wall-crossing triangle), and since $\chi(\partial U)=0$ the index does not change (consistent with Corollary \ref{cor:wall}).
\item[(I4)] Graded refinement: $P_{\HH}(t)$ is given by the form of Theorem \ref{thm:K} ($\chi(X)t^n$ and the jump terms determined by the kernels of the indicial roots).
\end{itemize}
\end{conj}
\paragraph{Correspondence with what is already proved.} The value in (I2) is proved by Theorems \ref{thm:E}, \ref{thm:K1} ($\chi(X)$, \tP). The triangle in (I3) is used in this paper at the topological level as $j_!L_U\to Rj_*L_U\to i_*i^*Rj_*L_U$ (completely checked for $n=2$ \tP \tV), and its index-theoretic counterpart is unproved. (I4) is Theorem \ref{thm:B} for $n=2$ \tP, and Theorem \ref{thm:K} (under (GEN$_Q$)) in general dimension \tP.
\paragraph{What remains to be proved.} (1) Fredholm theory on the complex Lie algebroid (without $J$): construction of a parametrix and the ``$\Leftrightarrow$'' of (I1). Extending the framework of Ammann--Lauter--Nistor to the weighted setting seems natural but has not been done. (2) The index formula (I2): an analogue of the Atiyah--Singer formula on the $b$-tangent bundle \cite{BLS} for complexes with $\wedge^\bullet T$ coefficients and weights. The contribution of corners (strata of higher codimension) is the problem. (3) The APS-type wall-crossing of (I3): identification with $H^*(\partial U;L)$ beyond smooth $D$ with $n=2$ (Gysin sequence). (4) The graded version (I4): graded refinement of spectral flow (no literature found). (5) A general sufficient condition for (GEN$_Q$).

\subsection{The three treasures: where we stand}
\begin{center}\small
\begin{tabular}{>{\raggedright\arraybackslash}p{2.4cm}>{\raggedright\arraybackslash}p{6.5cm}>{\raggedright\arraybackslash}p{6.0cm}}
\toprule Treasure & What holds in another form & What fails / is open\\\midrule
Riemann--Roch & Total Euler characteristic (Thms.\ \ref{thm:E}, \ref{thm:K1}), Poincar\'e polynomial formulas (Thms.\ \ref{thm:B}, \ref{thm:K}), degree-wise polynomials for regular type (Cor.\ \ref{cor:L}) & The naive HRR fails. Unification by an index theorem is a conjecture (Conj.\ \ref{conj:index})\\\midrule
Serre vanishing & Vanishing in transversal degree $a>0$ (Thm.\ \ref{thm:serre}) & Vanishing in all degrees fails (Remark \ref{rem:total})\\\midrule
Kodaira embedding & Leaf-space version (Thm.\ \ref{thm:K1leaf}), toric Poisson--Kodaira (Thm.\ \ref{thm:PKtoric}) & Embedding by GH sections is impossible (Prop.\ \ref{prop:one}). General Poisson--Kodaira is a conjecture (Conj.\ \ref{conj:PK})\\\bottomrule
\end{tabular}
\end{center}

%==============================================================
\section{Open problems}\label{sec:open}
%==============================================================
\begin{enumerate}
\item \textbf{General Poisson--Kodaira} (Conjecture \ref{conj:PK}): $\mathbb P^2$ with $m\ge3$, cases other than the Segre embedding of $\mathbb P^1\times\mathbb P^1$, singular $E$. Vanishing of the obstruction class $\mathcal J_0\bmod\mathrm{im}\,\mathcal L$.
\item \textbf{General sufficient conditions for (GEN$_Q$)}, proof of Proposition \ref{prop:grQ} for general $n$, the holomorphic logarithmic Darboux theorem (Proposition \ref{prop:genNec}(b)).
\item \textbf{Vanishing for mixed residues}: $H^k(U;L_U)=0\ (k<n)$ when only some $r_i$ are integers (extension of Theorem \ref{thm:V}).
\item \textbf{Local formulas}: non-quasi-homogeneous isolated singularities, non-isolated singularities, $s=1$, non-reduced boundary (Conjecture \ref{conj:NR}), globalization of the refined Riemann--Roch (Conjecture \ref{conj:refined}). An explanation of the triple point anomalies at $m=-2,-4$.
\item \textbf{A general formula at resonance $r_i=1$} (Observation \ref{obs:res}) and its proof.
\item \textbf{The index conjecture} (Conjecture \ref{conj:index}): Fredholmness in weighted $b$-calculus, index $=\chi(U)$, relation with the cohomology of the boundary.
\item \textbf{Extensions of regular type}: non-flat bundles, fibres failing hard Lefschetz, a precise formulation of the unification conjecture (Conjecture \ref{conj:U}).
\item \textbf{Riemann--Roch and Kodaira-type theorems for general generalized complex structures} (not attempted in this study). A Bochner--Kodaira--Nakano identity with line bundle coefficients on the generalized K\"ahler side.
\item \textbf{Resolution of unverified literature items}: Mokri's theorem, re-derivation of LGSX Lemma 4.22, a close reading of Gualtieri \cite{Gua11} \S4, remaining bibliographic data (list in the appendix).
\item \textbf{Verification in dimension 4 beyond toric and product cases.} A candidate: an elliptic Poisson structure of Feigin--Odesskii type on $\mathbb P^4$ ($D$ is not SNC).
\end{enumerate}

%==============================================================
\appendix
\section{Changes from the comprehensive version}\label{app:changes}
%==============================================================
This paper rewrites the comprehensive version (Parts I--V stitched together, 80 pages, 2026-11-22 v1) as one paper organized by theme. No claim has been upgraded. The changes are listed below.

\subsection{Updates of status wording}
\begin{enumerate}
\item \textbf{Local formula.} Parts II and III of the comprehensive version treated it as an ``observation/conjecture'' (\tP \tV \tK), but Part V proved it for quasi-homogeneous germs with $s\ne1$. Here it is collected in one place as Theorem \ref{thm:local}, with a shorter proof in the $(e,\chi_h)$ basis (the unchecked sign-convention and convergence issues of the old proof are removed; what remains \tQ is the identification of the local model with the stalk of a global Poisson module and the number of Saito's criterion). The rule $R(s)$ is the corrected version of Part II v2. We note that the threshold rule $[s<\mathrm{lct}]\cdot(\cdots)$ of the old version is wrong for $\mathrm{lct}<1/2$, and the global refined conjecture (Conjecture \ref{conj:refined}) is restricted to $\mathrm{lct}\ge1/2$.
\item \textbf{``Unresolved for $n>2$''.} The reduction theorem (Theorem \ref{thm:R}) is reproved without the filtration of $\mathrm{gr}\,Q$, in terms of the chain map $\varepsilon$ (condition (GEN$_Q$)). The step (GEN$'$) $\Rightarrow$ (GEN$_Q$) for general $n$ (Proposition \ref{prop:grQ}) remains \tQ. Theorem V (which was in the addendum) is now in the text of \S\ref{sec:reduction}, connected to the vanishing part of Theorem K (Theorem \ref{thm:K}) through Corollary \ref{cor:Vgen}. The old statement ``vanishing of $H^k(U;L_U)$ for $k<n$ is unproved'' is resolved when all $r_i\notin\Z$ (the mixed case is open).
\item \textbf{Triple-point anomalies.} The cancellation at $m=-2,-4$ is kept as an observation/conjecture (not explained by Theorem \ref{thm:local}).
\item \textbf{Name of (GEN).} To distinguish the hypothesis (GEN) of Theorem \ref{thm:T} (arithmetic genericity of the constant matrix $A$) from the hypothesis of the reduction theorem, the latter is renamed (GEN$_Q$).
\end{enumerate}

\subsection{Structural changes}
\begin{itemize}
\item The overview, symbol table, status table and ``scope of verification'' of each part are merged into one status table and one symbol table in \S\ref{sec:intro}. Numbering is a single global sequence common to definitions, theorems, propositions, lemmas, corollaries, examples, remarks, observations and conjectures.
\item The section on the ``classical three treasures'' of old Part I is compressed into the table of classical theorems and hypotheses (\S\ref{sec:prelim}). The remark on integrability in old Part I and some hypothesis tables for classical theorems not used here (Kodaira--Nakano etc.) are omitted.
\item The toric re-proof of old Part IV (via (GEN$'$)) is omitted since Theorem \ref{thm:T} has a direct proof; the relation is stated in \S\ref{sec:reduction}.
\item The example of $\mathbb{CP}^2$ in old Part I is rewritten as a derivation from Theorem \ref{thm:B} (Example \ref{ex:cp2}).
\item The K\"unneth proposition of old Part III is merged into Example \ref{ex:prod}, and the non-reduced boundary into Observation \ref{obs:nr} and Conjecture \ref{conj:NR}.
\item The section ``Checking the effect on Part I (v1, v2)'' of old Part II is absorbed into the hypothesis tables of the theorems and no longer stands as a section.
\end{itemize}

\subsection{Correspondence between old claims and this paper}
\begin{center}\small
\begin{tabular}{>{\raggedright\arraybackslash}p{5.4cm}>{\raggedright\arraybackslash}p{9.2cm}}\toprule
Old claim (comprehensive version) & This paper \\\midrule
Part I: GH bundle $=$ Poisson module, flat log connection & Prop.\ \ref{prop:ident}, Thm.\ \ref{thm:main1}\\
Part I: main theorem for surfaces, closed formula, periodic stabilization & Thms.\ \ref{thm:main2}, \ref{thm:B}, \ref{thm:A2}\\
Part I: Serre-type theorem, duality & Thms.\ \ref{thm:serre}, \ref{thm:dual}\\
Part II: Euler characteristic theorem & Thm.\ \ref{thm:E}\\
Part II: toric (surfaces, general dimension), resonance & Thms.\ \ref{thm:weight}, \ref{thm:C1}, \ref{thm:C2}, \ref{thm:T}, Obs.\ \ref{obs:res}\\
Parts II, III, V: local formula, triple points, non-reduced & Thm.\ \ref{thm:local}, Obs.\ \ref{obs:triple}, \ref{obs:nr}, Conj.\ \ref{conj:refined}, \ref{conj:NR}\\
Part III: Kodaira type & Props.\ \ref{prop:one}, \ref{prop:leaf}, Thms.\ \ref{thm:K1leaf}, \ref{thm:PKtoric}, Conj.\ \ref{conj:PK}\\
Part IV: regular type, reduction theorem & Thm.\ \ref{thm:L}, Cor.\ \ref{cor:L}, Thm.\ \ref{thm:R}\\
Part V: Theorem K1 (log Euler characteristic) & Thm.\ \ref{thm:K1}\\
Part V: Theorem K, index conjecture & Thm.\ \ref{thm:K}, Conj.\ \ref{conj:index}\\
Addendum: Theorem V, additional Poisson--Kodaira checks & Thm.\ \ref{thm:V}, Obs.\ \ref{obs:Segre}\\\bottomrule
\end{tabular}
\end{center}
``Theorem K1 (leaf-space Kodaira)'' of Part III and ``Theorem K1 (log Euler characteristic)'' of Part V in the comprehensive version were different theorems; here the former is Theorem \ref{thm:K1leaf} and the latter Theorem \ref{thm:K1}.

\subsection{Corrections and additions}
\begin{itemize}
\item The clash of the symbol $\rho$ (toric $\rho\in M_\C$ versus the dimension of the span of the $[D_i]$ for surfaces) is resolved by renaming the latter $\rho_D$.
\item For the quasi-homogeneous $D_5=z^2w+w^4$ the weights are $(3,2)$ and the degree is 8 (the table and text of the comprehensive version are correct; different weights in a research memo were wrong). The last column for $A_4$ in the local table is filled in.
\item The title of Hu's second paper follows the notation of Parts I, III of the comprehensive version (Lie groupoid Riemann--Roch--Hirzebruch theorem and applications, arXiv:2403.13424) (final check of the title \tQ).
\end{itemize}

\subsection{Scope of verification and scripts}
All computations are exact linear algebra over the rationals or finite fields ($32003$, $2^{61}-1$), or fits to known formulas; they are not proofs. Dependence on the choice of curves, and whether reduction modulo $p$ represents the characteristic 0 result, have not been verified.
\paragraph{Re-run while writing this paper.}
\begin{itemize}
\item \texttt{duality\_ndim.py}, \texttt{duality\_ndim\_exact.py}: $n\le4$, all agree.
\item \texttt{roundH\_proof\_check.py}: 14 quasi-homogeneous germs, 0 discrepancies.
\item \texttt{roundK\_check.py}: 175, 188 and 61 examples, 0 discrepancies.
\item \texttt{roundJ2\_vanish\_check.py}: 6 examples with $n=4$.
\item \texttt{roundD\_check.py}, \texttt{ncd\_local\_weights.py}, \texttt{ncd\_cochain\_check.py}, \texttt{serre\_rr\_check.py}.
\item Newly written, \texttt{local\_basis\_check.py}: 210 checks of formulas \eqref{eq:d0}, \eqref{eq:d1} in the $(e,\chi_h)$ basis, 0 discrepancies.
\end{itemize}
Outputs are saved in \texttt{logs/}.
\paragraph{Transcribed from the records of the comprehensive version, not re-run.} The large toric tables (\texttt{roundC\_*.py}, \texttt{toric\_hh.py}, \texttt{toric\_generic\_check2.py}), triple points (\texttt{local\_unequal*.py}), non-reduced (\texttt{roundF\_nonreduced.py}), numerics of Poisson--Kodaira (\texttt{roundE\_*.py}), the 108 examples of the resonance formula, and the regular-type examples (Kodaira--Thurston, Bieberbach).
\paragraph{Literature.} Many cited theorems remain whose range of applicability has not been checked in the original (the \tQ items in the text: re-derivation of LGSX Thm.~4.19 and Mokri's theorem, a close reading of Gualtieri 2011 \S4, the numbers of Saito's criterion and of Cox--Little--Schenck, the definition numbers of Pal--Poddar, Hu's bibliographic data, etc.).

\subsection{Low-confidence points (summary)}
(1) Proposition \ref{prop:ident} (depends on LGSX and Mokri, no re-derivation), (2) Proposition \ref{prop:grQ} (general-$n$ $\mathrm{gr}\,Q$ and chain-map property of $\varepsilon$), (3) the link between local model and global module in Theorem \ref{thm:local}, (4) consistency of Koszul signs in Theorem \ref{thm:dual} (hand computation only), (5) step (i) of Theorem \ref{thm:L} (identification with Gauss--Manin is a sketch), (6) the definition of $N_R$ in Observation \ref{obs:res}, (7) Lemma \ref{lem:lift} (sketch).

%==============================================================
\begin{thebibliography}{99}\small
\bibitem{AB} M.\ Abouzaid, M.\ Boyarchenko, Local structure of generalized complex manifolds, J.\ Symplectic Geom.\ 4 (2006) 43--62, arXiv:math/0511049.
\bibitem{ALN} B.\ Ammann, R.\ Lauter, V.\ Nistor, Pseudodifferential operators on manifolds with a Lie structure at infinity, arXiv:math/0304044.
\bibitem{AN} Y.\ Akizuki, S.\ Nakano, Note on Kodaira--Spencer's proof of Lefschetz's theorems, Proc.\ Japan Acad.\ 30 (1954) 266--272.
\bibitem{AS} M.\ Atiyah, I.\ Singer, The index of elliptic operators III, Ann.\ Math.\ 87 (1968) 546--604.
\bibitem{Ara} D.\ Arapura, Geometry of cohomology support loci for local systems I, arXiv:alg-geom/9612006; J.\ Algebraic Geom.\ 6 (1997) 563--597 [volume/pages from memory \tQ].
\bibitem{Atiyah} M.\ Atiyah, Vector bundles over an elliptic curve, Proc.\ London Math.\ Soc.\ (3) 7 (1957) 414--452.
\bibitem{BG} A.\ Bailey, M.\ Gualtieri, Integration of generalized complex structures, arXiv:1611.03850.
\bibitem{BLS} M.\ Braverman, Y.\ Loizides, Y.\ Song, Geometric quantization of $b$-symplectic manifolds, arXiv:1910.10016.
\bibitem{Bai} W.\ L.\ Baily, On the imbedding of V-manifolds in projective space, Amer.\ J.\ Math.\ 79 (1957) 403--430 [details \tQ].
\bibitem{Bre} G.\ Bredon, \emph{Sheaf Theory}, 2nd ed., Springer GTM 170, 1997 [details \tQ].
\bibitem{CF} M.\ Crainic, R.\ L.\ Fernandes, Integrability of Lie brackets, Ann.\ Math.\ 157 (2003) 575--620, arXiv:math/0105033.
\bibitem{CLS} D.\ Cox, J.\ Little, H.\ Schenck, \emph{Toric Varieties}, Graduate Studies in Math.\ 124, AMS 2011 [details \tQ].
\bibitem{CMNM} F.\ J.\ Calder\'on-Moreno, L.\ Narv\'aez-Macarro, Logarithmic comparison theorems for free divisors (Dualit\'e et comparaison sur les complexes de de Rham logarithmiques par rapport aux diviseurs libres), Proc.\ London Math.\ Soc.\ 98 (2009), doi:10.1112/plms/pdn043, arXiv:math/0603003 [details \tQ].
\bibitem{Cav} G.\ Cavalcanti, New aspects of the $dd^c$-lemma, arXiv:math/0501406; The decomposition of forms and cohomology of generalized complex manifolds, J.\ Geom.\ Phys.\ 57 (2006) 121--132.
\bibitem{Del} P.\ Deligne, \'Equations diff\'erentielles \`a points singuliers r\'eguliers, Lecture Notes in Math.\ 163, Springer 1970.
\bibitem{Del68} P.\ Deligne, Th\'eor\`eme de Lefschetz et crit\`eres de d\'eg\'en\'erescence de suites spectrales, Publ.\ Math.\ IH\'ES 35 (1968) 259--278.
\bibitem{DelFG} P.\ Deligne, Le groupe fondamental du compl\'ement d'une courbe plane n'ayant que des points doubles ordinaires est ab\'elien, S\'em.\ Bourbaki 543 (1979/80).
\bibitem{EV} H.\ Esnault, E.\ Viehweg, Logarithmic de Rham complexes and vanishing theorems, Invent.\ Math.\ 86 (1986) 161--194.
\bibitem{Ful} W.\ Fulton, On the fundamental group of the complement of a node curve, Ann.\ Math.\ 111 (1980) 407--409.
\bibitem{GH} P.\ Griffiths, J.\ Harris, Principles of Algebraic Geometry, Wiley 1978.
\bibitem{GL} M.\ Gualtieri, S.\ Li, Symplectic groupoids of log symplectic manifolds, IMRN 2014, arXiv:1206.3674.
\bibitem{Goto} R.\ Goto, Unobstructed deformations of generalized complex structures induced by $C^\infty$ logarithmic symplectic structures, arXiv:1501.03398 [journal and pages unchecked \tQ].
\bibitem{Gua11} M.\ Gualtieri, Generalized complex geometry, Ann.\ Math.\ 174 (2011) 75--123, arXiv:math/0703298.
\bibitem{GuaBr} M.\ Gualtieri, Branes on Poisson varieties, arXiv:0710.2719 (v2, 2010-07; no publication data on the arXiv page; The Many Facets of Geometry, Oxford Univ.\ Press 2010, from memory \tQ).
\bibitem{GuaHodge} M.\ Gualtieri, Generalized geometry and the Hodge decomposition, arXiv:math/0409093.
\bibitem{GuaK} M.\ Gualtieri, Generalized K\"ahler geometry, Comm.\ Math.\ Phys.\ 331 (2014) 297--331, arXiv:1007.3485.
\bibitem{HHM} T.\ Hausel, E.\ Hunsicker, R.\ Mazzeo, Hodge cohomology of gravitational instantons, Duke Math.\ J.\ 122 (2004) 485--548, arXiv:math/0207169.
\bibitem{Hart} R.\ Hartshorne, Algebraic Geometry, GTM 52, Springer 1977.
\bibitem{Hit03} N.\ Hitchin, Generalized Calabi--Yau manifolds, Q.\ J.\ Math.\ 54 (2003) 281--308, arXiv:math/0209099.
\bibitem{Hu1} T.\ Hu, Kodaira vanishing theorems for K\"ahler Lie algebroids, arXiv:2403.10876.
\bibitem{Hu2} T.\ Hu, Lie groupoid Riemann--Roch--Hirzebruch theorem and applications, arXiv:2403.13424.
\bibitem{Huy} D.\ Huybrechts, Complex Geometry: An Introduction, Springer 2005.
\bibitem{Kod53} K.\ Kodaira, On a differential-geometric method in the theory of analytic stacks, Proc.\ Nat.\ Acad.\ Sci.\ USA 39 (1953) 1268--1273.
\bibitem{Kod54} K.\ Kodaira, On K\"ahler varieties of restricted type, Ann.\ Math.\ 60 (1954) 28--48.
\bibitem{LGSX1} C.\ Laurent-Gengoux, M.\ Sti\'enon, P.\ Xu, Holomorphic Poisson manifolds and holomorphic Lie algebroids, IMRN 2008, arXiv:0707.4253.
\bibitem{LGSX2} C.\ Laurent-Gengoux, M.\ Sti\'enon, P.\ Xu, Integration of holomorphic Lie algebroids, arXiv:0803.2031.
\bibitem{MM} R.\ Mazzeo, R.\ Melrose, Pseudodifferential operators on manifolds with fibred boundaries, arXiv:math/9812120.
\bibitem{Milnor} J.\ Milnor, \emph{Singular Points of Complex Hypersurfaces}, Annals of Math.\ Studies 61, Princeton Univ.\ Press 1968 [details \tQ].
\bibitem{Mok} T.\ Mokri, Matched pairs of Lie algebroids, Glasgow Math.\ J.\ 39 (1997) 167--181, doi:10.1017/S0017089500032055 [year and DOI checked on the publisher page, volume/pages from memory \tQ].
\bibitem{Mum} D.\ Mumford, Varieties defined by quadratic equations, in: Questions on Algebraic Varieties (C.I.M.E., Varenna 1969), Cremonese 1970 [details \tQ].
\bibitem{PP} D.\ Pal, M.\ Poddar, Strong generalized holomorphic principal bundles, arXiv:2404.18113; Ann.\ Sc.\ Norm.\ Super.\ Pisa Cl.\ Sci., doi:10.2422/2036-2145.202501\_026 (volume/issue \tQ).
\bibitem{PPT} M.\ Pflaum, H.\ Posthuma, X.\ Tang, The index of geometric operators on Lie groupoids, Indag.\ Math.\ 25 (2014) 1--18, arXiv:1308.0236.
\bibitem{Pym} B.\ Pym, Constructions and classifications of projective Poisson varieties, Lett.\ Math.\ Phys.\ 108 (2018), no.\ 3, 573--632, doi:10.1007/s11005-017-0984-5.
\bibitem{QS} Y.\ Qiao, Y.\ So, Deformation and K-theoretic index formulae on boundary groupoids, arXiv:2304.13345.
\bibitem{SGA4} M.\ Artin, Th\'eor\`eme de finitude pour un morphisme propre; dimension cohomologique des sch\'emas alg\'ebriques affines, in SGA 4, Exp.\ XIV, Lecture Notes in Math.\ 305 [details \tQ].
\bibitem{Saito} K.\ Saito, Theory of logarithmic differential forms and logarithmic vector fields, J.\ Fac.\ Sci.\ Univ.\ Tokyo Sect.\ IA Math.\ 27 (1980) 265--291 [theorem number of the criterion from memory \tQ].
\bibitem{Serre} J.-P.\ Serre, Faisceaux alg\'ebriques coh\'erents, Ann.\ Math.\ 61 (1955) 197--278.
\bibitem{Tim} K.\ Timmerscheidt, Mixed Hodge theory for unitary local systems, J.\ reine angew.\ Math.\ 379 (1987) 152--171.
\bibitem{Hirz} F.\ Hirzebruch, Neue topologische Methoden in der algebraischen Geometrie, Springer 1956.
\end{thebibliography}

\end{document}
